%% Generated by Sphinx.
\def\sphinxdocclass{report}
\documentclass[letterpaper,10pt,english]{sphinxmanual}
\ifdefined\pdfpxdimen
   \let\sphinxpxdimen\pdfpxdimen\else\newdimen\sphinxpxdimen
\fi \sphinxpxdimen=.75bp\relax
\ifdefined\pdfimageresolution
    \pdfimageresolution= \numexpr \dimexpr1in\relax/\sphinxpxdimen\relax
\fi
%% let collapsible pdf bookmarks panel have high depth per default
\PassOptionsToPackage{bookmarksdepth=5}{hyperref}

\PassOptionsToPackage{booktabs}{sphinx}
\PassOptionsToPackage{colorrows}{sphinx}

\PassOptionsToPackage{warn}{textcomp}
\usepackage[utf8]{inputenc}
\ifdefined\DeclareUnicodeCharacter
% support both utf8 and utf8x syntaxes
  \ifdefined\DeclareUnicodeCharacterAsOptional
    \def\sphinxDUC#1{\DeclareUnicodeCharacter{"#1}}
  \else
    \let\sphinxDUC\DeclareUnicodeCharacter
  \fi
  \sphinxDUC{00A0}{\nobreakspace}
  \sphinxDUC{2500}{\sphinxunichar{2500}}
  \sphinxDUC{2502}{\sphinxunichar{2502}}
  \sphinxDUC{2514}{\sphinxunichar{2514}}
  \sphinxDUC{251C}{\sphinxunichar{251C}}
  \sphinxDUC{2572}{\textbackslash}
\fi
\usepackage{cmap}
\usepackage[T1]{fontenc}
\usepackage{amsmath,amssymb,amstext}
\usepackage{babel}



\usepackage{tgtermes}
\usepackage{tgheros}
\renewcommand{\ttdefault}{txtt}



\usepackage[Bjarne]{fncychap}
\usepackage{sphinx}

\fvset{fontsize=auto}
\usepackage{geometry}


% Include hyperref last.
\usepackage{hyperref}
% Fix anchor placement for figures with captions.
\usepackage{hypcap}% it must be loaded after hyperref.
% Set up styles of URL: it should be placed after hyperref.
\urlstyle{same}

\addto\captionsenglish{\renewcommand{\contentsname}{Contents:}}

\usepackage{sphinxmessages}
\setcounter{tocdepth}{1}



\title{BacLipidAPP}
\date{Jul 22, 2026}
\release{1.0.0}
\author{Lilia CHALES}
\newcommand{\sphinxlogo}{\vbox{}}
\renewcommand{\releasename}{Release}
\makeindex
\begin{document}

\ifdefined\shorthandoff
  \ifnum\catcode`\=\string=\active\shorthandoff{=}\fi
  \ifnum\catcode`\"=\active\shorthandoff{"}\fi
\fi

\pagestyle{empty}
\sphinxmaketitle
\pagestyle{plain}
\sphinxtableofcontents
\pagestyle{normal}
\phantomsection\label{\detokenize{index::doc}}


\sphinxAtStartPar
Add your content using \sphinxcode{\sphinxupquote{reStructuredText}} syntax. See the
\sphinxhref{https://www.sphinx-doc.org/en/master/usage/restructuredtext/index.html}{reStructuredText}
documentation for details.

\sphinxstepscope


\chapter{Home}
\label{\detokenize{Home:home}}\label{\detokenize{Home::doc}}
\sphinxAtStartPar
Fichier source : \sphinxcode{\sphinxupquote{app/Home.py}}

\begin{sphinxadmonition}{note}{Note:}
\sphinxAtStartPar
Cette page contient la documentation du module Streamlit (\sphinxcode{\sphinxupquote{Home.py}}).
La majorité du code n’est pas écrit sous forme de fonction, il n’est donc pas possible d’utiliser \sphinxcode{\sphinxupquote{automodule}}
pour générer cette page automatiquement. La documentation ci\sphinxhyphen{}dessous est rédigée manuellement
à partir du code source.
\end{sphinxadmonition}


\section{Vue d’ensemble}
\label{\detokenize{Home:vue-d-ensemble}}
\sphinxAtStartPar
Cette page est la page d’accueil de \sphinxstylestrong{BacLipidAPP}. Elle affiche le logo de
l’application, une description générale, des liens vers les cinq modules
disponibles ({\hyperref[\detokenize{1_Integration::doc}]{\sphinxcrossref{\DUrole{doc}{Data\_Integration}}}}, {\hyperref[\detokenize{2_BacLipidDB::doc}]{\sphinxcrossref{\DUrole{doc}{BacLipidDB}}}}, {\hyperref[\detokenize{3_Export::doc}]{\sphinxcrossref{\DUrole{doc}{Export}}}},
{\hyperref[\detokenize{4_Resources::doc}]{\sphinxcrossref{\DUrole{doc}{Resources}}}}, {\hyperref[\detokenize{5_Management::doc}]{\sphinxcrossref{\DUrole{doc}{Management}}}}) ainsi qu’un aperçu synthétique du
contenu de \sphinxstylestrong{BacLipidDB}.


\section{Logo et description}
\label{\detokenize{Home:logo-et-description}}
\sphinxAtStartPar
Le logo (\sphinxcode{\sphinxupquote{assets/APP.svg}}) est affiché centré (\sphinxcode{\sphinxupquote{st.columns}}).
Une description de l’application est ensuite affichée via \sphinxcode{\sphinxupquote{st.markdown}}
(HTML autorisé via \sphinxcode{\sphinxupquote{unsafe\_allow\_html=True}}).


\section{Modules disponibles}
\label{\detokenize{Home:modules-disponibles}}
\sphinxAtStartPar
Cinq encadrés (\sphinxcode{\sphinxupquote{st.container(border=True)}}), répartis sur \sphinxcode{\sphinxupquote{st.columns(5)}},
présentent chacun un module et pointent vers sa page (\sphinxcode{\sphinxupquote{st.page\_link}}) :
\begin{itemize}
\item {} 
\sphinxAtStartPar
\sphinxstylestrong{MODULE 1} : lien vers \sphinxcode{\sphinxupquote{pages/1\_Integration.py}}, intégration des données.

\item {} 
\sphinxAtStartPar
\sphinxstylestrong{MODULE 2} : lien vers \sphinxcode{\sphinxupquote{pages/2\_BacLipidDB.py}}, exploration de la base de données \sphinxstylestrong{BacLipidDB}

\item {} 
\sphinxAtStartPar
\sphinxstylestrong{MODULE 3} : lien vers \sphinxcode{\sphinxupquote{pages/3\_Export.py}}, export des données dans des formats compatibles avec MZmine.

\item {} 
\sphinxAtStartPar
\sphinxstylestrong{MODULE 4} : lien vers \sphinxcode{\sphinxupquote{pages/4\_Resources.py}}, guide des colonnes pour les fichiers d’intégration.

\item {} 
\sphinxAtStartPar
\sphinxstylestrong{MODULE 5} : lien vers \sphinxcode{\sphinxupquote{pages/5\_Management.py}}, correction/suppression des enregistrements déjà intégrés.

\end{itemize}


\section{Aperçu de la base de données}
\label{\detokenize{Home:apercu-de-la-base-de-donnees}}
\sphinxAtStartPar
La fonction \sphinxcode{\sphinxupquote{statistics()}} (mise en cache via \sphinxcode{\sphinxupquote{st.cache\_data(ttl=60)}})
interroge la base de données via une session SQLAlchemy et retourne :
\begin{itemize}
\item {} 
\sphinxAtStartPar
le nombre de détections \sphinxcode{\sphinxupquote{MS1}} (\sphinxcode{\sphinxupquote{Detection.MS\_level == "MS1"}}) ;

\item {} 
\sphinxAtStartPar
le nombre de détections \sphinxcode{\sphinxupquote{MS2}} (\sphinxcode{\sphinxupquote{Detection.MS\_level == "MS2"}}) ;

\item {} 
\sphinxAtStartPar
la taille du fichier \sphinxcode{\sphinxupquote{BacLipidDB.db}} (formatée en Ko ou Mo selon sa taille).

\end{itemize}

\sphinxAtStartPar
En cas d’erreur de connexion, la fonction retourne \sphinxcode{\sphinxupquote{None}} et le message
\sphinxcode{\sphinxupquote{st.warning("Unable to connect to the database.")}} est affiché à la place
des métriques. Sinon, les trois valeurs sont affichées via \sphinxcode{\sphinxupquote{st.metric}} dans
des encadrés (\sphinxcode{\sphinxupquote{st.container(border=True)}}), avec pour libellés respectifs
\sphinxcode{\sphinxupquote{"MS1 detections"}}, \sphinxcode{\sphinxupquote{"MS2 detections"}} et \sphinxcode{\sphinxupquote{"Database size"}}.


\section{Fonctions internes}
\label{\detokenize{Home:fonctions-internes}}

\begin{savenotes}\sphinxattablestart
\sphinxthistablewithglobalstyle
\centering
\begin{tabular}[t]{\X{30}{100}\X{70}{100}}
\sphinxtoprule
\sphinxstyletheadfamily 
\sphinxAtStartPar
Fonction
&\sphinxstyletheadfamily 
\sphinxAtStartPar
Description
\\
\sphinxmidrule
\sphinxtableatstartofbodyhook
\sphinxAtStartPar
\sphinxcode{\sphinxupquote{statistics() \sphinxhyphen{}\textgreater{} dict | None}}
&
\sphinxAtStartPar
Calcule et met en cache le nombre de détections \sphinxcode{\sphinxupquote{MS1}}/\sphinxcode{\sphinxupquote{MS2}} et la
taille du fichier de base de données. Retourne \sphinxcode{\sphinxupquote{None}} en cas d’erreur.
\\
\sphinxbottomrule
\end{tabular}
\sphinxtableafterendhook\par
\sphinxattableend\end{savenotes}


\section{Dépendances internes}
\label{\detokenize{Home:dependances-internes}}\begin{itemize}
\item {} 
\sphinxAtStartPar
{\hyperref[\detokenize{api:module-models.model}]{\sphinxcrossref{\sphinxcode{\sphinxupquote{models.model}}}}} (\sphinxcode{\sphinxupquote{Detection}})

\item {} 
\sphinxAtStartPar
{\hyperref[\detokenize{api:module-config}]{\sphinxcrossref{\sphinxcode{\sphinxupquote{config}}}}} (\sphinxcode{\sphinxupquote{get\_engine}})

\end{itemize}


\section{Fonctionnement général de Streamlit}
\label{\detokenize{Home:fonctionnement-general-de-streamlit}}
\sphinxAtStartPar
Cette section décrit deux mécanismes de Streamlit qui reviennent dans
plusieurs modules de l’application.

\sphinxAtStartPar
\sphinxstylestrong{Le rerun}

\sphinxAtStartPar
À chaque interaction de l’utilisateur (clic sur un bouton, changement de
valeur d’un widget, etc.), Streamlit \sphinxstylestrong{réexécute l’intégralité du script
Python de haut en bas}. Il n’y a pas de gestion d’événements comme dans une
application web classique, ce qui signifie que le code d’une page est rejoué à chaque
interaction, y compris les widgets déjà affichés.

\sphinxAtStartPar
\sphinxstylestrong{\textasciigrave{}\textasciigrave{}st.session\_state\textasciigrave{}\textasciigrave{} et la \textasciigrave{}\textasciigrave{}key\textasciigrave{}\textasciigrave{} des widgets}

\sphinxAtStartPar
Puisque le script est rejoué en entier à chaque interaction, Streamlit a
besoin d’un moyen d’associer un widget (ex. \sphinxcode{\sphinxupquote{st.file\_uploader(...)}})
rappelé lors d’un rerun à la valeur qu’il avait lors du rerun précédent.
Cette association se fait via la \sphinxcode{\sphinxupquote{key}} du widget : c’est elle qui permet
de retrouver la valeur précédente (fichier déposé, texte saisi, tableur…) et de la
lui réattribuer, plutôt que de le réinitialiser à chaque fois.

\sphinxAtStartPar
C’est ce que permet \sphinxcode{\sphinxupquote{st.session\_state}} : c’est un dictionnaire persistant,
propre à chaque session utilisateur qui survit aux reruns (contrairement
aux variables Python classiques). On peut se le représenter comme une armoire à casiers, où
chaque \sphinxcode{\sphinxupquote{key}} de widget désigne un casier précis. Par exemple, quand le script appelle
\sphinxcode{\sphinxupquote{st.file\_uploader(key="mon\_widget")}} Streamlit vérifie si le casier
\sphinxcode{\sphinxupquote{"mon\_widget"}} contient déjà une valeur (issue d’un rerun précédent) :
si oui, il l’utilise pour réafficher le widget dans son état existant
sinon, il crée le casier et l’initialise.

\sphinxAtStartPar
Pour la plupart des widgets (\sphinxcode{\sphinxupquote{st.selectbox}}, \sphinxcode{\sphinxupquote{st.text\_input}}…), on
peut aussi écrire directement dans \sphinxcode{\sphinxupquote{st.session\_state}} pour changer leur
valeur par programmation (ex. les remettre à \sphinxcode{\sphinxupquote{None}}/\sphinxcode{\sphinxupquote{""}} depuis un
bouton Reset). Mais certains widgets (\sphinxcode{\sphinxupquote{st.file\_uploader}}, \sphinxcode{\sphinxupquote{st.button}}) \sphinxstylestrong{interdisent} cette écriture directe (Streamlit
lève une exception). Pour les vider malgré tout depuis le code, la seule
solution est de changer leur \sphinxcode{\sphinxupquote{key}} : Streamlit regarde alors dans un
casier différent vide. Par exemple l’usage d’un compteur
de version (ex. \sphinxcode{\sphinxupquote{UPLOADER\_VERSION}}) inclus dans la \sphinxcode{\sphinxupquote{key}} pour forcer la
recréation d’un widget vide.

\sphinxAtStartPar
Exemple concret tiré de la page \sphinxstylestrong{Integration} (\sphinxcode{\sphinxupquote{pages/1\_Integration.py}}) :
\begin{itemize}
\item {} 
\sphinxAtStartPar
Au premier passage, \sphinxcode{\sphinxupquote{st.session\_state}} ne contient pas encore la clé
\sphinxcode{\sphinxupquote{"uploader\_version"}}, donc \sphinxcode{\sphinxupquote{st.session\_state.get(UPLOADER\_VERSION, 0)}}
vaut \sphinxcode{\sphinxupquote{0}}. Le widget est créé avec \sphinxcode{\sphinxupquote{key="file\_uploader\_0"}}. Si
l’utilisateur dépose un fichier, celui\sphinxhyphen{}ci est rangé dans le casier
\sphinxcode{\sphinxupquote{"file\_uploader\_0"}}.

\item {} 
\sphinxAtStartPar
Quand l’utilisateur clique sur \sphinxstylestrong{Reset 🔄}, le code exécute
\sphinxcode{\sphinxupquote{st.session\_state{[}UPLOADER\_VERSION{]} = st.session\_state.get(UPLOADER\_VERSION, 0) + 1}},
ce qui fait passer le compteur de \sphinxcode{\sphinxupquote{0}} à \sphinxcode{\sphinxupquote{1}}, puis appelle \sphinxcode{\sphinxupquote{st.rerun()}}.

\item {} 
\sphinxAtStartPar
Au rerun suivant, la \sphinxcode{\sphinxupquote{key}} est réévaluée et vaut désormais
\sphinxcode{\sphinxupquote{"file\_uploader\_1"}}. Ce casier n’a jamais été utilisé : il est vide, et le
widget s’affiche donc vide. Le casier \sphinxcode{\sphinxupquote{"file\_uploader\_0"}} existe toujours
quelque part (avec l’ancien fichier dedans), mais n’est plus consulté
puisque la clé du widget a changé.

\end{itemize}

\sphinxstepscope


\chapter{Data\_Integration}
\label{\detokenize{1_Integration:data-integration}}\label{\detokenize{1_Integration::doc}}
\sphinxAtStartPar
Fichier source : \sphinxcode{\sphinxupquote{app/pages/1\_Integration.py}}

\begin{sphinxadmonition}{note}{Note:}
\sphinxAtStartPar
Cette page contient la documentation du module Streamlit (\sphinxcode{\sphinxupquote{1\_Integration.py}}).
La majorité du code n’est pas écrit sous forme de fonction, il n’est donc pas possible d’utiliser \sphinxcode{\sphinxupquote{automodule}}
pour générer cette page automatiquement. La documentation ci\sphinxhyphen{}dessous est rédigée manuellement
à partir du code source.
\end{sphinxadmonition}


\section{Vue d’ensemble}
\label{\detokenize{1_Integration:vue-d-ensemble}}
\sphinxAtStartPar
Ce module Streamlit (\sphinxstylestrong{MODULE 1}) permet d’intégrer un fichier d’annotations
lipidiques, \sphinxstylestrong{MS1} ou \sphinxstylestrong{MS2}, dans la base de données \sphinxstylestrong{BacLipidDB}, via un
workflow guidé en 5 étapes affiché à la fois dans le corps de la page et dans
la barre latérale (icônes ✅ / ⬜). Le niveau MS choisi à l’étape 1 détermine
le format de fichier attendu et le chemin de traitement des étapes suivantes.


\section{Colonnes du fichier d’entrée}
\label{\detokenize{1_Integration:colonnes-du-fichier-d-entree}}
\sphinxAtStartPar
\sphinxstylestrong{Fichier MS1} (\sphinxcode{\sphinxupquote{.csv}}/\sphinxcode{\sphinxupquote{.tsv}}/\sphinxcode{\sphinxupquote{.xlsx}}, une ligne = une annotation)

\sphinxAtStartPar
Colonnes obligatoires (\sphinxcode{\sphinxupquote{REQUIRED\_COLUMNS}}), doivent être présentes dans le fichier importé
(voir \sphinxcode{\sphinxupquote{4\_Resources.py}}) :
\begin{itemize}
\item {} 
\sphinxAtStartPar
\sphinxcode{\sphinxupquote{Lipid\_Name}}

\item {} 
\sphinxAtStartPar
\sphinxcode{\sphinxupquote{Formula}}

\item {} 
\sphinxAtStartPar
\sphinxcode{\sphinxupquote{Precursor\_MZ}}

\item {} 
\sphinxAtStartPar
\sphinxcode{\sphinxupquote{Lipid\_category}}

\item {} 
\sphinxAtStartPar
\sphinxcode{\sphinxupquote{Lipid\_class}}

\item {} 
\sphinxAtStartPar
\sphinxcode{\sphinxupquote{Lipid\_subclass}}

\item {} 
\sphinxAtStartPar
\sphinxcode{\sphinxupquote{Adduct}}

\end{itemize}

\sphinxAtStartPar
Parmi celles\sphinxhyphen{}ci, \sphinxcode{\sphinxupquote{Lipid\_Name}}, \sphinxcode{\sphinxupquote{Formula}}, \sphinxcode{\sphinxupquote{Precursor\_MZ}} et \sphinxcode{\sphinxupquote{Adduct}}
(\sphinxcode{\sphinxupquote{NON\_EMPTY\_COLUMNS}}) ne doivent contenir aucune cellule vide.

\sphinxAtStartPar
Colonnes optionnelles : \sphinxcode{\sphinxupquote{RT}}, \sphinxcode{\sphinxupquote{CCS}}. Si elles sont présentes dans le
fichier importé, elles sont conservées telles quelles dans le tableau et
transmises à {\hyperref[\detokenize{api:utils.loading_MS1.DB_MS1}]{\sphinxcrossref{\sphinxcode{\sphinxupquote{utils.loading\_MS1.DB\_MS1()}}}}}, qui les insère dans la table
\sphinxcode{\sphinxupquote{Detection}} (\sphinxcode{\sphinxupquote{None}} si absentes ou vides).

\sphinxAtStartPar
\sphinxstylestrong{Fichier MS2} (\sphinxcode{\sphinxupquote{.csv}}/\sphinxcode{\sphinxupquote{.tsv}} organisé en blocs empilés, un bloc = un scan)

\sphinxAtStartPar
Le fichier est parsé par \sphinxcode{\sphinxupquote{utils.loading\_MS2.ms2\_parsing()}}, qui valide
elle\sphinxhyphen{}même chaque bloc. Chaque bloc commence par une ligne d’en\sphinxhyphen{}tête
\sphinxcode{\sphinxupquote{Scan \#\textless{}identifiant\textgreater{}}} dont l’identifiant est extrait sous la clé
\sphinxcode{\sphinxupquote{scan\_id}}. Ce \sphinxcode{\sphinxupquote{scan\_id}} n’est utilisé qu’en interne (voir
{\hyperref[\detokenize{1_Integration:workflow}]{\sphinxcrossref{Workflow}}}, étape 4) : il n’est ni affiché ni modifiable dans l’application.
Champs obligatoires par scan : \sphinxcode{\sphinxupquote{Lipid\_Name}},
\sphinxcode{\sphinxupquote{Precursor\_MZ}}, \sphinxcode{\sphinxupquote{Formula}}, \sphinxcode{\sphinxupquote{Adduct}}, \sphinxcode{\sphinxupquote{Lipid\_category}},
\sphinxcode{\sphinxupquote{Lipid\_class}}, \sphinxcode{\sphinxupquote{Lipid\_subclass}}, \sphinxcode{\sphinxupquote{Num\_peaks}}, suivis d’un en\sphinxhyphen{}tête
\sphinxcode{\sphinxupquote{m/z}} / \sphinxcode{\sphinxupquote{Intensity}} et d’au moins une ligne de fragment. Champs optionnels :
\sphinxcode{\sphinxupquote{FA\_composition}}, \sphinxcode{\sphinxupquote{RT}}, \sphinxcode{\sphinxupquote{CCS}}.


\section{Éléments d’interface communs}
\label{\detokenize{1_Integration:elements-d-interface-communs}}\begin{itemize}
\item {} 
\sphinxAtStartPar
\sphinxstylestrong{Barre latérale} : récapitulatif du workflow (icône par étape) et bouton
\sphinxstylestrong{Reset 🔄}, qui effectue trois actions :
\begin{itemize}
\item {} 
\sphinxAtStartPar
remet les widgets de l’étape 1 à \sphinxcode{\sphinxupquote{None}} (\sphinxcode{\sphinxupquote{ms\_level}}, \sphinxcode{\sphinxupquote{ion\_mode}}) et à \sphinxcode{\sphinxupquote{""}}
(\sphinxcode{\sphinxupquote{integrator\_name}}, \sphinxcode{\sphinxupquote{file\_row\_name}}) ;

\item {} 
\sphinxAtStartPar
supprime de \sphinxcode{\sphinxupquote{st.session\_state}} toutes les clés calculées listées dans
\sphinxcode{\sphinxupquote{RESET\_KEYS}}, ce qui efface l’état des étapes 2 à 5 ;

\item {} 
\sphinxAtStartPar
incrémente \sphinxcode{\sphinxupquote{UPLOADER\_VERSION}}, ce qui change la \sphinxcode{\sphinxupquote{key}} du widget
\sphinxcode{\sphinxupquote{st.file\_uploader}} de l’étape 2
(\sphinxcode{\sphinxupquote{key=f"file\_uploader\_\{st.session\_state.get(UPLOADER\_VERSION, 0)\}"}}) et
le vide ainsi indirectement \sphinxhyphen{} voir {\hyperref[\detokenize{Home::doc}]{\sphinxcrossref{\DUrole{doc}{Home}}}}, section
\sphinxstyleemphasis{Fonctionnement général de Streamlit}, pour le détail de ce mécanisme
(\sphinxcode{\sphinxupquote{st.file\_uploader}} n’autorisant pas d’écriture directe dans
\sphinxcode{\sphinxupquote{st.session\_state}}).

\end{itemize}

\item {} 
\sphinxAtStartPar
\sphinxstylestrong{En\sphinxhyphen{}tête} : bandeau \sphinxcode{\sphinxupquote{MODULE 1 : Data integration page}}.

\item {} 
\sphinxAtStartPar
\sphinxstylestrong{Barre de progression} : nombre d’étapes complétées sur 5.

\item {} 
\sphinxAtStartPar
\sphinxstylestrong{Encadré d’aide} (\sphinxcode{\sphinxupquote{st.expander}}) : résumé du workflow en 5 points.

\end{itemize}


\section{Workflow}
\label{\detokenize{1_Integration:workflow}}\begin{enumerate}
\sphinxsetlistlabels{\arabic}{enumi}{enumii}{}{.}%
\item {} 
\sphinxAtStartPar
\sphinxstylestrong{Settings}
Sélection du niveau MS (\sphinxcode{\sphinxupquote{ms\_level}} : \sphinxcode{\sphinxupquote{MS1}} ou \sphinxcode{\sphinxupquote{MS2}}), du mode
d’ionisation (\sphinxcode{\sphinxupquote{ion\_mode}} : \sphinxcode{\sphinxupquote{Positive}} / \sphinxcode{\sphinxupquote{Negative}}), du nom de
l’intégrateur (\sphinxcode{\sphinxupquote{integrator\_name}}) et du nom du fichier expérimental
d’origine (\sphinxcode{\sphinxupquote{file\_row\_name}}). Tant que ces quatre champs ne sont pas
renseignés, la page s’arrête (\sphinxcode{\sphinxupquote{st.stop()}}).

\sphinxAtStartPar
Une fois la complétion automatique de l’étape 4 effectuée (\sphinxcode{\sphinxupquote{df\_complete}}
présent en session), ces quatre champs sont verrouillés (\sphinxcode{\sphinxupquote{disabled}})
afin de garantir la cohérence entre les colonnes calculées et les réglages
utilisés pour les calculer.

\item {} 
\sphinxAtStartPar
\sphinxstylestrong{File upload}
Chargement d’un fichier \sphinxcode{\sphinxupquote{.xlsx}}, \sphinxcode{\sphinxupquote{.csv}} ou \sphinxcode{\sphinxupquote{.tsv}} via
\sphinxcode{\sphinxupquote{st.file\_uploader}} (clé versionnée par \sphinxcode{\sphinxupquote{UPLOADER\_VERSION}} pour pouvoir
être réinitialisé par le bouton \sphinxstylestrong{Reset}). Le rechargement du fichier (parsing) n’est
déclenché que si son identifiant  (\sphinxcode{\sphinxupquote{uploaded\_file.file\_id}}) ou le
niveau MS (\sphinxcode{\sphinxupquote{df\_ms\_level}}) diffère de la dernière valeur connue.
\begin{itemize}
\item {} 
\sphinxAtStartPar
\sphinxcode{\sphinxupquote{MS2}} \(\rightarrow\) \sphinxcode{\sphinxupquote{utils.loading\_MS2.ms2\_parsing()}} (liste de scans stockée
sous la clé \sphinxcode{\sphinxupquote{ms2}}) ;

\item {} 
\sphinxAtStartPar
\sphinxcode{\sphinxupquote{MS1}} \(\rightarrow\) \sphinxcode{\sphinxupquote{pandas.read\_csv()}} (\sphinxcode{\sphinxupquote{.csv}}/\sphinxcode{\sphinxupquote{.tsv}}) ou
\sphinxcode{\sphinxupquote{pandas.read\_excel()}} (\sphinxcode{\sphinxupquote{.xlsx}}), stocké sous la clé \sphinxcode{\sphinxupquote{ms1}}.

\end{itemize}

\sphinxAtStartPar
Le chargement d’un nouveau fichier efface les clés \sphinxcode{\sphinxupquote{RELOAD\_RESET\_KEYS}} (résultats
calculés en aval) ainsi que la donnée de l’autre niveau MS (\sphinxcode{\sphinxupquote{ms1}} si l’on
bascule vers \sphinxcode{\sphinxupquote{MS2}}, \sphinxcode{\sphinxupquote{ms2}} si l’on bascule vers \sphinxcode{\sphinxupquote{MS1}}), tout en
conservant les réglages de l’étape 1.

\item {} 
\sphinxAtStartPar
\sphinxstylestrong{Verification}
\begin{itemize}
\item {} 
\sphinxAtStartPar
\sphinxstylestrong{MS2} : la validation est effectuée par \sphinxcode{\sphinxupquote{ms2\_parsing}}. La page affiche
directement un tableau des scans parsés (\sphinxcode{\sphinxupquote{st.expander}},
colonnes \sphinxcode{\sphinxupquote{Precursor\_MZ}}, \sphinxcode{\sphinxupquote{Formula}}, \sphinxcode{\sphinxupquote{Lipid\_Name}},
\sphinxcode{\sphinxupquote{FA\_composition}}, \sphinxcode{\sphinxupquote{Adduct}}, \sphinxcode{\sphinxupquote{Lipid\_category}}, \sphinxcode{\sphinxupquote{Lipid\_class}},
\sphinxcode{\sphinxupquote{Lipid\_subclass}}, \sphinxcode{\sphinxupquote{RT}}, \sphinxcode{\sphinxupquote{CCS}}, \sphinxcode{\sphinxupquote{Num\_Peaks}}).

\item {} 
\sphinxAtStartPar
\sphinxstylestrong{MS1} : contrôle, dans l’ordre, la présence de toutes les colonnes de
\sphinxcode{\sphinxupquote{REQUIRED\_COLUMNS}}, l’absence de cellules vides dans
\sphinxcode{\sphinxupquote{NON\_EMPTY\_COLUMNS}}, puis la validité numérique de la colonne
\sphinxcode{\sphinxupquote{Precursor\_MZ}} (\sphinxcode{\sphinxupquote{pandas.to\_numeric()}}, conversion en \sphinxcode{\sphinxupquote{float}}).

\end{itemize}

\sphinxAtStartPar
En cas d’erreur, un message détaillé est affiché et la page s’arrête.
L’utilisateur doit corriger puis recharger son fichier.

\item {} 
\sphinxAtStartPar
\sphinxstylestrong{Completion}
Au clic sur \sphinxstylestrong{Run automatic completion}, calcul des colonnes dérivées :
\begin{itemize}
\item {} 
\sphinxAtStartPar
\sphinxcode{\sphinxupquote{Molecular\_weight}} (via {\hyperref[\detokenize{api:utils.molecular_weight.molecular_weight}]{\sphinxcrossref{\sphinxcode{\sphinxupquote{utils.molecular\_weight.molecular\_weight()}}}}})
et \sphinxcode{\sphinxupquote{Monoisotopic\_mass}} (via {\hyperref[\detokenize{api:utils.monoisotopic.monoisotopic_mass}]{\sphinxcrossref{\sphinxcode{\sphinxupquote{utils.monoisotopic.monoisotopic\_mass()}}}}}),
à partir de \sphinxcode{\sphinxupquote{Formula}}. Si une formule est invalide, la page affiche une erreur et s’arrête sans stocker le
résultat.

\item {} 
\sphinxAtStartPar
\sphinxcode{\sphinxupquote{Adduct}} et \sphinxcode{\sphinxupquote{Neutral\_mass}}, résolus pour toutes les lignes/tous les
scans d’un coup par la fonction interne \sphinxcode{\sphinxupquote{\_adducts()}} (voir
{\hyperref[\detokenize{1_Integration:fonctions-internes}]{\sphinxcrossref{Fonctions internes}}}). Si un ou plusieurs adduits sont invalides, la
page liste chaque ligne/scan en erreur et s’arrête sans stocker les
résultats, même ceux valides. L’utilisateur doit corriger les adduits dans le fichier et recharger.

\item {} 
\sphinxAtStartPar
Pour \sphinxcode{\sphinxupquote{MS1}} uniquement : \sphinxcode{\sphinxupquote{MS\_level}} (valeur fixe \sphinxcode{\sphinxupquote{ms\_level}}),
\sphinxcode{\sphinxupquote{Num\_Peaks}} (fixé à \sphinxcode{\sphinxupquote{0}}), \sphinxcode{\sphinxupquote{Ionisation\_mode}} (valeur fixe
\sphinxcode{\sphinxupquote{ion\_mode}}). Pour \sphinxcode{\sphinxupquote{MS2}}, seul \sphinxcode{\sphinxupquote{Num\_Peaks}} provient réellement du
fichier parsé (nombre de fragments de chaque scan, issu du champ
\sphinxcode{\sphinxupquote{Num\_peaks}} de son bloc). \sphinxcode{\sphinxupquote{Ionisation\_mode}} reste fixé à la valeur
choisie en étape 1 (\sphinxcode{\sphinxupquote{ion\_mode}}), exactement comme pour MS1.
\sphinxcode{\sphinxupquote{MS\_level}} vaut simplement \sphinxcode{\sphinxupquote{"MS2"}}.

\end{itemize}

\sphinxAtStartPar
Le tableau résultat (\sphinxcode{\sphinxupquote{df\_complete}}) est ensuite modifiable via
\sphinxcode{\sphinxupquote{st.data\_editor}}. Les colonnes calculées ou fixées (\sphinxcode{\sphinxupquote{Precursor\_MZ}},
\sphinxcode{\sphinxupquote{Neutral\_mass}}, \sphinxcode{\sphinxupquote{Molecular\_weight}}, \sphinxcode{\sphinxupquote{Monoisotopic\_mass}}, \sphinxcode{\sphinxupquote{Formula}},
\sphinxcode{\sphinxupquote{MS\_level}}, \sphinxcode{\sphinxupquote{Num\_Peaks}}, \sphinxcode{\sphinxupquote{Ionisation\_mode}}, \sphinxcode{\sphinxupquote{Adduct}}) sont en
lecture seule. Restent donc modifiables :
\begin{itemize}
\item {} 
\sphinxAtStartPar
\sphinxstylestrong{MS1} : \sphinxcode{\sphinxupquote{Lipid\_Name}}, \sphinxcode{\sphinxupquote{Lipid\_category}}, \sphinxcode{\sphinxupquote{Lipid\_class}},
\sphinxcode{\sphinxupquote{Lipid\_subclass}}, \sphinxcode{\sphinxupquote{RT}}, \sphinxcode{\sphinxupquote{CCS}} ;

\item {} 
\sphinxAtStartPar
\sphinxstylestrong{MS2} : \sphinxcode{\sphinxupquote{Lipid\_Name}}, \sphinxcode{\sphinxupquote{FA\_composition}}, \sphinxcode{\sphinxupquote{Lipid\_category}},
\sphinxcode{\sphinxupquote{Lipid\_class}}, \sphinxcode{\sphinxupquote{Lipid\_subclass}}, \sphinxcode{\sphinxupquote{RT}}, \sphinxcode{\sphinxupquote{CCS}}.

\end{itemize}

\sphinxAtStartPar
Le \sphinxcode{\sphinxupquote{scan\_id}} de chaque bloc MS2 (voir {\hyperref[\detokenize{1_Integration:colonnes-du-fichier-d-entree}]{\sphinxcrossref{Colonnes du fichier d’entrée}}})
n’apparaît pas dans ce tableau : il reste uniquement dans
\sphinxcode{\sphinxupquote{st.session\_state{[}"ms2"{]}}} et n’est utilisé qu’en interne, à l’étape 5,
pour associer chaque ligne éditée au scan correspondant par position.

\sphinxAtStartPar
Le nombre de lignes est figé pour MS2 (\sphinxcode{\sphinxupquote{num\_rows="fixed"}}, un scan = une
ligne) et modifiable pour MS1 (\sphinxcode{\sphinxupquote{num\_rows="dynamic"}}, ajout/suppression de
lignes possible).

\sphinxAtStartPar
Le bouton \sphinxstylestrong{Validate data} revérifie l’absence de cellules vides dans
\sphinxcode{\sphinxupquote{NON\_EMPTY\_COLUMNS}} sur les données éditées avant de les stocker dans
\sphinxcode{\sphinxupquote{df\_validated}}. Une fois validées, un résumé est affiché : métriques
(nombre total de lipides, de catégories, classes et sous\sphinxhyphen{}classes) et trois
histogrammes horizontaux (\sphinxcode{\sphinxupquote{plotly.express.bar()}}, couleur selon
\sphinxcode{\sphinxupquote{CHART\_COLOR\_SCALE}}) répartissant les lipides par catégorie, classe et
sous\sphinxhyphen{}classe.

\item {} 
\sphinxAtStartPar
\sphinxstylestrong{Integration}
Insertion des données validées dans les tables \sphinxcode{\sphinxupquote{Detection}}, \sphinxcode{\sphinxupquote{Lipid}},
\sphinxcode{\sphinxupquote{Annotation}} (et \sphinxcode{\sphinxupquote{Fragment}} pour MS2) :
\begin{itemize}
\item {} 
\sphinxAtStartPar
\sphinxstylestrong{MS1} : {\hyperref[\detokenize{api:utils.loading_MS1.DB_MS1}]{\sphinxcrossref{\sphinxcode{\sphinxupquote{utils.loading\_MS1.DB\_MS1()}}}}} reçoit directement le
DataFrame validé (\sphinxcode{\sphinxupquote{df\_validated}}).

\item {} 
\sphinxAtStartPar
\sphinxstylestrong{MS2} : le tableau \sphinxcode{\sphinxupquote{df\_validated}} ne contient que les colonnes de
métadonnées affichées à l’étape 4. Les fragments (m/z/intensité) de
chaque scan, eux, n’existent que dans \sphinxcode{\sphinxupquote{st.session\_state{[}"ms2"{]}}} (liste
de dictionnaires produite par le parsing) : ils n’apparaissent jamais
dans ce tableau.

\sphinxAtStartPar
Or \sphinxcode{\sphinxupquote{utils.loading\_MS2.DB\_MS2()}} a besoin des dictionnaires complets
pour insérer les fragments en base. La page reporte donc d’abord les
valeurs de \sphinxcode{\sphinxupquote{df\_validated}} (potentiellement modifiées par
l’utilisateur) dans les dictionnaires de \sphinxcode{\sphinxupquote{st.session\_state{[}"ms2"{]}}}.
Chaque ligne du tableau est associée au scan de même position
(\sphinxcode{\sphinxupquote{scans{[}pos{]}}}). Cette correspondance par position est fiable car
\sphinxcode{\sphinxupquote{num\_rows="fixed"}} empêche tout ajout ou suppression de ligne à
l’étape 4 : l’ordre reste donc identique entre les deux structures.

\sphinxAtStartPar
Seuls les champs présents dans le tableau (\sphinxcode{\sphinxupquote{lipid\_name}},
\sphinxcode{\sphinxupquote{fa\_composition}}, \sphinxcode{\sphinxupquote{lipid\_category}}, \sphinxcode{\sphinxupquote{lipid\_class}},
\sphinxcode{\sphinxupquote{lipid\_subclass}}, \sphinxcode{\sphinxupquote{rt}}, \sphinxcode{\sphinxupquote{ccs}}, \sphinxcode{\sphinxupquote{molecular\_weight}},
\sphinxcode{\sphinxupquote{monoisotopic\_mass}}, \sphinxcode{\sphinxupquote{neutral\_mass}}, \sphinxcode{\sphinxupquote{adduct}}, \sphinxcode{\sphinxupquote{ion\_mode}}) sont
ainsi écrasés dans chaque dictionnaire de scan. Les fragments, absents
du tableau, restent donc inchangés. \sphinxcode{\sphinxupquote{utils.loading\_MS2.DB\_MS2()}}
est ensuite appelée avec cette liste de scans mise à jour.

\end{itemize}

\sphinxAtStartPar
Dans les deux cas, la fonction reçoit également le nom du fichier importé,
\sphinxcode{\sphinxupquote{integrator\_name}} et \sphinxcode{\sphinxupquote{file\_row\_name}}, et se charge en interne de
l’historique et de la sauvegarde de la base. En cas de succès,
\sphinxcode{\sphinxupquote{integration\_done}} passe à \sphinxcode{\sphinxupquote{True}} et des ballons (\sphinxcode{\sphinxupquote{st.balloons}})
sont affichés une seule fois (flag \sphinxcode{\sphinxupquote{show\_balloons}} consommé via
\sphinxcode{\sphinxupquote{st.session\_state.pop}}). Une fois l’intégration finie, une nouvelle intégration est proposée
(bouton \sphinxstylestrong{Reset} de la barre latérale).

\end{enumerate}


\section{État de session (\sphinxstyleliteralintitle{\sphinxupquote{st.session\_state}})}
\label{\detokenize{1_Integration:etat-de-session-st-session-state}}

\begin{savenotes}\sphinxattablestart
\sphinxthistablewithglobalstyle
\centering
\begin{tabular}[t]{\X{20}{100}\X{20}{100}\X{60}{100}}
\sphinxtoprule
\sphinxstyletheadfamily 
\sphinxAtStartPar
Constante
&\sphinxstyletheadfamily 
\sphinxAtStartPar
Clé \sphinxcode{\sphinxupquote{session\_state}}
&\sphinxstyletheadfamily 
\sphinxAtStartPar
Description
\\
\sphinxmidrule
\sphinxtableatstartofbodyhook
\sphinxAtStartPar
—
&
\sphinxAtStartPar
\sphinxcode{\sphinxupquote{ms\_level}}, \sphinxcode{\sphinxupquote{ion\_mode}}, \sphinxcode{\sphinxupquote{integrator\_name}}, \sphinxcode{\sphinxupquote{file\_row\_name}}
&
\sphinxAtStartPar
Paramètres choisis à l’étape 1 (clés des widgets correspondants)
\\
\sphinxhline
\sphinxAtStartPar
\sphinxcode{\sphinxupquote{MS1}}
&
\sphinxAtStartPar
\sphinxcode{\sphinxupquote{ms1}}
&
\sphinxAtStartPar
DataFrame brut chargé à l’étape 2 (MS1 uniquement)
\\
\sphinxhline
\sphinxAtStartPar
\sphinxcode{\sphinxupquote{MS2}}
&
\sphinxAtStartPar
\sphinxcode{\sphinxupquote{ms2}}
&
\sphinxAtStartPar
Liste de scans parsés à l’étape 2 (MS2 uniquement)
\\
\sphinxhline
\sphinxAtStartPar
\sphinxcode{\sphinxupquote{DF\_ID}}
&
\sphinxAtStartPar
\sphinxcode{\sphinxupquote{df\_file\_id}}
&
\sphinxAtStartPar
Identifiant du fichier chargé, utilisé pour détecter un nouvel upload
\\
\sphinxhline
\sphinxAtStartPar
\sphinxcode{\sphinxupquote{DF\_MS\_LEVEL}}
&
\sphinxAtStartPar
\sphinxcode{\sphinxupquote{df\_ms\_level}}
&
\sphinxAtStartPar
Niveau MS utilisé lors du dernier parsing, pour détecter un changement de niveau
\\
\sphinxhline
\sphinxAtStartPar
\sphinxcode{\sphinxupquote{COLUMNS\_VALID}}
&
\sphinxAtStartPar
\sphinxcode{\sphinxupquote{columns\_valid}}
&
\sphinxAtStartPar
Booléen indiquant que la vérification des colonnes obligatoires (étape 3) a réussi
\\
\sphinxhline
\sphinxAtStartPar
\sphinxcode{\sphinxupquote{DF\_COMPLETE}}
&
\sphinxAtStartPar
\sphinxcode{\sphinxupquote{df\_complete}}
&
\sphinxAtStartPar
DataFrame avec les colonnes dérivées calculées (étape 4)
\\
\sphinxhline
\sphinxAtStartPar
\sphinxcode{\sphinxupquote{DF\_VALID}}
&
\sphinxAtStartPar
\sphinxcode{\sphinxupquote{df\_validated}}
&
\sphinxAtStartPar
DataFrame validé par l’utilisateur, prêt pour l’intégration
\\
\sphinxhline
\sphinxAtStartPar
\sphinxcode{\sphinxupquote{INTEGRATION\_DONE}}
&
\sphinxAtStartPar
\sphinxcode{\sphinxupquote{integration\_done}}
&
\sphinxAtStartPar
Booléen indiquant que l’intégration en base a été effectuée
\\
\sphinxhline
\sphinxAtStartPar
\sphinxcode{\sphinxupquote{EDITOR}}
&
\sphinxAtStartPar
\sphinxcode{\sphinxupquote{editor\_integration}}
&
\sphinxAtStartPar
Clé du widget \sphinxcode{\sphinxupquote{st.data\_editor}} utilisé à l’étape 4
\\
\sphinxhline
\sphinxAtStartPar
\sphinxcode{\sphinxupquote{UPLOADER\_VERSION}}
&
\sphinxAtStartPar
\sphinxcode{\sphinxupquote{uploader\_version}}
&
\sphinxAtStartPar
Compteur incrémenté par le bouton \sphinxstylestrong{Reset} pour recréer un \sphinxcode{\sphinxupquote{st.file\_uploader}} vide
\\
\sphinxbottomrule
\end{tabular}
\sphinxtableafterendhook\par
\sphinxattableend\end{savenotes}

\sphinxAtStartPar
\sphinxcode{\sphinxupquote{RESET\_KEYS}} (utilisé par le bouton \sphinxstylestrong{Reset}) regroupe l’ensemble des clés
calculées (\sphinxcode{\sphinxupquote{MS1}}, \sphinxcode{\sphinxupquote{MS2}}, \sphinxcode{\sphinxupquote{DF\_ID}}, \sphinxcode{\sphinxupquote{DF\_MS\_LEVEL}}, \sphinxcode{\sphinxupquote{DF\_COMPLETE}},
\sphinxcode{\sphinxupquote{DF\_VALID}}, \sphinxcode{\sphinxupquote{INTEGRATION\_DONE}}, \sphinxcode{\sphinxupquote{COLUMNS\_VALID}}, \sphinxcode{\sphinxupquote{EDITOR}}) ;
\sphinxcode{\sphinxupquote{RELOAD\_RESET\_KEYS}} n’en contient qu’un sous\sphinxhyphen{}ensemble (\sphinxcode{\sphinxupquote{DF\_COMPLETE}},
\sphinxcode{\sphinxupquote{DF\_VALID}}, \sphinxcode{\sphinxupquote{INTEGRATION\_DONE}}, \sphinxcode{\sphinxupquote{COLUMNS\_VALID}}, \sphinxcode{\sphinxupquote{EDITOR}}), utilisé
lors du chargement d’un nouveau fichier (ou d’un changement de niveau MS) pour
conserver les réglages de l’étape 1 et le fichier tout juste chargé.


\section{Fonctions internes}
\label{\detokenize{1_Integration:fonctions-internes}}

\begin{savenotes}\sphinxattablestart
\sphinxthistablewithglobalstyle
\centering
\begin{tabular}[t]{\X{30}{100}\X{70}{100}}
\sphinxtoprule
\sphinxstyletheadfamily 
\sphinxAtStartPar
Fonction
&\sphinxstyletheadfamily 
\sphinxAtStartPar
Description
\\
\sphinxmidrule
\sphinxtableatstartofbodyhook
\sphinxAtStartPar
\sphinxcode{\sphinxupquote{\_workflow(done: bool) \sphinxhyphen{}\textgreater{} str}}
&
\sphinxAtStartPar
Retourne l’icône \sphinxcode{\sphinxupquote{"✅"}} si \sphinxcode{\sphinxupquote{done}} est vrai, sinon \sphinxcode{\sphinxupquote{"⬜"}}. Utilisée
pour afficher l’état de chaque étape dans la barre latérale.
\\
\sphinxhline
\sphinxAtStartPar
\sphinxcode{\sphinxupquote{\_adducts(records) \sphinxhyphen{}\textgreater{} tuple{[}list, list, list{[}str{]}{]}}}
&
\sphinxAtStartPar
Résout, pour un lot de lignes/scans, l’adduit standardisé et la masse
neutre correspondants (voir détail ci\sphinxhyphen{}dessous).
\\
\sphinxbottomrule
\end{tabular}
\sphinxtableafterendhook\par
\sphinxattableend\end{savenotes}


\subsection{\sphinxstyleliteralintitle{\sphinxupquote{\_adducts(records)}}}
\label{\detokenize{1_Integration:adducts-records}}
\sphinxAtStartPar
Résout l’adduit et la masse neutre pour un lot d’enregistrements
\sphinxcode{\sphinxupquote{(label, lipid\_name, precursor\_mz, raw\_adduct)}}, en s’appuyant sur
{\hyperref[\detokenize{api:utils.neutral_mass.adduct}]{\sphinxcrossref{\sphinxcode{\sphinxupquote{utils.neutral\_mass.adduct()}}}}} (normalisation/validation de l’adduit) puis
{\hyperref[\detokenize{api:utils.neutral_mass.neutral_mass}]{\sphinxcrossref{\sphinxcode{\sphinxupquote{utils.neutral\_mass.neutral\_mass()}}}}} (calcul de la masse neutre à partir du
\sphinxcode{\sphinxupquote{Precursor\_MZ}} et de l’adduit résolu).
\begin{description}
\sphinxlineitem{Paramètres}\begin{description}
\sphinxlineitem{\sphinxcode{\sphinxupquote{records}}}
\sphinxAtStartPar
Itérable de tuples \sphinxcode{\sphinxupquote{(label, lipid\_name, precursor\_mz, raw\_adduct)}} :
\begin{itemize}
\item {} 
\sphinxAtStartPar
\sphinxcode{\sphinxupquote{label}} (\sphinxstyleemphasis{str}) : identifiant lisible de la ligne/du scan pour les
messages d’erreur (ex. \sphinxcode{\sphinxupquote{"Row 3"}} pour MS1, \sphinxcode{\sphinxupquote{"Scan 12"}} pour MS2) ;

\item {} 
\sphinxAtStartPar
\sphinxcode{\sphinxupquote{lipid\_name}} (\sphinxstyleemphasis{str}) : nom du lipide, ré\sphinxhyphen{}inséré dans le message
d’erreur pour faciliter la correction du fichier ;

\item {} 
\sphinxAtStartPar
\sphinxcode{\sphinxupquote{precursor\_mz}} (\sphinxstyleemphasis{float}) : m/z du précurseur, transmis à
{\hyperref[\detokenize{api:utils.neutral_mass.neutral_mass}]{\sphinxcrossref{\sphinxcode{\sphinxupquote{utils.neutral\_mass.neutral\_mass()}}}}} ;

\item {} 
\sphinxAtStartPar
\sphinxcode{\sphinxupquote{raw\_adduct}} (\sphinxstyleemphasis{str} ou \sphinxstyleemphasis{None}) : valeur brute de l’adduit telle que
lue dans le fichier, avant normalisation.

\end{itemize}

\end{description}

\sphinxlineitem{Retour}
\sphinxAtStartPar
Un triplet \sphinxcode{\sphinxupquote{(adducts, masses, errors)}} de listes alignées sur
\sphinxcode{\sphinxupquote{records}} :
\begin{itemize}
\item {} 
\sphinxAtStartPar
\sphinxcode{\sphinxupquote{adducts}} : adduit normalisé (ex. \sphinxcode{\sphinxupquote{"{[}M+H{]}+"}}), ou \sphinxcode{\sphinxupquote{None}} si la
résolution a échoué pour cette entrée ;

\item {} 
\sphinxAtStartPar
\sphinxcode{\sphinxupquote{masses}} : masse neutre en Daltons (6 décimales), ou \sphinxcode{\sphinxupquote{None}} en cas
d’échec ;

\item {} 
\sphinxAtStartPar
\sphinxcode{\sphinxupquote{errors}} : liste des messages d’erreur (un par entrée en échec),
préformatés en \sphinxcode{\sphinxupquote{"\{label\} (\{lipid\_name\}) : \{message\}"}}, prête à être
affichée telle quelle par la page (une entrée par ligne dans
\sphinxcode{\sphinxupquote{st.error}}).

\end{itemize}

\sphinxAtStartPar
Une entrée échoue (adduit/masse à \sphinxcode{\sphinxupquote{None}}, message ajouté à \sphinxcode{\sphinxupquote{errors}}) si
{\hyperref[\detokenize{api:utils.neutral_mass.adduct}]{\sphinxcrossref{\sphinxcode{\sphinxupquote{utils.neutral\_mass.adduct()}}}}} lève une \sphinxcode{\sphinxupquote{ValueError}} — adduit manquant
ou non reconnu (valeurs acceptées : \sphinxcode{\sphinxupquote{"{[}M+H{]}+"}}, \sphinxcode{\sphinxupquote{"{[}M+NH4{]}+"}},
\sphinxcode{\sphinxupquote{"{[}M\sphinxhyphen{}H{]}\sphinxhyphen{}"}}). Les autres entrées du lot continuent d’être traitées : un
seul appel à \sphinxcode{\sphinxupquote{\_adducts}} permet donc de collecter \sphinxstyleemphasis{toutes} les erreurs
d’adduit du fichier en une fois, plutôt que de s’arrêter à la première.

\sphinxlineitem{Exemple}
\begin{sphinxVerbatim}[commandchars=\\\{\}]
\PYG{g+gp}{\PYGZgt{}\PYGZgt{}\PYGZgt{} }\PYG{n}{records} \PYG{o}{=} \PYG{p}{[}
\PYG{g+gp}{... }    \PYG{p}{(}\PYG{l+s+s2}{\PYGZdq{}}\PYG{l+s+s2}{Row 0}\PYG{l+s+s2}{\PYGZdq{}}\PYG{p}{,} \PYG{l+s+s2}{\PYGZdq{}}\PYG{l+s+s2}{PE 34:1}\PYG{l+s+s2}{\PYGZdq{}}\PYG{p}{,} \PYG{l+m+mf}{700.500000}\PYG{p}{,} \PYG{l+s+s2}{\PYGZdq{}}\PYG{l+s+s2}{[M+H]+}\PYG{l+s+s2}{\PYGZdq{}}\PYG{p}{)}\PYG{p}{,}
\PYG{g+gp}{... }    \PYG{p}{(}\PYG{l+s+s2}{\PYGZdq{}}\PYG{l+s+s2}{Row 1}\PYG{l+s+s2}{\PYGZdq{}}\PYG{p}{,} \PYG{l+s+s2}{\PYGZdq{}}\PYG{l+s+s2}{PC 32:0}\PYG{l+s+s2}{\PYGZdq{}}\PYG{p}{,} \PYG{l+m+mf}{750.500000}\PYG{p}{,} \PYG{l+s+s2}{\PYGZdq{}}\PYG{l+s+s2}{XYZ}\PYG{l+s+s2}{\PYGZdq{}}\PYG{p}{)}\PYG{p}{,}   \PYG{c+c1}{\PYGZsh{} adduit invalide}
\PYG{g+gp}{... }\PYG{p}{]}
\PYG{g+gp}{\PYGZgt{}\PYGZgt{}\PYGZgt{} }\PYG{n}{adducts}\PYG{p}{,} \PYG{n}{masses}\PYG{p}{,} \PYG{n}{errors} \PYG{o}{=} \PYG{n}{\PYGZus{}adducts}\PYG{p}{(}\PYG{n}{records}\PYG{p}{)}
\PYG{g+gp}{\PYGZgt{}\PYGZgt{}\PYGZgt{} }\PYG{n}{adducts}
\PYG{g+go}{[\PYGZsq{}[M+H]+\PYGZsq{}, None]}
\PYG{g+gp}{\PYGZgt{}\PYGZgt{}\PYGZgt{} }\PYG{n}{masses}
\PYG{g+go}{[699.492724, None]}
\PYG{g+gp}{\PYGZgt{}\PYGZgt{}\PYGZgt{} }\PYG{n}{errors}
\PYG{g+go}{[\PYGZdq{}Row 1 (PC 32:0) : Unsupported adduct \PYGZsq{}XYZ\PYGZsq{}. Accepted values : [M+H]+, [M+NH4]+, [M\PYGZhy{}H]\PYGZhy{}.\PYGZdq{}]}
\end{sphinxVerbatim}

\sphinxAtStartPar
Dans l’étape 4 de la page, \sphinxcode{\sphinxupquote{records}} est construit différemment selon le
niveau MS :

\begin{sphinxVerbatim}[commandchars=\\\{\}]
\PYG{c+c1}{\PYGZsh{} MS1 : une ligne de DataFrame par enregistrement}
\PYG{n}{adducts}\PYG{p}{,} \PYG{n}{neutral\PYGZus{}masses}\PYG{p}{,} \PYG{n}{e\PYGZus{}adduct} \PYG{o}{=} \PYG{n}{\PYGZus{}adducts}\PYG{p}{(}
    \PYG{p}{(}\PYG{l+s+sa}{f}\PYG{l+s+s2}{\PYGZdq{}}\PYG{l+s+s2}{Row }\PYG{l+s+si}{\PYGZob{}}\PYG{n}{idx}\PYG{l+s+si}{\PYGZcb{}}\PYG{l+s+s2}{\PYGZdq{}}\PYG{p}{,} \PYG{n}{row}\PYG{p}{[}\PYG{l+s+s2}{\PYGZdq{}}\PYG{l+s+s2}{Lipid\PYGZus{}Name}\PYG{l+s+s2}{\PYGZdq{}}\PYG{p}{]}\PYG{p}{,} \PYG{n}{row}\PYG{p}{[}\PYG{l+s+s2}{\PYGZdq{}}\PYG{l+s+s2}{Precursor\PYGZus{}MZ}\PYG{l+s+s2}{\PYGZdq{}}\PYG{p}{]}\PYG{p}{,} \PYG{n}{row}\PYG{o}{.}\PYG{n}{get}\PYG{p}{(}\PYG{l+s+s2}{\PYGZdq{}}\PYG{l+s+s2}{Adduct}\PYG{l+s+s2}{\PYGZdq{}}\PYG{p}{)}\PYG{p}{)}
    \PYG{k}{for} \PYG{n}{idx}\PYG{p}{,} \PYG{n}{row} \PYG{o+ow}{in} \PYG{n}{df}\PYG{o}{.}\PYG{n}{iterrows}\PYG{p}{(}\PYG{p}{)}
\PYG{p}{)}

\PYG{c+c1}{\PYGZsh{} MS2 : un scan par enregistrement}
\PYG{n}{adducts}\PYG{p}{,} \PYG{n}{n\PYGZus{}mass}\PYG{p}{,} \PYG{n}{e\PYGZus{}adduct} \PYG{o}{=} \PYG{n}{\PYGZus{}adducts}\PYG{p}{(}
    \PYG{p}{(}\PYG{l+s+sa}{f}\PYG{l+s+s2}{\PYGZdq{}}\PYG{l+s+s2}{Scan }\PYG{l+s+si}{\PYGZob{}}\PYG{n}{data}\PYG{p}{[}\PYG{l+s+s1}{\PYGZsq{}}\PYG{l+s+s1}{scan\PYGZus{}id}\PYG{l+s+s1}{\PYGZsq{}}\PYG{p}{]}\PYG{l+s+si}{\PYGZcb{}}\PYG{l+s+s2}{\PYGZdq{}}\PYG{p}{,} \PYG{n}{data}\PYG{p}{[}\PYG{l+s+s2}{\PYGZdq{}}\PYG{l+s+s2}{lipid\PYGZus{}name}\PYG{l+s+s2}{\PYGZdq{}}\PYG{p}{]}\PYG{p}{,} \PYG{n}{data}\PYG{p}{[}\PYG{l+s+s2}{\PYGZdq{}}\PYG{l+s+s2}{precursor\PYGZus{}mz}\PYG{l+s+s2}{\PYGZdq{}}\PYG{p}{]}\PYG{p}{,} \PYG{n}{data}\PYG{o}{.}\PYG{n}{get}\PYG{p}{(}\PYG{l+s+s2}{\PYGZdq{}}\PYG{l+s+s2}{adduct}\PYG{l+s+s2}{\PYGZdq{}}\PYG{p}{)}\PYG{p}{)}
    \PYG{k}{for} \PYG{n}{data} \PYG{o+ow}{in} \PYG{n}{ms2}
\PYG{p}{)}
\end{sphinxVerbatim}

\end{description}


\section{Dépendances internes}
\label{\detokenize{1_Integration:dependances-internes}}\begin{itemize}
\item {} 
\sphinxAtStartPar
{\hyperref[\detokenize{api:module-utils.loading_MS1}]{\sphinxcrossref{\sphinxcode{\sphinxupquote{utils.loading\_MS1}}}}}

\item {} 
\sphinxAtStartPar
\sphinxcode{\sphinxupquote{utils.loading\_MS2}}

\item {} 
\sphinxAtStartPar
{\hyperref[\detokenize{api:module-utils.molecular_weight}]{\sphinxcrossref{\sphinxcode{\sphinxupquote{utils.molecular\_weight}}}}}

\item {} 
\sphinxAtStartPar
{\hyperref[\detokenize{api:module-utils.monoisotopic}]{\sphinxcrossref{\sphinxcode{\sphinxupquote{utils.monoisotopic}}}}}

\item {} 
\sphinxAtStartPar
{\hyperref[\detokenize{api:module-utils.neutral_mass}]{\sphinxcrossref{\sphinxcode{\sphinxupquote{utils.neutral\_mass}}}}}

\end{itemize}

\sphinxstepscope


\chapter{BacLipidDB}
\label{\detokenize{2_BacLipidDB:baclipiddb}}\label{\detokenize{2_BacLipidDB::doc}}
\sphinxAtStartPar
Fichier source : \sphinxcode{\sphinxupquote{app/pages/2\_BacLipidDB.py}}

\begin{sphinxadmonition}{note}{Note:}
\sphinxAtStartPar
Cette page contient la documentation du module Streamlit (\sphinxcode{\sphinxupquote{2\_BacLipidDB.py}}).
La majorité du code n’est pas écrit sous forme de fonction, il n’est donc pas possible d’utiliser \sphinxcode{\sphinxupquote{automodule}}
pour générer cette page automatiquement. La documentation ci\sphinxhyphen{}dessous est rédigée manuellement
à partir du code source.
\end{sphinxadmonition}


\section{Vue d’ensemble}
\label{\detokenize{2_BacLipidDB:vue-d-ensemble}}
\sphinxAtStartPar
Ce module Streamlit (\sphinxstylestrong{MODULE 2}) permet d’explorer le contenu de la base de données
\sphinxstylestrong{BacLipidDB}, via un sélecteur (\sphinxcode{\sphinxupquote{st.radio}}) proposant :
\begin{itemize}
\item {} 
\sphinxAtStartPar
une vue des 3 tables de la base de données : \sphinxcode{\sphinxupquote{Detection}}, \sphinxcode{\sphinxupquote{Fragment}}, \sphinxcode{\sphinxupquote{Lipid}} ;

\item {} 
\sphinxAtStartPar
\sphinxstylestrong{Full view} : une vue jointe combinant les champs des tables \sphinxcode{\sphinxupquote{Detection\textasciigrave{}\textasciigrave{}et \textasciigrave{}\textasciigrave{}Lipid}},
colorée par lot d’import et accompagnée de 4 graphiques ;

\item {} 
\sphinxAtStartPar
\sphinxstylestrong{History} : l’historique des imports passés.

\end{itemize}


\section{Tables brutes}
\label{\detokenize{2_BacLipidDB:tables-brutes}}
\sphinxAtStartPar
Pour \sphinxcode{\sphinxupquote{Detection}}, \sphinxcode{\sphinxupquote{Fragment}} et \sphinxcode{\sphinxupquote{Lipid}}, le contenu de la table est lu
directement depuis BacLipidDB via \sphinxcode{\sphinxupquote{pandas.read\_sql\_table()}} puis affiché avec
\sphinxcode{\sphinxupquote{st.dataframe()}}. Les colonnes numériques suivantes sont formatées via
\sphinxcode{\sphinxupquote{st.column\_config.NumberColumn}} :
\begin{itemize}
\item {} 
\sphinxAtStartPar
\sphinxcode{\sphinxupquote{Precursor\_MZ}}, \sphinxcode{\sphinxupquote{Neutral\_mass}}, \sphinxcode{\sphinxupquote{MZ}}, \sphinxcode{\sphinxupquote{Monoisotopic\_mass}}, \sphinxcode{\sphinxupquote{Intensity}} : 6 décimales

\item {} 
\sphinxAtStartPar
\sphinxcode{\sphinxupquote{RT}}, \sphinxcode{\sphinxupquote{CCS}} : 4 décimales

\item {} 
\sphinxAtStartPar
\sphinxcode{\sphinxupquote{Molecular\_weight}} : 0 décimale

\end{itemize}

\sphinxAtStartPar
Si la table ne contient aucune ligne, un message \sphinxcode{\sphinxupquote{st.info("This table contains no data yet.")}}
est affiché à la place.


\section{Full view}
\label{\detokenize{2_BacLipidDB:full-view}}

\subsection{Importation des données}
\label{\detokenize{2_BacLipidDB:importation-des-donnees}}
\sphinxAtStartPar
La fonction \sphinxcode{\sphinxupquote{\_load\_data(\_engine)}} s’appuie sur \sphinxcode{\sphinxupquote{utils.data\_access.load\_database()}}
(requête SQLAlchemy commune sur \sphinxcode{\sphinxupquote{Annotation}}, avec chargement des relations avec les tables
\sphinxcode{\sphinxupquote{lipid}} et \sphinxcode{\sphinxupquote{detection}}, partagée avec les pages \sphinxstylestrong{3\_Export} et \sphinxstylestrong{5\_Management} pour éviter de
dupliquer cette requête). Cette fonction ne conserve que les colonnes :
\sphinxcode{\sphinxupquote{Detection\_ID}}, \sphinxcode{\sphinxupquote{Lipid\_name}}, \sphinxcode{\sphinxupquote{Formula}}, \sphinxcode{\sphinxupquote{FA\_composition}}, \sphinxcode{\sphinxupquote{Lipid\_class}},
\sphinxcode{\sphinxupquote{Lipid\_subclass}}, \sphinxcode{\sphinxupquote{Lipid\_category}}, \sphinxcode{\sphinxupquote{Precursor\_MZ}}, \sphinxcode{\sphinxupquote{Adduct}}, \sphinxcode{\sphinxupquote{Neutral\_mass}},
\sphinxcode{\sphinxupquote{Molecular\_weight}}, \sphinxcode{\sphinxupquote{Monoisotopic\_mass}}, \sphinxcode{\sphinxupquote{MS\_level}}, \sphinxcode{\sphinxupquote{Ionisation\_mode}}, \sphinxcode{\sphinxupquote{Num\_Peaks}},
\sphinxcode{\sphinxupquote{RT}}, \sphinxcode{\sphinxupquote{CCS}}. Les autres clé primaires (\_ID) et étrangères (\_id) ne sont pas affichées pour question de lisibilité.


\begin{savenotes}\sphinxattablestart
\sphinxthistablewithglobalstyle
\centering
\sphinxcapstartof{table}
\sphinxthecaptionisattop
\sphinxcaption{Exemple}\label{\detokenize{2_BacLipidDB:id2}}
\sphinxaftertopcaption
\begin{tabular}[t]{\X{30}{100}\X{70}{100}}
\sphinxtoprule
\sphinxstyletheadfamily 
\sphinxAtStartPar
Appel
&\sphinxstyletheadfamily 
\sphinxAtStartPar
Résultat
\\
\sphinxmidrule
\sphinxtableatstartofbodyhook
\sphinxAtStartPar
\sphinxcode{\sphinxupquote{load\_database(engine)}}
&
\sphinxAtStartPar
DataFrame complet : une ligne par \sphinxcode{\sphinxupquote{Annotation}}, avec en plus \sphinxcode{\sphinxupquote{Annotation\_ID}} et
\sphinxcode{\sphinxupquote{Lipid\_ID}}
\\
\sphinxhline
\sphinxAtStartPar
\sphinxcode{\sphinxupquote{\_load\_data(engine)}}
&
\sphinxAtStartPar
Le même DataFrame, restreint aux 17 colonnes listées ci\sphinxhyphen{}dessus
\\
\sphinxbottomrule
\end{tabular}
\sphinxtableafterendhook\par
\sphinxattableend\end{savenotes}


\subsection{Colorisation par lot d’import}
\label{\detokenize{2_BacLipidDB:colorisation-par-lot-d-import}}
\sphinxAtStartPar
Chaque ligne du tableau \sphinxstylestrong{Full view} est colorée selon le fichier d’import (\sphinxcode{\sphinxupquote{History}})
auquel appartient sa \sphinxcode{\sphinxupquote{Detection\_ID}} :
\begin{itemize}
\item {} 
\sphinxAtStartPar
\sphinxcode{\sphinxupquote{utils.history.load\_history()}} charge l’historique des imports depuis \sphinxcode{\sphinxupquote{history.json}}
(\sphinxcode{\sphinxupquote{config.HISTORY\_PATH}}, lu indirectement par \sphinxcode{\sphinxupquote{load\_history}}) ;

\item {} 
\sphinxAtStartPar
\sphinxcode{\sphinxupquote{\_batch\_map(history)}} associe chaque \sphinxcode{\sphinxupquote{Detection\_ID}} à l’indice du lot
d’import auquel il appartient (champ \sphinxcode{\sphinxupquote{detection\_ids}} de chaque entrée de l’historique) ;

\item {} 
\sphinxAtStartPar
la fonction \sphinxcode{\sphinxupquote{\_row\_color(row)}} applique, via \sphinxcode{\sphinxupquote{df.style.apply(..., axis=1)}}, une couleur de
fond calculée par \sphinxcode{\sphinxupquote{\_color(index)}} à chaque ligne selon son lot.

\end{itemize}


\subsubsection{Exemple (\sphinxstyleliteralintitle{\sphinxupquote{\_batch\_map}})}
\label{\detokenize{2_BacLipidDB:exemple-batch-map}}
\begin{sphinxVerbatim}[commandchars=\\\{\}]
\PYG{g+gp}{\PYGZgt{}\PYGZgt{}\PYGZgt{} }\PYG{n}{history} \PYG{o}{=} \PYG{p}{[}
\PYG{g+gp}{... }    \PYG{p}{\PYGZob{}}\PYG{l+s+s2}{\PYGZdq{}}\PYG{l+s+s2}{filename}\PYG{l+s+s2}{\PYGZdq{}}\PYG{p}{:} \PYG{l+s+s2}{\PYGZdq{}}\PYG{l+s+s2}{batch1.csv}\PYG{l+s+s2}{\PYGZdq{}}\PYG{p}{,} \PYG{l+s+s2}{\PYGZdq{}}\PYG{l+s+s2}{detection\PYGZus{}ids}\PYG{l+s+s2}{\PYGZdq{}}\PYG{p}{:} \PYG{p}{[}\PYG{l+m+mi}{101}\PYG{p}{,} \PYG{l+m+mi}{102}\PYG{p}{,} \PYG{l+m+mi}{103}\PYG{p}{]}\PYG{p}{\PYGZcb{}}\PYG{p}{,}
\PYG{g+gp}{... }    \PYG{p}{\PYGZob{}}\PYG{l+s+s2}{\PYGZdq{}}\PYG{l+s+s2}{filename}\PYG{l+s+s2}{\PYGZdq{}}\PYG{p}{:} \PYG{l+s+s2}{\PYGZdq{}}\PYG{l+s+s2}{batch2.csv}\PYG{l+s+s2}{\PYGZdq{}}\PYG{p}{,} \PYG{l+s+s2}{\PYGZdq{}}\PYG{l+s+s2}{detection\PYGZus{}ids}\PYG{l+s+s2}{\PYGZdq{}}\PYG{p}{:} \PYG{p}{[}\PYG{l+m+mi}{104}\PYG{p}{,} \PYG{l+m+mi}{105}\PYG{p}{]}\PYG{p}{\PYGZcb{}}\PYG{p}{,}
\PYG{g+gp}{... }\PYG{p}{]}
\PYG{g+gp}{\PYGZgt{}\PYGZgt{}\PYGZgt{} }\PYG{n}{\PYGZus{}batch\PYGZus{}map}\PYG{p}{(}\PYG{n}{history}\PYG{p}{)}
\PYG{g+go}{\PYGZob{}101: 0, 102: 0, 103: 0, 104: 1, 105: 1\PYGZcb{}}
\end{sphinxVerbatim}

\sphinxAtStartPar
\sphinxcode{\sphinxupquote{\_color(index)}} ne repose pas sur une palette figée : la teinte (\sphinxcode{\sphinxupquote{hue}}) est calculée
par \sphinxcode{\sphinxupquote{(index * 0.618) \% 1.0}} (technique du \sphinxstyleemphasis{golden angle} avec le nombre d’or conjugé \sphinxstyleemphasis{0.618}) puis convertie en couleur pastel via \sphinxcode{\sphinxupquote{colorsys.hls\_to\_rgb()}} (luminosité 0.85,
saturation 0.55 fixes).


\subsubsection{Exemple (\sphinxstyleliteralintitle{\sphinxupquote{\_color(1)}})}
\label{\detokenize{2_BacLipidDB:exemple-color-1}}\begin{enumerate}
\sphinxsetlistlabels{\arabic}{enumi}{enumii}{}{.}%
\item {} 
\sphinxAtStartPar
\sphinxcode{\sphinxupquote{hue = (1 * 0.618) \% 1.0 = 0.618}}

\item {} 
\sphinxAtStartPar
\sphinxcode{\sphinxupquote{colorsys.hls\_to\_rgb()}} (0.618, 0.85, 0.55) convertit HLS en RGB (flottants dans {[}0, 1)) :
\sphinxcode{\sphinxupquote{r = 0.7675}}, \sphinxcode{\sphinxupquote{g = 0.8157}}, \sphinxcode{\sphinxupquote{b = 0.9325}}

\item {} 
\sphinxAtStartPar
chaque flottant est mis à l’échelle {[}0, 255{]} et arrondi à l’entier le plus proche :
\sphinxcode{\sphinxupquote{r = 196}}, \sphinxcode{\sphinxupquote{g = 208}}, \sphinxcode{\sphinxupquote{b = 238}}

\item {} 
\sphinxAtStartPar
chaque entier est formaté en paire hexadécimale sur 2 chiffres : \sphinxcode{\sphinxupquote{196 = c4}},
\sphinxcode{\sphinxupquote{208 = d0}}, \sphinxcode{\sphinxupquote{238 = ee}} \(\rightarrow\) \sphinxcode{\sphinxupquote{\#c4d0ee}}

\end{enumerate}

\sphinxAtStartPar
Un encadré \sphinxstylestrong{color caption} (\sphinxcode{\sphinxupquote{st.expander}}) rappelle, pour chaque fichier importé, sa
couleur, son nom, sa date d’import, son nombre de lignes et son intégrateur.

\begin{sphinxadmonition}{note}{Note:}
\sphinxAtStartPar
Cette coloration par lot ne concerne que les lignes du tableau \sphinxstylestrong{Full view}. Les graphiques 1 à 3 (masses),
plus bas, utilisent un code couleur différent, à 3 catégories fixes : voir {\hyperref[\detokenize{2_BacLipidDB:groupe-d-import}]{\sphinxcrossref{Groupe d’import}}}.
\end{sphinxadmonition}


\subsubsection{Exemple}
\label{\detokenize{2_BacLipidDB:exemple}}
\sphinxAtStartPar
Soit un historique de 2 imports, et \sphinxcode{\sphinxupquote{df}} (après \sphinxcode{\sphinxupquote{reset\_index}}) contenant 4 lignes. Pour
chaque ligne, \sphinxcode{\sphinxupquote{caption\_values{[}row.name{]}}} donne le lot d’import de sa \sphinxcode{\sphinxupquote{Detection\_ID}} :


\begin{savenotes}\sphinxattablestart
\sphinxthistablewithglobalstyle
\centering
\begin{tabular}[t]{\X{10}{90}\X{15}{90}\X{20}{90}\X{25}{90}\X{20}{90}}
\sphinxtoprule
\sphinxstyletheadfamily 
\sphinxAtStartPar
\sphinxcode{\sphinxupquote{row.name}}
&\sphinxstyletheadfamily 
\sphinxAtStartPar
\sphinxcode{\sphinxupquote{Detection\_ID}}
&\sphinxstyletheadfamily 
\sphinxAtStartPar
\sphinxcode{\sphinxupquote{caption\_values{[}row.name{]}}}
&\sphinxstyletheadfamily 
\sphinxAtStartPar
Lot d’import
&\sphinxstyletheadfamily 
\sphinxAtStartPar
Couleur appliquée
\\
\sphinxmidrule
\sphinxtableatstartofbodyhook
\sphinxAtStartPar
0
&
\sphinxAtStartPar
1
&
\sphinxAtStartPar
0.0
&
\sphinxAtStartPar
Import 1 (plus ancien)
&
\sphinxAtStartPar
\sphinxcode{\sphinxupquote{\#eec4c4}}
\\
\sphinxhline
\sphinxAtStartPar
1
&
\sphinxAtStartPar
2
&
\sphinxAtStartPar
1.0
&
\sphinxAtStartPar
Import 2 (dernier import)
&
\sphinxAtStartPar
\sphinxcode{\sphinxupquote{\#c4d0ee}}
\\
\sphinxhline
\sphinxAtStartPar
2
&
\sphinxAtStartPar
3
&
\sphinxAtStartPar
NaN
&
\sphinxAtStartPar
aucun (\sphinxcode{\sphinxupquote{Detection\_ID}} absente de l’historique)
&
\sphinxAtStartPar
aucune (cellule non colorée)
\\
\sphinxhline
\sphinxAtStartPar
3
&
\sphinxAtStartPar
4
&
\sphinxAtStartPar
0.0
&
\sphinxAtStartPar
Import 1 (plus ancien)
&
\sphinxAtStartPar
\sphinxcode{\sphinxupquote{\#eec4c4}}
\\
\sphinxbottomrule
\end{tabular}
\sphinxtableafterendhook\par
\sphinxattableend\end{savenotes}

\sphinxAtStartPar
Pour la ligne d’index 1, \sphinxcode{\sphinxupquote{\_row\_color}} lit \sphinxcode{\sphinxupquote{caption\_values{[}1{]} = 1.0}}, en déduit la couleur
du lot 1 via \sphinxcode{\sphinxupquote{\_color(1)}} et l’applique à toutes les cellules de cette ligne.

\begin{sphinxadmonition}{note}{Note:}
\sphinxAtStartPar
En usage normal, ce cas (\sphinxcode{\sphinxupquote{Detection\_ID}} absente de l’historique) ne devrait pas se
produire : chaque intégration (\sphinxcode{\sphinxupquote{1\_Integration.py}}) ajoute automatiquement une entrée à
\sphinxcode{\sphinxupquote{history.json}} listant les \sphinxcode{\sphinxupquote{Detection\_ID}} insérées. Il ne survient qu’en cas
d’incohérence entre la base et \sphinxcode{\sphinxupquote{history.json}}. La ligne 3 de l’exemple ci\sphinxhyphen{}dessus illustre le comportement du code face à un scénario inattendu.
\end{sphinxadmonition}


\subsection{Groupe d’import}
\label{\detokenize{2_BacLipidDB:groupe-d-import}}
\sphinxAtStartPar
Les graphiques 1 à 3 ne colorent pas les points par lot d’import individuel
comme le tableau \sphinxstylestrong{Full view}, mais par groupe, via une colonne \sphinxcode{\sphinxupquote{Import\_group}} calculée par
\sphinxcode{\sphinxupquote{\_group(batch\_id)}} à partir de \sphinxcode{\sphinxupquote{caption\_values}} et de l’indice du dernier lot
(\sphinxcode{\sphinxupquote{last\_batch\_id = len(history) \sphinxhyphen{} 1}}) :


\begin{savenotes}\sphinxattablestart
\sphinxthistablewithglobalstyle
\centering
\begin{tabular}[t]{\X{30}{100}\X{40}{100}\X{30}{100}}
\sphinxtoprule
\sphinxstyletheadfamily 
\sphinxAtStartPar
\sphinxcode{\sphinxupquote{\_group(batch\_id)}}
&\sphinxstyletheadfamily 
\sphinxAtStartPar
Condition
&\sphinxstyletheadfamily 
\sphinxAtStartPar
Couleur (\sphinxcode{\sphinxupquote{COLORS}})
\\
\sphinxmidrule
\sphinxtableatstartofbodyhook
\sphinxAtStartPar
\sphinxcode{\sphinxupquote{"Unknown"}}
&
\sphinxAtStartPar
\sphinxcode{\sphinxupquote{batch\_id}} est \sphinxcode{\sphinxupquote{NaN}} (\sphinxcode{\sphinxupquote{Detection\_ID}} absente de l’historique)
&
\sphinxAtStartPar
\sphinxcode{\sphinxupquote{\#898781}} (gris)
\\
\sphinxhline
\sphinxAtStartPar
\sphinxcode{\sphinxupquote{"Latest import"}}
&
\sphinxAtStartPar
\sphinxcode{\sphinxupquote{int(batch\_id) == last\_batch\_id}}
&
\sphinxAtStartPar
\sphinxcode{\sphinxupquote{\#e34948}} (rouge)
\\
\sphinxhline
\sphinxAtStartPar
\sphinxcode{\sphinxupquote{"Other imports"}}
&
\sphinxAtStartPar
tout autre \sphinxcode{\sphinxupquote{batch\_id}}
&
\sphinxAtStartPar
\sphinxcode{\sphinxupquote{\#2a78d6}} (bleu)
\\
\sphinxbottomrule
\end{tabular}
\sphinxtableafterendhook\par
\sphinxattableend\end{savenotes}


\subsubsection{Exemple}
\label{\detokenize{2_BacLipidDB:id1}}
\sphinxAtStartPar
Avec le même historique de 2 imports que ci\sphinxhyphen{}dessus (\sphinxcode{\sphinxupquote{last\_batch\_id = 1}}) :


\begin{savenotes}\sphinxattablestart
\sphinxthistablewithglobalstyle
\centering
\begin{tabular}[t]{\X{15}{85}\X{20}{85}\X{30}{85}\X{20}{85}}
\sphinxtoprule
\sphinxstyletheadfamily 
\sphinxAtStartPar
\sphinxcode{\sphinxupquote{Detection\_ID}}
&\sphinxstyletheadfamily 
\sphinxAtStartPar
\sphinxcode{\sphinxupquote{caption\_values}}
&\sphinxstyletheadfamily 
\sphinxAtStartPar
\sphinxcode{\sphinxupquote{\_group(...)}}
&\sphinxstyletheadfamily 
\sphinxAtStartPar
Couleur
\\
\sphinxmidrule
\sphinxtableatstartofbodyhook
\sphinxAtStartPar
1
&
\sphinxAtStartPar
0.0
&
\sphinxAtStartPar
\sphinxcode{\sphinxupquote{"Other imports"}}
&
\sphinxAtStartPar
\sphinxcode{\sphinxupquote{\#2a78d6}}
\\
\sphinxhline
\sphinxAtStartPar
2
&
\sphinxAtStartPar
1.0
&
\sphinxAtStartPar
\sphinxcode{\sphinxupquote{"Latest import"}}
&
\sphinxAtStartPar
\sphinxcode{\sphinxupquote{\#e34948}}
\\
\sphinxhline
\sphinxAtStartPar
3
&
\sphinxAtStartPar
NaN
&
\sphinxAtStartPar
\sphinxcode{\sphinxupquote{"Unknown"}}
&
\sphinxAtStartPar
\sphinxcode{\sphinxupquote{\#898781}}
\\
\sphinxbottomrule
\end{tabular}
\sphinxtableafterendhook\par
\sphinxattableend\end{savenotes}


\subsection{Graphique 1 \sphinxhyphen{} Masse mesurée vs masse théorique}
\label{\detokenize{2_BacLipidDB:graphique-1-masse-mesuree-vs-masse-theorique}}
\sphinxAtStartPar
Nuage de points (\sphinxcode{\sphinxupquote{plotly.express.scatter()}}) de \sphinxcode{\sphinxupquote{Neutral\_mass}} (masse neutre mesurée,
calculée à partir de \sphinxcode{\sphinxupquote{Precursor\_MZ}} et de l’adduit \sphinxcode{\sphinxupquote{Adduct}}) en fonction de \sphinxcode{\sphinxupquote{Monoisotopic\_mass}} (masse monoisotopique théorique
du lipide de référence calculée à partir de la formule chimique), coloré par \sphinxcode{\sphinxupquote{Import\_group}}
(voir \sphinxtitleref{Groupe d’import}). Le survol affiche \sphinxcode{\sphinxupquote{Lipid\_name}}, \sphinxcode{\sphinxupquote{Lipid\_class}}, l’écart
\sphinxcode{\sphinxupquote{Delta (Da)}} (\sphinxcode{\sphinxupquote{Neutral\_mass \sphinxhyphen{} Monoisotopic\_mass}}) et l’erreur \sphinxcode{\sphinxupquote{Error (ppm)}}
(\sphinxcode{\sphinxupquote{Delta / Monoisotopic\_mass * 1e6}}). Une droite \sphinxcode{\sphinxupquote{y = x}} (\sphinxcode{\sphinxupquote{plotly.graph\_objects.Scatter}},
trait pointillé) sert de référence visuelle pour l’écart entre les masses.


\subsection{Graphique 2 \sphinxhyphen{} Erreur de masse (ppm) vs masse neutre}
\label{\detokenize{2_BacLipidDB:graphique-2-erreur-de-masse-ppm-vs-masse-neutre}}
\sphinxAtStartPar
Nuage de points (\sphinxcode{\sphinxupquote{plotly.express.scatter()}}) de l’erreur relative \sphinxcode{\sphinxupquote{Error (ppm)}} (même
calcul que pour le graphique 1) en fonction de \sphinxcode{\sphinxupquote{Neutral\_mass}}, coloré par \sphinxcode{\sphinxupquote{Import\_group}}. Le
survol affiche \sphinxcode{\sphinxupquote{Lipid\_name}}, \sphinxcode{\sphinxupquote{Lipid\_class}} et \sphinxcode{\sphinxupquote{Error (ppm)}}.


\subsection{Graphique 3 \sphinxhyphen{} Erreur de masse (Da) vs masse neutre}
\label{\detokenize{2_BacLipidDB:graphique-3-erreur-de-masse-da-vs-masse-neutre}}
\sphinxAtStartPar
Même principe que le graphique 2, mais en ordonnée l’écart absolu \sphinxcode{\sphinxupquote{Delta (Da)}} (au lieu de
l’erreur relative \sphinxcode{\sphinxupquote{Error (ppm)}}), toujours en fonction de \sphinxcode{\sphinxupquote{Neutral\_mass}} et coloré par
\sphinxcode{\sphinxupquote{Import\_group}}. Le survol affiche \sphinxcode{\sphinxupquote{Lipid\_name}}, \sphinxcode{\sphinxupquote{Lipid\_class}} et \sphinxcode{\sphinxupquote{Delta (Da)}}.

\sphinxAtStartPar
Ce graphique est suivi d’un tableau \sphinxstylestrong{Mass error table} listant \sphinxcode{\sphinxupquote{Lipid\_name}}, \sphinxcode{\sphinxupquote{Formula}},
\sphinxcode{\sphinxupquote{Neutral\_mass}}, \sphinxcode{\sphinxupquote{Monoisotopic\_mass}}, \sphinxcode{\sphinxupquote{Delta (Da)}} et \sphinxcode{\sphinxupquote{Error (ppm)}}, trié par valeur
absolue de \sphinxcode{\sphinxupquote{Error (ppm)}} décroissante, pour faire ressortir en premier les annotations dont la
masse mesurée s’écarte le plus de la masse théorique.


\subsection{Graphique 4 \sphinxhyphen{} Diagramme de Van Krevelen}
\label{\detokenize{2_BacLipidDB:graphique-4-diagramme-de-van-krevelen}}
\sphinxAtStartPar
Nuage de points H/C en fonction de O/C, coloré par \sphinxcode{\sphinxupquote{Lipid\_class}}. Les ratios atomiques
sont calculés à partir de la colonne \sphinxcode{\sphinxupquote{Formula}} : \sphinxcode{\sphinxupquote{\_formula(formula)}} parse la formule
chimique (regex \sphinxcode{\sphinxupquote{({[}A\sphinxhyphen{}Z{]}{[}a\sphinxhyphen{}z{]}?)(\textbackslash{}d*)}}) en un dictionnaire élément \(\rightarrow\) nombre d’atomes, dont on
extrait les comptes de C, H et O (0 si absent). Seules les lignes avec au moins un atome de
carbone (\sphinxcode{\sphinxupquote{C \textgreater{} 0}}) sont conservées, afin d’éviter une division par zéro sur \sphinxcode{\sphinxupquote{H/C}} et \sphinxcode{\sphinxupquote{O/C}}.


\subsubsection{Exemple (parsing de formule)}
\label{\detokenize{2_BacLipidDB:exemple-parsing-de-formule}}

\begin{savenotes}\sphinxattablestart
\sphinxthistablewithglobalstyle
\centering
\begin{tabular}[t]{\X{25}{100}\X{30}{100}\X{10}{100}\X{10}{100}\X{10}{100}\X{15}{100}}
\sphinxtoprule
\sphinxstyletheadfamily 
\sphinxAtStartPar
\sphinxcode{\sphinxupquote{Formula}}
&\sphinxstyletheadfamily 
\sphinxAtStartPar
\sphinxcode{\sphinxupquote{\_formula(formula)}}
&\sphinxstyletheadfamily 
\sphinxAtStartPar
C
&\sphinxstyletheadfamily 
\sphinxAtStartPar
H
&\sphinxstyletheadfamily 
\sphinxAtStartPar
O
&\sphinxstyletheadfamily 
\sphinxAtStartPar
H/C, O/C
\\
\sphinxmidrule
\sphinxtableatstartofbodyhook
\sphinxAtStartPar
\sphinxcode{\sphinxupquote{C42H82NO8P}}
&
\sphinxAtStartPar
\sphinxcode{\sphinxupquote{\{"C": 42, "H": 82, "N": 1, "O": 8, "P": 1\}}}
&
\sphinxAtStartPar
42
&
\sphinxAtStartPar
82
&
\sphinxAtStartPar
8
&
\sphinxAtStartPar
1.952, 0.190
\\
\sphinxbottomrule
\end{tabular}
\sphinxtableafterendhook\par
\sphinxattableend\end{savenotes}


\section{History}
\label{\detokenize{2_BacLipidDB:history}}
\sphinxAtStartPar
Affiche, via \sphinxcode{\sphinxupquote{pandas.DataFrame()}} construit à partir de \sphinxcode{\sphinxupquote{utils.history.load\_history()}},
la liste des imports passés : nom de fichier, date (\sphinxcode{\sphinxupquote{pandas.to\_datetime()}}), niveau MS, mode
d’ionisation, nombre de lignes insérées et nom de l’intégrateur. Si
l’historique est vide, un message \sphinxcode{\sphinxupquote{st.info("No imports recorded yet.")}} est affiché à la place.


\section{Fonctions internes}
\label{\detokenize{2_BacLipidDB:fonctions-internes}}

\begin{savenotes}\sphinxattablestart
\sphinxthistablewithglobalstyle
\centering
\begin{tabular}[t]{\X{30}{100}\X{70}{100}}
\sphinxtoprule
\sphinxstyletheadfamily 
\sphinxAtStartPar
Fonction
&\sphinxstyletheadfamily 
\sphinxAtStartPar
Description
\\
\sphinxmidrule
\sphinxtableatstartofbodyhook
\sphinxAtStartPar
\sphinxcode{\sphinxupquote{\_load\_data(\_engine: Engine) \sphinxhyphen{}\textgreater{} pd.DataFrame}}
&
\sphinxAtStartPar
Charge la vue jointe complète via \sphinxcode{\sphinxupquote{utils.data\_access.load\_database()}} et n’en
conserve que certaines colonnes.
\\
\sphinxhline
\sphinxAtStartPar
\sphinxcode{\sphinxupquote{\_formula(formula: str) \sphinxhyphen{}\textgreater{} dict}}
&
\sphinxAtStartPar
Parse une formule chimique en un dictionnaire symbole d’élément \(\rightarrow\) nombre d’atomes.
\\
\sphinxhline
\sphinxAtStartPar
\sphinxcode{\sphinxupquote{\_batch\_map(history: list{[}dict{]}) \sphinxhyphen{}\textgreater{} dict{[}int, int{]}}}
&
\sphinxAtStartPar
Associe chaque \sphinxcode{\sphinxupquote{Detection\_ID}} à l’indice du lot d’import auquel il appartient
(utilisé pour la colorisation de \sphinxstylestrong{Full view}).
\\
\sphinxhline
\sphinxAtStartPar
\sphinxcode{\sphinxupquote{\_group(batch\_id: float) \sphinxhyphen{}\textgreater{} str}}
&
\sphinxAtStartPar
Classe le lot d’import d’une ligne dans l’une des 3 catégories \sphinxcode{\sphinxupquote{"Latest import"}},
\sphinxcode{\sphinxupquote{"Other imports"}} ou \sphinxcode{\sphinxupquote{"Unknown"}}, utilisées pour colorer les graphiques 1 à 3 (voir
{\hyperref[\detokenize{2_BacLipidDB:groupe-d-import}]{\sphinxcrossref{Groupe d’import}}}).
\\
\sphinxhline
\sphinxAtStartPar
\sphinxcode{\sphinxupquote{\_row\_color(row: pd.Series) \sphinxhyphen{}\textgreater{} list{[}str{]}}}
&
\sphinxAtStartPar
Callback de style (\sphinxcode{\sphinxupquote{df.style.apply(..., axis=1)}}) qui renvoie la couleur de fond CSS
à appliquer à chaque cellule d’une ligne de \sphinxstylestrong{Full view}, selon son lot d’import.
\\
\sphinxhline
\sphinxAtStartPar
\sphinxcode{\sphinxupquote{\_color(index: int) \sphinxhyphen{}\textgreater{} str}}
&
\sphinxAtStartPar
Génère une couleur pastel (hex) pour un indice de lot, en espaçant les teintes par le
nombre d’or conjugué (\sphinxstyleemphasis{golden angle}) plutôt que via une palette figée (voir
{\hyperref[\detokenize{2_BacLipidDB:colorisation-par-lot-d-import}]{\sphinxcrossref{Colorisation par lot d’import}}}).
\\
\sphinxbottomrule
\end{tabular}
\sphinxtableafterendhook\par
\sphinxattableend\end{savenotes}


\section{Dépendances internes}
\label{\detokenize{2_BacLipidDB:dependances-internes}}\begin{itemize}
\item {} 
\sphinxAtStartPar
\sphinxcode{\sphinxupquote{utils.data\_access.load\_database()}}

\item {} 
\sphinxAtStartPar
\sphinxcode{\sphinxupquote{utils.history.load\_history()}}

\item {} 
\sphinxAtStartPar
{\hyperref[\detokenize{api:module-config}]{\sphinxcrossref{\sphinxcode{\sphinxupquote{config}}}}} (\sphinxcode{\sphinxupquote{get\_engine}}, \sphinxcode{\sphinxupquote{PROJECT\_ROOT}} importés directement ; \sphinxcode{\sphinxupquote{HISTORY\_PATH}} est
utilisé indirectement, à l’intérieur de \sphinxcode{\sphinxupquote{load\_history}})

\end{itemize}

\sphinxstepscope


\chapter{Export}
\label{\detokenize{3_Export:export}}\label{\detokenize{3_Export::doc}}
\sphinxAtStartPar
Fichier source : \sphinxcode{\sphinxupquote{app/pages/3\_Export.py}}

\begin{sphinxadmonition}{note}{Note:}
\sphinxAtStartPar
Cette page contient la documentation du module Streamlit (\sphinxcode{\sphinxupquote{3\_Export.py}}).
La majorité du code n’est pas écrit sous forme de fonction, il n’est donc pas possible d’utiliser \sphinxcode{\sphinxupquote{automodule}}
pour générer cette page automatiquement. La documentation ci\sphinxhyphen{}dessous est rédigée manuellement
à partir du code source.
\end{sphinxadmonition}


\section{Vue d’ensemble}
\label{\detokenize{3_Export:vue-d-ensemble}}
\sphinxAtStartPar
Ce module Streamlit (\sphinxstylestrong{MODULE 3}) permet de filtrer les annotations de
\sphinxstylestrong{BacLipidDB} par catégorie, classe et sous\sphinxhyphen{}classe lipidique, niveau MS,
mode d’ionisation et plage de m/z. Si des RT et des CCS sont renseignés,
il est également possible de filtrer sur ces valeurs. Les annotations filtrées peuvent ensuite être visualisées
puis téléchargées dans des formats compatible avec MZmine .csv, .msp.


\section{Chargement des données}
\label{\detokenize{3_Export:chargement-des-donnees}}
\sphinxAtStartPar
\sphinxcode{\sphinxupquote{\_load\_data(\_engine)}} s’appuie sur \sphinxcode{\sphinxupquote{utils.data\_access.load\_database()}} (la vue jointe
commune \sphinxcode{\sphinxupquote{Annotation}} + \sphinxcode{\sphinxupquote{Lipid}} + \sphinxcode{\sphinxupquote{Detection}}, mise en cache via \sphinxcode{\sphinxupquote{st.cache\_data(ttl=60)}}
et partagée avec les pages 2\_BacLipidDB et 5\_Management), renomme certaines colonnes (\sphinxcode{\sphinxupquote{Lipid\_name}} \(\rightarrow\) \sphinxcode{\sphinxupquote{name}},
\sphinxcode{\sphinxupquote{Formula}} \(\rightarrow\) \sphinxcode{\sphinxupquote{formula}}, \sphinxcode{\sphinxupquote{Precursor\_MZ}} \(\rightarrow\) \sphinxcode{\sphinxupquote{mz}}, \sphinxcode{\sphinxupquote{Neutral\_mass}} \(\rightarrow\) \sphinxcode{\sphinxupquote{neutral\_mass}}). Les lignes conservées sont :
\sphinxcode{\sphinxupquote{name}}, \sphinxcode{\sphinxupquote{formula}}, \sphinxcode{\sphinxupquote{Lipid\_category}}, \sphinxcode{\sphinxupquote{Lipid\_class}}, \sphinxcode{\sphinxupquote{Lipid\_subclass}}, \sphinxcode{\sphinxupquote{mz}}, \sphinxcode{\sphinxupquote{neutral\_mass}}, \sphinxcode{\sphinxupquote{MS\_level}}, \sphinxcode{\sphinxupquote{Ionisation\_mode}}, \sphinxcode{\sphinxupquote{Adduct}}, \sphinxcode{\sphinxupquote{RT}}, \sphinxcode{\sphinxupquote{CCS}} et
\sphinxcode{\sphinxupquote{Num\_Peaks}}. Les colonnes \sphinxcode{\sphinxupquote{RT}}/\sphinxcode{\sphinxupquote{CCS}}/\sphinxcode{\sphinxupquote{Num\_Peaks}} valent \sphinxcode{\sphinxupquote{None}} si non renseignés en base, \sphinxcode{\sphinxupquote{Num\_Peaks}} n’est renseigné que pour les détections \sphinxcode{\sphinxupquote{MS2}}. Si le chargement échoue, un message \sphinxcode{\sphinxupquote{st.error}} est affiché
(« \sphinxstyleemphasis{Unable to load data from the database: \{e\}} ») et la page s’arrête (\sphinxcode{\sphinxupquote{st.stop()}}). Sinon si la base est vide, un message \sphinxcode{\sphinxupquote{st.info}} est affiché à la place (« \sphinxstyleemphasis{The database contains no data yet.} »).


\section{Filtres}
\label{\detokenize{3_Export:filtres}}
\sphinxAtStartPar
Les filtres sont répartis sur deux lignes de 3 colonnes (\sphinxcode{\sphinxupquote{st.columns}}) :
\begin{itemize}
\item {} 
\sphinxAtStartPar
\sphinxstylestrong{Lipid category}, \sphinxstylestrong{Lipid class}, \sphinxstylestrong{Lipid subclass} (\sphinxcode{\sphinxupquote{st.multiselect}}) : les valeurs
possibles sont extraites des colonnes présente dans la base de données chargée. Une option
\sphinxcode{\sphinxupquote{(None)}} est ajoutée si la colonne contient des valeurs manquantes.

\item {} 
\sphinxAtStartPar
\sphinxstylestrong{MS level} et \sphinxstylestrong{Ionisation mode} (\sphinxcode{\sphinxupquote{st.multiselect}}) : mêmes principes mais sans option
{\color{red}\bfseries{}\textasciigrave{}\textasciigrave{}}(None)\textasciigrave{}\textasciigrave{}car toutes les lignes ont une valeur pour ces colonnes.

\item {} 
\sphinxAtStartPar
\sphinxstylestrong{Precursor m/z range} : deux \sphinxcode{\sphinxupquote{st.number\_input}} (\sphinxcode{\sphinxupquote{Min m/z}} / \sphinxcode{\sphinxupquote{Max m/z}},
format \sphinxcode{\sphinxupquote{\%.6f}}), initialisés au min/max de la colonne \sphinxcode{\sphinxupquote{mz}}. Les deux valeurs sont triées (\sphinxcode{\sphinxupquote{min}}/\sphinxcode{\sphinxupquote{max}}) pour former \sphinxcode{\sphinxupquote{mz\_range}}.

\end{itemize}


\section{Application des filtres}
\label{\detokenize{3_Export:application-des-filtres}}
\sphinxAtStartPar
Le DataFrame \sphinxcode{\sphinxupquote{df\_all}} est copié dans \sphinxcode{\sphinxupquote{df\_filtered}}, puis chaque filtre est appliqué
l’un après l’autre (chaque étape réduit \sphinxcode{\sphinxupquote{df\_filtered}} obtenu à l’étape précédente) :
\begin{itemize}
\item {} 
\sphinxAtStartPar
\sphinxstylestrong{Lipid\_category}, \sphinxstylestrong{Lipid\_class}, \sphinxstylestrong{Lipid\_subclass} : si aucune valeur n’est
sélectionnée pour la colonne, aucun filtre n’est appliqué (toutes les lignes sont
conservées). Sinon, seules les lignes dont la valeur figure parmi les filtres sélectionnés
sont conservées. Si l’option \sphinxcode{\sphinxupquote{(None)}} fait partie de la sélection, les lignes où la
colonne est manquante (\sphinxcode{\sphinxupquote{NaN}}) sont conservées en plus.

\item {} 
\sphinxAtStartPar
\sphinxstylestrong{MS\_level}, \sphinxstylestrong{Ionisation\_mode} : même principe, mais sans option \sphinxcode{\sphinxupquote{(None)}} (ces
colonnes sont toujours renseignées). Aucun filtre n’est appliqué si la sélection est vide,
sinon seules les lignes dont la valeur figure parmi les filtres sélectionnés sont conservées.

\item {} 
\sphinxAtStartPar
\sphinxstylestrong{mz} : seules les lignes dont la valeur de \sphinxcode{\sphinxupquote{mz}} est comprise entre les deux bornes
de \sphinxcode{\sphinxupquote{mz\_range}} (bornes incluses) sont conservées. Ce filtre est toujours appliqué
(les bornes valent par défaut le min/max de la colonne, donc rien n’est exclu tant que
l’utilisateur ne les modifie pas).

\end{itemize}


\section{Aperçu}
\label{\detokenize{3_Export:apercu}}
\sphinxAtStartPar
\sphinxcode{\sphinxupquote{st.subheader}} affiche le nombre de lignes filtrées. Si aucune ligne ne correspond,
un message \sphinxcode{\sphinxupquote{st.warning}} est affiché et la page s’arrête. Sinon, un extrait des colonnes
\sphinxcode{\sphinxupquote{name}}, \sphinxcode{\sphinxupquote{formula}}, \sphinxcode{\sphinxupquote{mz}}, \sphinxcode{\sphinxupquote{MS\_level}}, \sphinxcode{\sphinxupquote{RT}}, \sphinxcode{\sphinxupquote{CCS}} et \sphinxcode{\sphinxupquote{Num\_Peaks}} est affiché
via \sphinxcode{\sphinxupquote{st.dataframe\textasciigrave{}. \textasciigrave{}\textasciigrave{}Num\_Peaks}} vaut \sphinxcode{\sphinxupquote{None}} pour les lignes \sphinxcode{\sphinxupquote{MS1}} (le nombre de pics n’a de sens qu’en MS2).


\section{Téléchargement}
\label{\detokenize{3_Export:telechargement}}
\sphinxAtStartPar
Trois boutons de téléchargement (\sphinxcode{\sphinxupquote{st.columns}} à 3 colonnes) :
\begin{itemize}
\item {} 
\sphinxAtStartPar
\sphinxstylestrong{CSV (MS1 + MS2)} : colonnes \sphinxcode{\sphinxupquote{neutral\_mass}}, \sphinxcode{\sphinxupquote{mz}}, \sphinxcode{\sphinxupquote{formula}}, \sphinxcode{\sphinxupquote{name}} de
l’ensemble de \sphinxcode{\sphinxupquote{df\_filtered}}, exportées en CSV. Toujours disponible dès que des données sont présentes dans BacLipidDB.

\item {} 
\sphinxAtStartPar
\sphinxstylestrong{MSP (MS2)} : disponible uniquement si \sphinxcode{\sphinxupquote{df\_filtered}} contient au moins une ligne
\sphinxcode{\sphinxupquote{MS2}} (bouton \sphinxcode{\sphinxupquote{st.button}} désactivé sinon). Généré par \sphinxcode{\sphinxupquote{generate\_msp()}} (mise en
cache via \sphinxcode{\sphinxupquote{st.cache\_data(ttl=60)}}), qui délègue à {\hyperref[\detokenize{api:utils.msp_export.generate_msp}]{\sphinxcrossref{\sphinxcode{\sphinxupquote{utils.msp\_export.generate\_msp()}}}}}
avec les catégories, classes et sous\sphinxhyphen{}classes sélectionnées, \sphinxcode{\sphinxupquote{mz\_range}} ainsi que les
modes d’ionisation sélectionnés (\sphinxcode{\sphinxupquote{selected\_ionisation\_modes}}). Le filtre \sphinxstylestrong{MS level}
choisi par l’utilisateur n’est en revanche pas transmis à cette fonction : elle filtre
elle\sphinxhyphen{}même en dur sur \sphinxcode{\sphinxupquote{Detection.MS\_level == "MS2"}}, si bien que ce filtre n’aurait de
toute façon aucun effet sur cet export. Cette fonction ne conserve que les annotations \sphinxcode{\sphinxupquote{MS2}} possédant au moins un fragment, et écrit
pour chacune un bloc \sphinxcode{\sphinxupquote{Name}}/\sphinxcode{\sphinxupquote{PrecursorMZ}}/\sphinxcode{\sphinxupquote{MW}}/\sphinxcode{\sphinxupquote{ExactMass}}/\sphinxcode{\sphinxupquote{Ion\_mode}}/\sphinxcode{\sphinxupquote{RT}}/\sphinxcode{\sphinxupquote{CCS}}/\sphinxcode{\sphinxupquote{Num Peaks}}
(\sphinxcode{\sphinxupquote{RT}}/\sphinxcode{\sphinxupquote{CCS}} omis si absents) suivi de la liste des fragments (\sphinxcode{\sphinxupquote{MZ Intensity}}).
\sphinxcode{\sphinxupquote{Ion\_mode}} reprend le mode d’ionisation de la détection, encodé comme dans les fichiers
MSP de LipidMaps (\sphinxcode{\sphinxupquote{P}} pour \sphinxcode{\sphinxupquote{Positive}}, \sphinxcode{\sphinxupquote{N}} pour \sphinxcode{\sphinxupquote{Negative}}).

\item {} 
\sphinxAtStartPar
\sphinxstylestrong{CSV (CCS + RT)} : disponible uniquement si \sphinxcode{\sphinxupquote{df\_filtered}} contient au moins une ligne
avec à la fois \sphinxcode{\sphinxupquote{RT}} et \sphinxcode{\sphinxupquote{CCS}} renseignés (bouton désactivé sinon). Colonnes
\sphinxcode{\sphinxupquote{neutral\_mass}}, \sphinxcode{\sphinxupquote{mz}}, \sphinxcode{\sphinxupquote{formula}}, \sphinxcode{\sphinxupquote{name}}, \sphinxcode{\sphinxupquote{RT}}, \sphinxcode{\sphinxupquote{CCS}} des lignes concernées,
qu’il s’agisse d’annotations \sphinxcode{\sphinxupquote{MS1}} ou \sphinxcode{\sphinxupquote{MS2}}. MZmine n’acceptant pas un CSV MS1 où
certaines lignes ont des valeurs RT/CCS et d’autres non, cet export est volontairement
séparé du CSV MS1 + MS2 ci\sphinxhyphen{}dessus.

\end{itemize}

\sphinxAtStartPar
Chaque fichier est nommé \sphinxcode{\sphinxupquote{annotation\_export\_\textless{}niveau\textgreater{}\_\textless{}date du jour AAAAMMJJ\textgreater{}.\textless{}ext\textgreater{}}}
(\sphinxcode{\sphinxupquote{MS1+MS2}}, \sphinxcode{\sphinxupquote{MS2}} ou \sphinxcode{\sphinxupquote{CCS\_RT}}).


\section{Fonctions internes}
\label{\detokenize{3_Export:fonctions-internes}}

\begin{savenotes}\sphinxattablestart
\sphinxthistablewithglobalstyle
\centering
\begin{tabular}[t]{\X{30}{100}\X{70}{100}}
\sphinxtoprule
\sphinxstyletheadfamily 
\sphinxAtStartPar
Fonction
&\sphinxstyletheadfamily 
\sphinxAtStartPar
Description
\\
\sphinxmidrule
\sphinxtableatstartofbodyhook
\sphinxAtStartPar
\sphinxcode{\sphinxupquote{\_load\_data(\_engine: Engine) \sphinxhyphen{}\textgreater{} pandas.DataFrame}}
&
\sphinxAtStartPar
Recadre et renomme les colonnes utiles de la vue jointe renvoyée par
\sphinxcode{\sphinxupquote{utils.data\_access.load\_database()}}.
\\
\sphinxhline
\sphinxAtStartPar
\sphinxcode{\sphinxupquote{generate\_msp(\_engine, categories, classes, sub\_classes, mz\_range, ionisation\_modes) \sphinxhyphen{}\textgreater{} str}}
&
\sphinxAtStartPar
Met en cache l’appel à {\hyperref[\detokenize{api:utils.msp_export.generate_msp}]{\sphinxcrossref{\sphinxcode{\sphinxupquote{utils.msp\_export.generate\_msp()}}}}} pour l’export MSP
des détections MS2 correspondant aux filtres.
\\
\sphinxbottomrule
\end{tabular}
\sphinxtableafterendhook\par
\sphinxattableend\end{savenotes}


\section{Dépendances internes}
\label{\detokenize{3_Export:dependances-internes}}\begin{itemize}
\item {} 
\sphinxAtStartPar
\sphinxcode{\sphinxupquote{utils.data\_access.load\_database()}}

\item {} 
\sphinxAtStartPar
{\hyperref[\detokenize{api:module-config}]{\sphinxcrossref{\sphinxcode{\sphinxupquote{config}}}}} (\sphinxcode{\sphinxupquote{get\_engine}})

\item {} 
\sphinxAtStartPar
{\hyperref[\detokenize{api:utils.msp_export.generate_msp}]{\sphinxcrossref{\sphinxcode{\sphinxupquote{utils.msp\_export.generate\_msp()}}}}}

\end{itemize}

\sphinxstepscope


\chapter{Resources}
\label{\detokenize{4_Resources:resources}}\label{\detokenize{4_Resources::doc}}
\sphinxAtStartPar
Fichier source : \sphinxcode{\sphinxupquote{app/pages/4\_Resources.py}}

\begin{sphinxadmonition}{note}{Note:}
\sphinxAtStartPar
Cette page contient la documentation du module Streamlit (\sphinxcode{\sphinxupquote{4\_Resources.py}}).
La majorité du code n’est pas écrit sous forme de fonction, il n’est donc pas possible d’utiliser \sphinxcode{\sphinxupquote{automodule}}
pour générer cette page automatiquement. La documentation ci\sphinxhyphen{}dessous est rédigée manuellement
à partir du code source.
\end{sphinxadmonition}


\section{Vue d’ensemble}
\label{\detokenize{4_Resources:vue-d-ensemble}}
\sphinxAtStartPar
Ce module Streamlit (\sphinxstylestrong{MODULE 4}) sert de guide pour préparer les fichiers
d’intégration destinés au {\hyperref[\detokenize{1_Integration::doc}]{\sphinxcrossref{\DUrole{doc}{Data\_Integration}}}} (\sphinxstylestrong{MODULE 1}). Il présente,
pour les niveaux \sphinxstylestrong{MS1} et \sphinxstylestrong{MS2}, la liste des colonnes attendues avec
leur type et leur caractère obligatoire/optionnel, propose un fichier modèle
téléchargeable pour chaque niveau. Il liste également les adduits acceptés dans la
colonne \sphinxcode{\sphinxupquote{Adduct}} et permet d’en ajouter ou d’en supprimer.


\section{Guide des colonnes MS1}
\label{\detokenize{4_Resources:guide-des-colonnes-ms1}}
\sphinxAtStartPar
\sphinxcode{\sphinxupquote{st.dataframe}} affiche un tableau statique (\sphinxcode{\sphinxupquote{guide\_ms1}}) décrivant les
colonnes attendues dans un fichier d’intégration \sphinxstylestrong{MS1} : \sphinxcode{\sphinxupquote{Precursor\_MZ}},
\sphinxcode{\sphinxupquote{Adduct}}, \sphinxcode{\sphinxupquote{Formula}}, \sphinxcode{\sphinxupquote{FA\_composition}}, \sphinxcode{\sphinxupquote{Lipid\_Name}},
\sphinxcode{\sphinxupquote{Lipid\_category}}, \sphinxcode{\sphinxupquote{Lipid\_class}}, \sphinxcode{\sphinxupquote{Lipid\_subclass}}, \sphinxcode{\sphinxupquote{RT}} et \sphinxcode{\sphinxupquote{CCS}},
avec pour chacune son type et si elle est obligatoire (✅) ou optionnelle (❌).
Seules \sphinxcode{\sphinxupquote{FA\_composition}}, \sphinxcode{\sphinxupquote{RT}} et \sphinxcode{\sphinxupquote{CCS}} sont optionnelles (à laisser
vides si non disponibles).


\section{Modèle MS1}
\label{\detokenize{4_Resources:modele-ms1}}
\sphinxAtStartPar
Un \sphinxcode{\sphinxupquote{st.download\_button}} génère et propose au téléchargement
\sphinxcode{\sphinxupquote{BacLipidDB\_MS1\_template.csv}} : un DataFrame d’une seule ligne (\sphinxcode{\sphinxupquote{template\_ms1}})
illustrant les 7 colonnes obligatoires et les 2 colonnes optionnelles
(\sphinxcode{\sphinxupquote{RT}}, \sphinxcode{\sphinxupquote{CCS}}) laissées vides, converti en CSV en mémoire via
\sphinxcode{\sphinxupquote{io.StringIO}}.


\section{Guide des colonnes MS2}
\label{\detokenize{4_Resources:guide-des-colonnes-ms2}}
\sphinxAtStartPar
Un fichier d’intégration \sphinxstylestrong{MS2} est organisé en blocs empilés verticalement,
un bloc par scan, séparés par une ligne vide. Chaque bloc commence par une
ligne \sphinxcode{\sphinxupquote{Scan \#\textless{}id\textgreater{}}}, suivie d’une ligne \sphinxcode{\sphinxupquote{Label,Value}} par champ (ordre
libre), puis de l’en\sphinxhyphen{}tête \sphinxcode{\sphinxupquote{m/z,Intensity}} et des lignes de fragments
(\sphinxcode{\sphinxupquote{st.code}} affiche ce format en exemple).

\sphinxAtStartPar
\sphinxcode{\sphinxupquote{st.dataframe}} affiche ensuite un tableau statique (\sphinxcode{\sphinxupquote{guide\_ms2}}) décrivant
les colonnes/champs attendus : \sphinxcode{\sphinxupquote{Scan \#\textless{}id\textgreater{}}}, \sphinxcode{\sphinxupquote{Precursor\_MZ}}, \sphinxcode{\sphinxupquote{Adduct}},
\sphinxcode{\sphinxupquote{Formula}}, \sphinxcode{\sphinxupquote{FA\_composition}}, \sphinxcode{\sphinxupquote{Lipid\_Name}}, \sphinxcode{\sphinxupquote{Lipid\_category}},
\sphinxcode{\sphinxupquote{Lipid\_class}}, \sphinxcode{\sphinxupquote{Lipid\_subclass}}, \sphinxcode{\sphinxupquote{Num\_peaks}}, \sphinxcode{\sphinxupquote{m/z}}, \sphinxcode{\sphinxupquote{Intensity}},
\sphinxcode{\sphinxupquote{RT}} et \sphinxcode{\sphinxupquote{CCS}}. Seules \sphinxcode{\sphinxupquote{FA\_composition}}, \sphinxcode{\sphinxupquote{RT}} et \sphinxcode{\sphinxupquote{CCS}} sont
optionnelles (à supprimer entièrement du bloc si non disponibles, contrairement
au MS1 où elles sont laissées vides). \sphinxcode{\sphinxupquote{Num\_peaks}} doit correspondre au
nombre de lignes \sphinxcode{\sphinxupquote{m/z}}/\sphinxcode{\sphinxupquote{Intensity}} du bloc.


\section{Modèle MS2}
\label{\detokenize{4_Resources:modele-ms2}}
\sphinxAtStartPar
Un \sphinxcode{\sphinxupquote{st.download\_button}} génère et propose au téléchargement
\sphinxcode{\sphinxupquote{BacLipidDB\_MS2\_template.csv}} : deux scans d’exemple (\sphinxcode{\sphinxupquote{template\_ms2}},
liste de lignes écrites via \sphinxcode{\sphinxupquote{csv.writer}} dans un \sphinxcode{\sphinxupquote{io.StringIO}}) séparés
par une ligne vide. Le premier scan inclut les champs optionnels \sphinxcode{\sphinxupquote{RT}} et
\sphinxcode{\sphinxupquote{CCS}} ; le second montre que ces lignes peuvent être omises entièrement
quand ces valeurs ne sont pas disponibles.


\section{Adduits}
\label{\detokenize{4_Resources:adduits}}

\subsection{Table des adduits}
\label{\detokenize{4_Resources:table-des-adduits}}
\sphinxAtStartPar
\sphinxcode{\sphinxupquote{st.dataframe}} affiche la table des adduits acceptés dans la colonne
\sphinxcode{\sphinxupquote{Adduct}} (intégration MS1 et MS2), construite à partir de
\sphinxcode{\sphinxupquote{utils.neutral\_mass.ADDUCT\_SHIFTS}} et \sphinxcode{\sphinxupquote{utils.neutral\_mass.ADDUCT\_SIGNS}} :
nom de l’adduit, décalage de masse en Da (\sphinxcode{\sphinxupquote{Mass shift}}), opération à
appliquer au m/z pour obtenir la masse neutre (\sphinxcode{\sphinxupquote{m/z \sphinxhyphen{} shift}} si
\sphinxcode{\sphinxupquote{ADDUCT\_SIGNS{[}name{]} == \sphinxhyphen{}1}}, \sphinxcode{\sphinxupquote{m/z + shift}} sinon) et indicateur \sphinxcode{\sphinxupquote{Custom}}
(✅ si l’adduit provient de \sphinxcode{\sphinxupquote{utils.neutral\_mass.list\_custom\_adducts()}},
❌ s’il s’agit d’un adduit natif).


\subsection{Ajout d’un adduit personnalisé}
\label{\detokenize{4_Resources:ajout-d-un-adduit-personnalise}}
\sphinxAtStartPar
Un formulaire (\sphinxcode{\sphinxupquote{st.form}}) permet de saisir un nom d’adduit
(\sphinxcode{\sphinxupquote{st.text\_input}}, ex. \sphinxcode{\sphinxupquote{{[}M+Cl{]}\sphinxhyphen{}}}), un décalage de masse en Da
(\sphinxcode{\sphinxupquote{st.number\_input}}) et le type de formation (\sphinxcode{\sphinxupquote{st.radio}}) : « Addition »
(ex. \sphinxcode{\sphinxupquote{{[}M+H{]}+}}, \sphinxcode{\sphinxupquote{{[}M+Na{]}+}}) ou « Loss » (ex. \sphinxcode{\sphinxupquote{{[}M\sphinxhyphen{}H{]}\sphinxhyphen{}}}). Le signe transmis
à \sphinxcode{\sphinxupquote{utils.neutral\_mass.add\_custom\_adduct()}} vaut \sphinxcode{\sphinxupquote{\sphinxhyphen{}1}} pour une addition
et \sphinxcode{\sphinxupquote{1}} pour une perte, conformément à la convention de
{\hyperref[\detokenize{api:utils.neutral_mass.neutral_mass}]{\sphinxcrossref{\sphinxcode{\sphinxupquote{utils.neutral\_mass.neutral\_mass()}}}}} (\sphinxcode{\sphinxupquote{mz + sign * shift}}). En cas de
succès, un message \sphinxcode{\sphinxupquote{st.success}} est affiché et la page est rechargée
(\sphinxcode{\sphinxupquote{st.rerun()}}) ; en cas d’erreur (nom vide/déjà utilisé, décalage non
strictement positif…), un message \sphinxcode{\sphinxupquote{st.error}} affiche l’exception levée
par \sphinxcode{\sphinxupquote{utils.neutral\_mass.add\_custom\_adduct()}}.

\sphinxAtStartPar
Les adduits personnalisés sont persistés dans \sphinxcode{\sphinxupquote{adducts.json}} (\sphinxcode{\sphinxupquote{config.ADDUCTS\_PATH}})
et restent donc disponibles d’une session à l’autre.


\subsection{Suppression d’un adduit personnalisé}
\label{\detokenize{4_Resources:suppression-d-un-adduit-personnalise}}
\sphinxAtStartPar
Si au moins un adduit personnalisé existe, un \sphinxcode{\sphinxupquote{st.selectbox}} (liste des
adduits renvoyés par \sphinxcode{\sphinxupquote{utils.neutral\_mass.list\_custom\_adducts()}}) et un
bouton \sphinxstylestrong{Remove} permettent de le retirer via
\sphinxcode{\sphinxupquote{utils.neutral\_mass.remove\_custom\_adduct()}}, puis de recharger la page
(\sphinxcode{\sphinxupquote{st.rerun()}}). Les adduits natifs (\sphinxcode{\sphinxupquote{BUILTIN\_ADDUCT\_SHIFTS}}) ne peuvent
pas être supprimés.


\section{Fonctions internes}
\label{\detokenize{4_Resources:fonctions-internes}}
\sphinxAtStartPar
Ce module ne définit aucune fonction interne : l’ensemble du code s’exécute
au niveau du script (délégation directe aux fonctions de
{\hyperref[\detokenize{api:module-utils.neutral_mass}]{\sphinxcrossref{\sphinxcode{\sphinxupquote{utils.neutral\_mass}}}}}).


\section{Dépendances internes}
\label{\detokenize{4_Resources:dependances-internes}}\begin{itemize}
\item {} 
\sphinxAtStartPar
{\hyperref[\detokenize{api:module-utils.neutral_mass}]{\sphinxcrossref{\sphinxcode{\sphinxupquote{utils.neutral\_mass}}}}} (\sphinxcode{\sphinxupquote{ADDUCT\_SHIFTS}}, \sphinxcode{\sphinxupquote{ADDUCT\_SIGNS}},
\sphinxcode{\sphinxupquote{add\_custom\_adduct}}, \sphinxcode{\sphinxupquote{list\_custom\_adducts}}, \sphinxcode{\sphinxupquote{remove\_custom\_adduct}})

\end{itemize}

\sphinxstepscope


\chapter{Management}
\label{\detokenize{5_Management:management}}\label{\detokenize{5_Management::doc}}
\sphinxAtStartPar
Fichier source : \sphinxcode{\sphinxupquote{app/pages/5\_Management.py}}

\begin{sphinxadmonition}{note}{Note:}
\sphinxAtStartPar
Cette page contient la documentation du module Streamlit (\sphinxcode{\sphinxupquote{5\_Management.py}}).
Le code c’est pas écrit sous forme de fonction, il n’est donc pas possible d’utiliser \sphinxcode{\sphinxupquote{automodule}}
pour générer cette page automatiquement. La documentation ci\sphinxhyphen{}dessous est rédigée manuellement
à partir du code source.
\end{sphinxadmonition}


\section{Vue d’ensemble}
\label{\detokenize{5_Management:vue-d-ensemble}}
\sphinxAtStartPar
Ce module Streamlit (\sphinxstylestrong{MODULE 5}) permet de corriger ou de supprimer des données déjà
intégrées dans \sphinxstylestrong{BacLipidDB}, via deux sections indépendantes :
\begin{itemize}
\item {} 
\sphinxAtStartPar
\sphinxstylestrong{Edit \& delete records} : édition cellule par cellule des annotations existantes, ou
suppression de lignes individuelles.

\item {} 
\sphinxAtStartPar
\sphinxstylestrong{Delete an entire import} : suppression en une fois de tout ce qu’un import donné a inséré
en base, à partir de l’historique.

\end{itemize}

\sphinxAtStartPar
Une sauvegarde de la base (\sphinxcode{\sphinxupquote{utils.db\_backup.backup\_database()}}) est systématiquement prise
avant toute modification, et chaque action destructrice ou modificatrice passe par une boîte de
dialogue de confirmation (\sphinxcode{\sphinxupquote{@st.dialog}}).


\section{Edit \& delete records}
\label{\detokenize{5_Management:edit-delete-records}}

\subsection{Chargement des données}
\label{\detokenize{5_Management:chargement-des-donnees}}
\sphinxAtStartPar
\sphinxcode{\sphinxupquote{\_load\_database()}} s’appuie sur \sphinxcode{\sphinxupquote{utils.data\_access.load\_annotations\_df()}} (la même
requête jointe \sphinxcode{\sphinxupquote{Annotation}} + \sphinxcode{\sphinxupquote{Lipid}} + \sphinxcode{\sphinxupquote{Detection}} que les pages \sphinxstylestrong{2\_BacLipidDB} et
\sphinxstylestrong{3\_Export}), dont elle ne conserve que les colonnes nécessaires à l’affichage et au report des
modifications en base : \sphinxcode{\sphinxupquote{Annotation\_ID}}, \sphinxcode{\sphinxupquote{Detection\_ID}}, \sphinxcode{\sphinxupquote{Lipid\_ID}}, \sphinxcode{\sphinxupquote{Lipid\_name}},
\sphinxcode{\sphinxupquote{Formula}}, \sphinxcode{\sphinxupquote{Lipid\_category}}, \sphinxcode{\sphinxupquote{Lipid\_class}}, \sphinxcode{\sphinxupquote{Lipid\_subclass}}, \sphinxcode{\sphinxupquote{Precursor\_MZ}},
\sphinxcode{\sphinxupquote{Ionisation\_mode}}, \sphinxcode{\sphinxupquote{Adduct}}, \sphinxcode{\sphinxupquote{Neutral\_mass}}, \sphinxcode{\sphinxupquote{Molecular\_weight}}, \sphinxcode{\sphinxupquote{Monoisotopic\_mass}},
\sphinxcode{\sphinxupquote{MS\_level}}, \sphinxcode{\sphinxupquote{Num\_Peaks}}, \sphinxcode{\sphinxupquote{RT}}, \sphinxcode{\sphinxupquote{CCS}}. Si le chargement échoue, un message \sphinxcode{\sphinxupquote{st.error}}
est affiché et la page s’arrête (\sphinxcode{\sphinxupquote{st.stop()}}) ; si la base est vide, un message \sphinxcode{\sphinxupquote{st.info}} est
affiché à la place.


\subsection{Édition du tableau}
\label{\detokenize{5_Management:edition-du-tableau}}
\sphinxAtStartPar
Une colonne \sphinxcode{\sphinxupquote{Delete}} (case à cocher) est ajoutée en tête du tableau affiché via
\sphinxcode{\sphinxupquote{st.data\_editor}}. Les colonnes \sphinxcode{\sphinxupquote{Annotation\_ID}}, \sphinxcode{\sphinxupquote{Detection\_ID}}, \sphinxcode{\sphinxupquote{Lipid\_ID}},
\sphinxcode{\sphinxupquote{Precursor\_MZ}}, \sphinxcode{\sphinxupquote{Ionisation\_mode}}, \sphinxcode{\sphinxupquote{Adduct}}, \sphinxcode{\sphinxupquote{Neutral\_mass}}, \sphinxcode{\sphinxupquote{Molecular\_weight}},
\sphinxcode{\sphinxupquote{Monoisotopic\_mass}} et \sphinxcode{\sphinxupquote{MS\_level}} sont verrouillées (\sphinxcode{\sphinxupquote{disabled}}). Restent donc modifiables :
\begin{itemize}
\item {} 
\sphinxAtStartPar
\sphinxcode{\sphinxupquote{Lipid\_name}}, \sphinxcode{\sphinxupquote{Lipid\_category}}, \sphinxcode{\sphinxupquote{Lipid\_class}}, \sphinxcode{\sphinxupquote{Lipid\_subclass}} (\sphinxcode{\sphinxupquote{EDITABLE\_LIPID\_FIELDS}}) ;

\item {} 
\sphinxAtStartPar
\sphinxcode{\sphinxupquote{Formula}} ;

\item {} 
\sphinxAtStartPar
\sphinxcode{\sphinxupquote{Num\_Peaks}}, \sphinxcode{\sphinxupquote{RT}}, \sphinxcode{\sphinxupquote{CCS}} (\sphinxcode{\sphinxupquote{EDITABLE\_DETECTION\_FIELDS}}).

\end{itemize}

\sphinxAtStartPar
\sphinxcode{\sphinxupquote{Precursor\_MZ}}, \sphinxcode{\sphinxupquote{Ionisation\_mode}} et \sphinxcode{\sphinxupquote{Adduct}} ne sont pas éditables ici car les modifier
nécessiterait de recalculer \sphinxcode{\sphinxupquote{Neutral\_mass}} (dérivée uniquement de ces deux dernières colonnes),
recalcul qui n’est pas implémenté sur cette page.


\subsection{Calcul du plan de modification}
\label{\detokenize{5_Management:calcul-du-plan-de-modification}}
\sphinxAtStartPar
\sphinxcode{\sphinxupquote{\_database\_modif(original, edited)}} compare le tableau édité à l’original ligne par ligne et
construit un plan \sphinxcode{\sphinxupquote{\{updates, deletes, errors\}}} :
\begin{itemize}
\item {} 
\sphinxAtStartPar
si la case \sphinxcode{\sphinxupquote{Delete}} est cochée, la ligne (avec ses \sphinxcode{\sphinxupquote{Annotation\_ID}}, \sphinxcode{\sphinxupquote{Detection\_ID}},
\sphinxcode{\sphinxupquote{Lipid\_ID}} et son nom) part directement dans \sphinxcode{\sphinxupquote{deletes}}, sans autre vérification ;

\item {} 
\sphinxAtStartPar
sinon, les colonnes de \sphinxcode{\sphinxupquote{NON\_EMPTY\_COLUMNS}} (\sphinxcode{\sphinxupquote{Lipid\_name}}, \sphinxcode{\sphinxupquote{Formula}}, \sphinxcode{\sphinxupquote{Precursor\_MZ}})
sont contrôlées : toute cellule vide ajoute un message dans \sphinxcode{\sphinxupquote{errors}} ;

\item {} 
\sphinxAtStartPar
les champs modifiés parmi \sphinxcode{\sphinxupquote{EDITABLE\_LIPID\_FIELDS}} et \sphinxcode{\sphinxupquote{EDITABLE\_DETECTION\_FIELDS}} sont
collectés séparément (\sphinxcode{\sphinxupquote{lipid\_fields}} / \sphinxcode{\sphinxupquote{detection\_fields}}), car ils seront reportés sur
deux tables différentes (\sphinxcode{\sphinxupquote{Lipid}} / \sphinxcode{\sphinxupquote{Detection}}) ;

\item {} 
\sphinxAtStartPar
si \sphinxcode{\sphinxupquote{Formula}} a changé, \sphinxcode{\sphinxupquote{Molecular\_weight}} et \sphinxcode{\sphinxupquote{Monoisotopic\_mass}} sont recalculées
({\hyperref[\detokenize{api:utils.molecular_weight.molecular_weight}]{\sphinxcrossref{\sphinxcode{\sphinxupquote{utils.molecular\_weight.molecular\_weight()}}}}}, {\hyperref[\detokenize{api:utils.monoisotopic.monoisotopic_mass}]{\sphinxcrossref{\sphinxcode{\sphinxupquote{utils.monoisotopic.monoisotopic\_mass()}}}}}),
exactement comme lors de la complétion automatique du module d’intégration (\sphinxstylestrong{1\_Integration}).
Si le calcul échoue, un message est ajouté à \sphinxcode{\sphinxupquote{errors}} et les masses ne sont pas mises à jour ;

\item {} 
\sphinxAtStartPar
une ligne sans aucun champ modifié n’est ajoutée ni à \sphinxcode{\sphinxupquote{updates}} ni à \sphinxcode{\sphinxupquote{deletes}}.

\end{itemize}

\sphinxAtStartPar
Le bouton \sphinxstylestrong{Apply changes} affiche les erreurs s’il y en a, un message \sphinxcode{\sphinxupquote{st.info}} si le plan est
vide, ou stocke le plan dans \sphinxcode{\sphinxupquote{st.session\_state{[}"pending\_plan"{]}}} et déclenche un \sphinxcode{\sphinxupquote{st.rerun()}}
pour ouvrir la confirmation.


\subsection{Confirmation et application}
\label{\detokenize{5_Management:confirmation-et-application}}
\sphinxAtStartPar
\sphinxcode{\sphinxupquote{\_confirm\_apply(plan)}} (\sphinxcode{\sphinxupquote{@st.dialog}}) résume le nombre de lignes à mettre à jour et à
supprimer, avertit que les suppressions sont irréversibles depuis l’application (\sphinxcode{\sphinxupquote{Fragment}},
\sphinxcode{\sphinxupquote{Annotation}}, \sphinxcode{\sphinxupquote{Detection}} et \sphinxcode{\sphinxupquote{Lipid}} liés) et liste les lignes concernées. Les boutons
\sphinxstylestrong{Confirm} / \sphinxstylestrong{Cancel} appellent respectivement \sphinxcode{\sphinxupquote{\_apply\_plan(plan)}} ou annulent, dans les deux
cas suivis d’un \sphinxcode{\sphinxupquote{st.rerun()}}.

\sphinxAtStartPar
\sphinxcode{\sphinxupquote{\_apply\_plan(plan)}} prend d’abord une sauvegarde (\sphinxcode{\sphinxupquote{backup\_database(label="manual\_edit")}}),
puis dans une session SQLAlchemy :
\begin{itemize}
\item {} 
\sphinxAtStartPar
pour chaque suppression, retire dans l’ordre les lignes liées de \sphinxcode{\sphinxupquote{Fragment}}, \sphinxcode{\sphinxupquote{Annotation}},
\sphinxcode{\sphinxupquote{Detection}} puis \sphinxcode{\sphinxupquote{Lipid}} (la base ne propageant pas ces suppressions automatiquement, elles
sont faites manuellement, dans cet ordre pour respecter les dépendances de clés étrangères) ;

\item {} 
\sphinxAtStartPar
pour chaque mise à jour, applique \sphinxcode{\sphinxupquote{lipid\_fields}} sur \sphinxcode{\sphinxupquote{Lipid}} et \sphinxcode{\sphinxupquote{detection\_fields}} sur
\sphinxcode{\sphinxupquote{Detection}} (une requête \sphinxcode{\sphinxupquote{update}} par table concernée, uniquement si le dictionnaire n’est
pas vide) ;

\item {} 
\sphinxAtStartPar
commit puis invalidation du cache (\sphinxcode{\sphinxupquote{load\_annotations\_df.clear()}}), pour que les pages
utilisant cette vue reflètent l’état à jour de la base.

\end{itemize}


\section{Delete an entire import}
\label{\detokenize{5_Management:delete-an-entire-import}}
\sphinxAtStartPar
L’historique des imports (\sphinxcode{\sphinxupquote{utils.history.load\_history()}}) est présenté dans un
\sphinxcode{\sphinxupquote{st.selectbox}} (\sphinxcode{\sphinxupquote{\_label(i)}} formate chaque entrée : nom de fichier, date d’insertion, nombre
de lignes et intégrateur). Une fois un import sélectionné, le nombre d’enregistrements de ce lot
encore présents en base est recalculé (\sphinxcode{\sphinxupquote{session.query(Detection)...count()}}) et affiché à côté
du nombre de lignes recensé au moment de l’import : les deux peuvent différer si des lignes de ce
lot ont déjà été supprimées individuellement via \sphinxstylestrong{Edit \& delete records}.

\sphinxAtStartPar
Le bouton \sphinxstylestrong{Delete this import} stocke l’index sélectionné dans
\sphinxcode{\sphinxupquote{st.session\_state{[}"pending\_history\_delete"{]}}} et déclenche un \sphinxcode{\sphinxupquote{st.rerun()}} pour ouvrir la
confirmation.

\sphinxAtStartPar
\sphinxcode{\sphinxupquote{\_confirm\_history\_delete(entry, index)}} (\sphinxcode{\sphinxupquote{@st.dialog}}) avertit du nombre d’enregistrements qui
seront supprimés et rappelle qu’une sauvegarde est prise automatiquement. Les boutons
\sphinxstylestrong{Confirm deletion} / \sphinxstylestrong{Cancel} appellent respectivement \sphinxcode{\sphinxupquote{\_delete\_import(entry, index)}} ou
annulent.

\sphinxAtStartPar
\sphinxcode{\sphinxupquote{\_delete\_import(entry, index)}} prend une sauvegarde
(\sphinxcode{\sphinxupquote{backup\_database(label=f"delete\_\{entry.get(\textquotesingle{}filename\textquotesingle{}, \textquotesingle{}import\textquotesingle{})\}")}}), puis dans une session
SQLAlchemy : récupère les \sphinxcode{\sphinxupquote{Lipid\_ID}} liés aux \sphinxcode{\sphinxupquote{detection\_ids}} de l’entrée (via \sphinxcode{\sphinxupquote{Annotation}}),
supprime les lignes de \sphinxcode{\sphinxupquote{Fragment}}, \sphinxcode{\sphinxupquote{Annotation}} et \sphinxcode{\sphinxupquote{Detection}} correspondant à ces
\sphinxcode{\sphinxupquote{detection\_ids}}, puis celles de \sphinxcode{\sphinxupquote{Lipid}} correspondant aux \sphinxcode{\sphinxupquote{Lipid\_ID}} récupérés, commit,
invalide le cache (\sphinxcode{\sphinxupquote{load\_annotations\_df.clear()}}), puis retire l’entrée de l’historique
(\sphinxcode{\sphinxupquote{utils.history.remove\_history()}}).


\section{Fonctions internes}
\label{\detokenize{5_Management:fonctions-internes}}

\begin{savenotes}\sphinxattablestart
\sphinxthistablewithglobalstyle
\centering
\begin{tabular}[t]{\X{30}{100}\X{70}{100}}
\sphinxtoprule
\sphinxstyletheadfamily 
\sphinxAtStartPar
Fonction
&\sphinxstyletheadfamily 
\sphinxAtStartPar
Description
\\
\sphinxmidrule
\sphinxtableatstartofbodyhook
\sphinxAtStartPar
\sphinxcode{\sphinxupquote{\_load\_database() \sphinxhyphen{}\textgreater{} pandas.DataFrame}}
&
\sphinxAtStartPar
Charge la vue jointe Detection/Lipid/Annotation utilisée pour l’édition, restreinte aux
colonnes nécessaires à l’affichage et au report des modifications.
\\
\sphinxhline
\sphinxAtStartPar
\sphinxcode{\sphinxupquote{\_database\_modif(original, edited) \sphinxhyphen{}\textgreater{} dict}}
&
\sphinxAtStartPar
Compare le tableau édité à l’original et construit le plan de modification
\sphinxcode{\sphinxupquote{\{updates, deletes, errors\}}}.
\\
\sphinxhline
\sphinxAtStartPar
\sphinxcode{\sphinxupquote{\_apply\_plan(plan: dict) \sphinxhyphen{}\textgreater{} None}}
&
\sphinxAtStartPar
Sauvegarde la base puis applique le plan (suppressions en cascade et mises à jour
ciblées) dans une session SQLAlchemy.
\\
\sphinxhline
\sphinxAtStartPar
\sphinxcode{\sphinxupquote{\_confirm\_apply(plan: dict) \sphinxhyphen{}\textgreater{} None}}
&
\sphinxAtStartPar
Boîte de dialogue de confirmation pour le plan d’édition/suppression de lignes.
\\
\sphinxhline
\sphinxAtStartPar
\sphinxcode{\sphinxupquote{\_delete\_import(entry: dict, index: int) \sphinxhyphen{}\textgreater{} None}}
&
\sphinxAtStartPar
Supprime tous les enregistrements d’un lot d’import (Fragment/Annotation/Detection/Lipid)
puis retire l’entrée de l’historique.
\\
\sphinxhline
\sphinxAtStartPar
\sphinxcode{\sphinxupquote{\_confirm\_history\_delete(entry: dict, index: int) \sphinxhyphen{}\textgreater{} None}}
&
\sphinxAtStartPar
Boîte de dialogue de confirmation pour la suppression d’un import entier.
\\
\sphinxbottomrule
\end{tabular}
\sphinxtableafterendhook\par
\sphinxattableend\end{savenotes}


\section{Dépendances internes}
\label{\detokenize{5_Management:dependances-internes}}\begin{itemize}
\item {} 
\sphinxAtStartPar
\sphinxcode{\sphinxupquote{utils.data\_access.load\_annotations\_df()}}

\item {} 
\sphinxAtStartPar
\sphinxcode{\sphinxupquote{utils.db\_backup.backup\_database()}}

\item {} 
\sphinxAtStartPar
\sphinxcode{\sphinxupquote{utils.history.load\_history()}}

\item {} 
\sphinxAtStartPar
\sphinxcode{\sphinxupquote{utils.history.remove\_history()}}

\item {} 
\sphinxAtStartPar
{\hyperref[\detokenize{api:utils.molecular_weight.molecular_weight}]{\sphinxcrossref{\sphinxcode{\sphinxupquote{utils.molecular\_weight.molecular\_weight()}}}}}

\item {} 
\sphinxAtStartPar
{\hyperref[\detokenize{api:utils.monoisotopic.monoisotopic_mass}]{\sphinxcrossref{\sphinxcode{\sphinxupquote{utils.monoisotopic.monoisotopic\_mass()}}}}}

\item {} 
\sphinxAtStartPar
{\hyperref[\detokenize{api:module-config}]{\sphinxcrossref{\sphinxcode{\sphinxupquote{config}}}}} (\sphinxcode{\sphinxupquote{get\_engine}})

\end{itemize}

\sphinxstepscope


\chapter{Modèle de données}
\label{\detokenize{model:modele-de-donnees}}\label{\detokenize{model::doc}}
\sphinxAtStartPar
Fichier source : \sphinxcode{\sphinxupquote{models/model.py}}

\begin{sphinxadmonition}{note}{Note:}
\sphinxAtStartPar
Ce module définit les modèles SQLAlchemy (ORM) de \sphinxstylestrong{BacLipidDB}. Contrairement aux
pages Streamlit, il est structuré en classes documentées : voir {\hyperref[\detokenize{api:models.model.Detection}]{\sphinxcrossref{\sphinxcode{\sphinxupquote{models.model.Detection}}}}},
{\hyperref[\detokenize{api:models.model.Fragment}]{\sphinxcrossref{\sphinxcode{\sphinxupquote{models.model.Fragment}}}}}, {\hyperref[\detokenize{api:models.model.Lipid}]{\sphinxcrossref{\sphinxcode{\sphinxupquote{models.model.Lipid}}}}} et {\hyperref[\detokenize{api:models.model.Annotation}]{\sphinxcrossref{\sphinxcode{\sphinxupquote{models.model.Annotation}}}}}
pour la référence générée automatiquement (\sphinxcode{\sphinxupquote{autoclass}}). Cette page complète cette
référence par une vue d’ensemble du schéma et des relations entre les tables.
\end{sphinxadmonition}


\section{Vue d’ensemble}
\label{\detokenize{model:vue-d-ensemble}}
\sphinxAtStartPar
Le module définit les 4 tables de la base de données via \sphinxcode{\sphinxupquote{DeclarativeBase}} :
\begin{itemize}
\item {} 
\sphinxAtStartPar
\sphinxcode{\sphinxupquote{Detection}}  : signal détecté par le spectromètre de masse.

\item {} 
\sphinxAtStartPar
\sphinxcode{\sphinxupquote{Fragment}}   : ions fragments associés à une détection (MS2 uniquement).

\item {} 
\sphinxAtStartPar
\sphinxcode{\sphinxupquote{Lipid}}      : lipide de référence (nom, classe, catégorie, formule).

\item {} 
\sphinxAtStartPar
\sphinxcode{\sphinxupquote{Annotation}} : association entre une détection et un lipide candidat, avec un
niveau de confiance sur l’identification.

\end{itemize}


\section{Relations}
\label{\detokenize{model:relations}}

\begin{savenotes}\sphinxattablestart
\sphinxthistablewithglobalstyle
\centering
\begin{tabular}[t]{\X{40}{100}\X{15}{100}\X{45}{100}}
\sphinxtoprule
\sphinxstyletheadfamily 
\sphinxAtStartPar
Relation
&\sphinxstyletheadfamily 
\sphinxAtStartPar
Cardinalité
&\sphinxstyletheadfamily 
\sphinxAtStartPar
Description
\\
\sphinxmidrule
\sphinxtableatstartofbodyhook
\sphinxAtStartPar
\sphinxcode{\sphinxupquote{Detection}} \(\rightarrow\) \sphinxcode{\sphinxupquote{Fragment}}
&
\sphinxAtStartPar
1 \(\rightarrow\) N
&
\sphinxAtStartPar
Un signal peut produire plusieurs fragments (MS2 uniquement).
\\
\sphinxhline
\sphinxAtStartPar
\sphinxcode{\sphinxupquote{Detection}} \(\rightarrow\) \sphinxcode{\sphinxupquote{Annotation}}
&
\sphinxAtStartPar
1 \(\rightarrow\) 1
&
\sphinxAtStartPar
Un signal correspond à une seule annotation.
\\
\sphinxhline
\sphinxAtStartPar
\sphinxcode{\sphinxupquote{Lipid}} \(\rightarrow\) \sphinxcode{\sphinxupquote{Annotation}}
&
\sphinxAtStartPar
1 \(\rightarrow\) 1
&
\sphinxAtStartPar
Un lipide est associé à une seule annotation.
\\
\sphinxbottomrule
\end{tabular}
\sphinxtableafterendhook\par
\sphinxattableend\end{savenotes}


\section{Table \sphinxstyleliteralintitle{\sphinxupquote{Detection}}}
\label{\detokenize{model:table-detection}}

\begin{savenotes}\sphinxattablestart
\sphinxthistablewithglobalstyle
\centering
\begin{tabular}[t]{\X{25}{100}\X{20}{100}\X{55}{100}}
\sphinxtoprule
\sphinxstyletheadfamily 
\sphinxAtStartPar
Colonne
&\sphinxstyletheadfamily 
\sphinxAtStartPar
Type
&\sphinxstyletheadfamily 
\sphinxAtStartPar
Description
\\
\sphinxmidrule
\sphinxtableatstartofbodyhook
\sphinxAtStartPar
\sphinxcode{\sphinxupquote{Detection\_ID}}
&
\sphinxAtStartPar
\sphinxcode{\sphinxupquote{int}}
&
\sphinxAtStartPar
Clé primaire.
\\
\sphinxhline
\sphinxAtStartPar
\sphinxcode{\sphinxupquote{Precursor\_MZ}}
&
\sphinxAtStartPar
\sphinxcode{\sphinxupquote{float}}
&
\sphinxAtStartPar
m/z de l’ion précurseur (Da).
\\
\sphinxhline
\sphinxAtStartPar
\sphinxcode{\sphinxupquote{MS\_level}}
&
\sphinxAtStartPar
\sphinxcode{\sphinxupquote{str}}
&
\sphinxAtStartPar
\sphinxcode{\sphinxupquote{"MS1"}} (précurseur seul) ou \sphinxcode{\sphinxupquote{"MS2"}} (avec spectre de fragmentation).
\\
\sphinxhline
\sphinxAtStartPar
\sphinxcode{\sphinxupquote{Ionisation\_mode}}
&
\sphinxAtStartPar
\sphinxcode{\sphinxupquote{str}}
&
\sphinxAtStartPar
\sphinxcode{\sphinxupquote{"Positive"}} ou \sphinxcode{\sphinxupquote{"Negative"}}.
\\
\sphinxhline
\sphinxAtStartPar
\sphinxcode{\sphinxupquote{Adduct}}
&
\sphinxAtStartPar
\sphinxcode{\sphinxupquote{str | None}}
&
\sphinxAtStartPar
Adduit du précurseur.
\\
\sphinxhline
\sphinxAtStartPar
\sphinxcode{\sphinxupquote{Num\_Peaks}}
&
\sphinxAtStartPar
\sphinxcode{\sphinxupquote{int | None}}
&
\sphinxAtStartPar
Nombre de pics de fragments (MS2 uniquement).
\\
\sphinxhline
\sphinxAtStartPar
\sphinxcode{\sphinxupquote{Neutral\_mass}}
&
\sphinxAtStartPar
\sphinxcode{\sphinxupquote{float}}
&
\sphinxAtStartPar
Masse neutre dérivée de \sphinxcode{\sphinxupquote{Precursor\_MZ}} et du mode d’ionisation (Da).
\\
\sphinxhline
\sphinxAtStartPar
\sphinxcode{\sphinxupquote{RT}}
&
\sphinxAtStartPar
\sphinxcode{\sphinxupquote{float | None}}
&
\sphinxAtStartPar
Temps de rétention (minutes).
\\
\sphinxhline
\sphinxAtStartPar
\sphinxcode{\sphinxupquote{CCS}}
&
\sphinxAtStartPar
\sphinxcode{\sphinxupquote{float | None}}
&
\sphinxAtStartPar
Section efficace de collision (Ų).
\\
\sphinxbottomrule
\end{tabular}
\sphinxtableafterendhook\par
\sphinxattableend\end{savenotes}


\section{Table \sphinxstyleliteralintitle{\sphinxupquote{Fragment}}}
\label{\detokenize{model:table-fragment}}

\begin{savenotes}\sphinxattablestart
\sphinxthistablewithglobalstyle
\centering
\begin{tabular}[t]{\X{25}{100}\X{20}{100}\X{55}{100}}
\sphinxtoprule
\sphinxstyletheadfamily 
\sphinxAtStartPar
Colonne
&\sphinxstyletheadfamily 
\sphinxAtStartPar
Type
&\sphinxstyletheadfamily 
\sphinxAtStartPar
Description
\\
\sphinxmidrule
\sphinxtableatstartofbodyhook
\sphinxAtStartPar
\sphinxcode{\sphinxupquote{Fragment\_ID}}
&
\sphinxAtStartPar
\sphinxcode{\sphinxupquote{int}}
&
\sphinxAtStartPar
Clé primaire.
\\
\sphinxhline
\sphinxAtStartPar
\sphinxcode{\sphinxupquote{Detection\_id}}
&
\sphinxAtStartPar
\sphinxcode{\sphinxupquote{int}}
&
\sphinxAtStartPar
Clé étrangère vers \sphinxcode{\sphinxupquote{Detection.Detection\_ID}}.
\\
\sphinxhline
\sphinxAtStartPar
\sphinxcode{\sphinxupquote{MZ}}
&
\sphinxAtStartPar
\sphinxcode{\sphinxupquote{float | None}}
&
\sphinxAtStartPar
m/z de l’ion fragment (Da).
\\
\sphinxhline
\sphinxAtStartPar
\sphinxcode{\sphinxupquote{Intensity}}
&
\sphinxAtStartPar
\sphinxcode{\sphinxupquote{float | None}}
&
\sphinxAtStartPar
Intensité du pic de fragment.
\\
\sphinxbottomrule
\end{tabular}
\sphinxtableafterendhook\par
\sphinxattableend\end{savenotes}


\section{Table \sphinxstyleliteralintitle{\sphinxupquote{Lipid}}}
\label{\detokenize{model:table-lipid}}

\begin{savenotes}\sphinxattablestart
\sphinxthistablewithglobalstyle
\centering
\begin{tabular}[t]{\X{25}{100}\X{20}{100}\X{55}{100}}
\sphinxtoprule
\sphinxstyletheadfamily 
\sphinxAtStartPar
Colonne
&\sphinxstyletheadfamily 
\sphinxAtStartPar
Type
&\sphinxstyletheadfamily 
\sphinxAtStartPar
Description
\\
\sphinxmidrule
\sphinxtableatstartofbodyhook
\sphinxAtStartPar
\sphinxcode{\sphinxupquote{Lipid\_ID}}
&
\sphinxAtStartPar
\sphinxcode{\sphinxupquote{int}}
&
\sphinxAtStartPar
Clé primaire.
\\
\sphinxhline
\sphinxAtStartPar
\sphinxcode{\sphinxupquote{Lipid\_name}}
&
\sphinxAtStartPar
\sphinxcode{\sphinxupquote{str}}
&
\sphinxAtStartPar
Nom du lipide.
\\
\sphinxhline
\sphinxAtStartPar
\sphinxcode{\sphinxupquote{Lipid\_category}}
&
\sphinxAtStartPar
\sphinxcode{\sphinxupquote{str | None}}
&
\sphinxAtStartPar
Catégorie LIPID MAPS.
\\
\sphinxhline
\sphinxAtStartPar
\sphinxcode{\sphinxupquote{Lipid\_class}}
&
\sphinxAtStartPar
\sphinxcode{\sphinxupquote{str | None}}
&
\sphinxAtStartPar
Classe LIPID MAPS.
\\
\sphinxhline
\sphinxAtStartPar
\sphinxcode{\sphinxupquote{Lipid\_subclass}}
&
\sphinxAtStartPar
\sphinxcode{\sphinxupquote{str | None}}
&
\sphinxAtStartPar
Sous\sphinxhyphen{}classe LIPID MAPS.
\\
\sphinxhline
\sphinxAtStartPar
\sphinxcode{\sphinxupquote{Formula}}
&
\sphinxAtStartPar
\sphinxcode{\sphinxupquote{str}}
&
\sphinxAtStartPar
Formule chimique.
\\
\sphinxhline
\sphinxAtStartPar
\sphinxcode{\sphinxupquote{Molecular\_weight}}
&
\sphinxAtStartPar
\sphinxcode{\sphinxupquote{float}}
&
\sphinxAtStartPar
Masse moléculaire moyenne (Da).
\\
\sphinxhline
\sphinxAtStartPar
\sphinxcode{\sphinxupquote{Monoisotopic\_mass}}
&
\sphinxAtStartPar
\sphinxcode{\sphinxupquote{float}}
&
\sphinxAtStartPar
Masse monoisotopique (Da).
\\
\sphinxbottomrule
\end{tabular}
\sphinxtableafterendhook\par
\sphinxattableend\end{savenotes}


\section{Table \sphinxstyleliteralintitle{\sphinxupquote{Annotation}}}
\label{\detokenize{model:table-annotation}}

\begin{savenotes}\sphinxattablestart
\sphinxthistablewithglobalstyle
\centering
\begin{tabular}[t]{\X{25}{100}\X{20}{100}\X{55}{100}}
\sphinxtoprule
\sphinxstyletheadfamily 
\sphinxAtStartPar
Colonne
&\sphinxstyletheadfamily 
\sphinxAtStartPar
Type
&\sphinxstyletheadfamily 
\sphinxAtStartPar
Description
\\
\sphinxmidrule
\sphinxtableatstartofbodyhook
\sphinxAtStartPar
\sphinxcode{\sphinxupquote{Annotation\_ID}}
&
\sphinxAtStartPar
\sphinxcode{\sphinxupquote{int}}
&
\sphinxAtStartPar
Clé primaire.
\\
\sphinxhline
\sphinxAtStartPar
\sphinxcode{\sphinxupquote{Lipid\_id}}
&
\sphinxAtStartPar
\sphinxcode{\sphinxupquote{int}}
&
\sphinxAtStartPar
Clé étrangère (unique) vers \sphinxcode{\sphinxupquote{Lipid.Lipid\_ID}}.
\\
\sphinxhline
\sphinxAtStartPar
\sphinxcode{\sphinxupquote{Detection\_id}}
&
\sphinxAtStartPar
\sphinxcode{\sphinxupquote{int}}
&
\sphinxAtStartPar
Clé étrangère (unique) vers \sphinxcode{\sphinxupquote{Detection.Detection\_ID}}.
\\
\sphinxbottomrule
\end{tabular}
\sphinxtableafterendhook\par
\sphinxattableend\end{savenotes}


\section{Conventions}
\label{\detokenize{model:conventions}}\begin{itemize}
\item {} 
\sphinxAtStartPar
Les clés primaires sont nommées avec le suffixe \sphinxcode{\sphinxupquote{\_ID}}.

\item {} 
\sphinxAtStartPar
Les clés étrangères sont nommées avec le suffixe \sphinxcode{\sphinxupquote{\_id}}.

\item {} 
\sphinxAtStartPar
Chaque classe définit un \sphinxcode{\sphinxupquote{\_\_repr\_\_}} pour afficher un objet de façon lisible lors du débogage.

\end{itemize}


\section{Voir aussi}
\label{\detokenize{model:voir-aussi}}
\sphinxAtStartPar
Documentation SQLAlchemy :
\begin{itemize}
\item {} 
\sphinxAtStartPar
\sphinxhref{https://docs.sqlalchemy.org/en/21/orm/declarative\_styles.html}{DeclarativeBase}

\item {} 
\sphinxAtStartPar
\sphinxhref{https://docs.sqlalchemy.org/en/21/orm/mapping\_styles.html\#orm-mapping-styles}{Mapped}

\item {} 
\sphinxAtStartPar
\sphinxhref{https://docs.sqlalchemy.org/en/21/orm/mapping\_api.html\#sqlalchemy.orm.Mapper.columns}{mapped\_column}

\item {} 
\sphinxAtStartPar
\sphinxhref{https://docs.sqlalchemy.org/en/20/orm/basic\_relationships.html}{relationship, ForeignKey}

\end{itemize}

\sphinxstepscope


\chapter{Référence API (générée automatiquement)}
\label{\detokenize{api:reference-api-generee-automatiquement}}\label{\detokenize{api::doc}}
\begin{sphinxadmonition}{note}{Note:}
\sphinxAtStartPar
Cette page est générée automatiquement à partir des docstrings du code source
(\sphinxcode{\sphinxupquote{sphinx.ext.autodoc}}). Contrairement aux pages Streamlit ({\hyperref[\detokenize{Home::doc}]{\sphinxcrossref{\DUrole{doc}{Home}}}},
{\hyperref[\detokenize{1_Integration::doc}]{\sphinxcrossref{\DUrole{doc}{Data\_Integration}}}}, {\hyperref[\detokenize{2_BacLipidDB::doc}]{\sphinxcrossref{\DUrole{doc}{BacLipidDB}}}}, {\hyperref[\detokenize{3_Export::doc}]{\sphinxcrossref{\DUrole{doc}{Export}}}}), les modules
ci\sphinxhyphen{}dessous sont écrits sous forme de fonctions/classes documentées, ce qui
permet une extraction automatique.
\end{sphinxadmonition}


\section{config}
\label{\detokenize{api:module-config}}\label{\detokenize{api:config}}\index{module@\spxentry{module}!config@\spxentry{config}}\index{config@\spxentry{config}!module@\spxentry{module}}\index{get\_engine() (in module config)@\spxentry{get\_engine()}\spxextra{in module config}}

\begin{fulllineitems}
\phantomsection\label{\detokenize{api:config.get_engine}}
\pysigstartsignatures
\pysiglinewithargsret
{\sphinxcode{\sphinxupquote{config.}}\sphinxbfcode{\sphinxupquote{get\_engine}}}
{}
{}
\pysigstopsignatures
\sphinxAtStartPar
Create and cache the SQLAlchemy engine for the SQLite database connection.

\end{fulllineitems}



\section{models.model}
\label{\detokenize{api:module-models.model}}\label{\detokenize{api:models-model}}\index{module@\spxentry{module}!models.model@\spxentry{models.model}}\index{models.model@\spxentry{models.model}!module@\spxentry{module}}\index{Base (class in models.model)@\spxentry{Base}\spxextra{class in models.model}}

\begin{fulllineitems}
\phantomsection\label{\detokenize{api:models.model.Base}}
\pysigstartsignatures
\pysiglinewithargsret
{\sphinxbfcode{\sphinxupquote{\DUrole{k}{class}\DUrole{w}{ }}}\sphinxcode{\sphinxupquote{models.model.}}\sphinxbfcode{\sphinxupquote{Base}}}
{\sphinxparam{\DUrole{o}{**}\DUrole{n}{kwargs}\DUrole{p}{:}\DUrole{w}{ }\DUrole{n}{Any}}}
{}
\pysigstopsignatures
\sphinxAtStartPar
Bases: \sphinxcode{\sphinxupquote{DeclarativeBase}}

\sphinxAtStartPar
Declarative base class for all BacLipidDB ORM models.

\sphinxAtStartPar
All mapped classes ({\hyperref[\detokenize{api:models.model.Detection}]{\sphinxcrossref{\sphinxcode{\sphinxupquote{Detection}}}}}, {\hyperref[\detokenize{api:models.model.Fragment}]{\sphinxcrossref{\sphinxcode{\sphinxupquote{Fragment}}}}}, {\hyperref[\detokenize{api:models.model.Lipid}]{\sphinxcrossref{\sphinxcode{\sphinxupquote{Lipid}}}}},
{\hyperref[\detokenize{api:models.model.Annotation}]{\sphinxcrossref{\sphinxcode{\sphinxupquote{Annotation}}}}}) inherit from this class so SQLAlchemy can collect
their metadata and generate the corresponding database tables.
\index{registry (models.model.Base attribute)@\spxentry{registry}\spxextra{models.model.Base attribute}}

\begin{fulllineitems}
\phantomsection\label{\detokenize{api:models.model.Base.registry}}
\pysigstartsignatures
\pysigline
{\sphinxbfcode{\sphinxupquote{registry}}\sphinxbfcode{\sphinxupquote{\DUrole{p}{:}\DUrole{w}{ }ClassVar\DUrole{p}{{[}}\_RegistryType\DUrole{p}{{]}}}}\sphinxbfcode{\sphinxupquote{\DUrole{w}{ }\DUrole{p}{=}\DUrole{w}{ }\textless{}sqlalchemy.orm.decl\_api.registry object\textgreater{}}}}
\pysigstopsignatures
\sphinxAtStartPar
Refers to the \sphinxcode{\sphinxupquote{\_orm.registry}} in use where new
\sphinxcode{\sphinxupquote{\_orm.Mapper}} objects will be associated.

\end{fulllineitems}


\end{fulllineitems}

\index{Detection (class in models.model)@\spxentry{Detection}\spxextra{class in models.model}}

\begin{fulllineitems}
\phantomsection\label{\detokenize{api:models.model.Detection}}
\pysigstartsignatures
\pysiglinewithargsret
{\sphinxbfcode{\sphinxupquote{\DUrole{k}{class}\DUrole{w}{ }}}\sphinxcode{\sphinxupquote{models.model.}}\sphinxbfcode{\sphinxupquote{Detection}}}
{\sphinxparam{\DUrole{o}{**}\DUrole{n}{kwargs}}}
{}
\pysigstopsignatures
\sphinxAtStartPar
Bases: {\hyperref[\detokenize{api:models.model.Base}]{\sphinxcrossref{\sphinxcode{\sphinxupquote{Base}}}}}

\sphinxAtStartPar
A single MS1 or MS2 detection (ion observed in a run).

\sphinxAtStartPar
Stores the precursor m/z, ionisation mode, adduct, retention time and
CCS. Linked to zero or more {\hyperref[\detokenize{api:models.model.Fragment}]{\sphinxcrossref{\sphinxcode{\sphinxupquote{Fragment}}}}} rows when \sphinxcode{\sphinxupquote{MS\_level}} is
“MS2”, and to one {\hyperref[\detokenize{api:models.model.Annotation}]{\sphinxcrossref{\sphinxcode{\sphinxupquote{Annotation}}}}}.
\index{Detection\_ID (models.model.Detection attribute)@\spxentry{Detection\_ID}\spxextra{models.model.Detection attribute}}

\begin{fulllineitems}
\phantomsection\label{\detokenize{api:models.model.Detection.Detection_ID}}
\pysigstartsignatures
\pysigline
{\sphinxbfcode{\sphinxupquote{Detection\_ID}}\sphinxbfcode{\sphinxupquote{\DUrole{p}{:}\DUrole{w}{ }Mapped\DUrole{p}{{[}}int\DUrole{p}{{]}}}}}
\pysigstopsignatures
\end{fulllineitems}

\index{Precursor\_MZ (models.model.Detection attribute)@\spxentry{Precursor\_MZ}\spxextra{models.model.Detection attribute}}

\begin{fulllineitems}
\phantomsection\label{\detokenize{api:models.model.Detection.Precursor_MZ}}
\pysigstartsignatures
\pysigline
{\sphinxbfcode{\sphinxupquote{Precursor\_MZ}}\sphinxbfcode{\sphinxupquote{\DUrole{p}{:}\DUrole{w}{ }Mapped\DUrole{p}{{[}}float\DUrole{p}{{]}}}}}
\pysigstopsignatures
\end{fulllineitems}

\index{MS\_level (models.model.Detection attribute)@\spxentry{MS\_level}\spxextra{models.model.Detection attribute}}

\begin{fulllineitems}
\phantomsection\label{\detokenize{api:models.model.Detection.MS_level}}
\pysigstartsignatures
\pysigline
{\sphinxbfcode{\sphinxupquote{MS\_level}}\sphinxbfcode{\sphinxupquote{\DUrole{p}{:}\DUrole{w}{ }Mapped\DUrole{p}{{[}}str\DUrole{p}{{]}}}}}
\pysigstopsignatures
\end{fulllineitems}

\index{Ionisation\_mode (models.model.Detection attribute)@\spxentry{Ionisation\_mode}\spxextra{models.model.Detection attribute}}

\begin{fulllineitems}
\phantomsection\label{\detokenize{api:models.model.Detection.Ionisation_mode}}
\pysigstartsignatures
\pysigline
{\sphinxbfcode{\sphinxupquote{Ionisation\_mode}}\sphinxbfcode{\sphinxupquote{\DUrole{p}{:}\DUrole{w}{ }Mapped\DUrole{p}{{[}}str\DUrole{p}{{]}}}}}
\pysigstopsignatures
\end{fulllineitems}

\index{Adduct (models.model.Detection attribute)@\spxentry{Adduct}\spxextra{models.model.Detection attribute}}

\begin{fulllineitems}
\phantomsection\label{\detokenize{api:models.model.Detection.Adduct}}
\pysigstartsignatures
\pysigline
{\sphinxbfcode{\sphinxupquote{Adduct}}\sphinxbfcode{\sphinxupquote{\DUrole{p}{:}\DUrole{w}{ }Mapped\DUrole{p}{{[}}str\DUrole{p}{{]}}}}}
\pysigstopsignatures
\end{fulllineitems}

\index{Num\_Peaks (models.model.Detection attribute)@\spxentry{Num\_Peaks}\spxextra{models.model.Detection attribute}}

\begin{fulllineitems}
\phantomsection\label{\detokenize{api:models.model.Detection.Num_Peaks}}
\pysigstartsignatures
\pysigline
{\sphinxbfcode{\sphinxupquote{Num\_Peaks}}\sphinxbfcode{\sphinxupquote{\DUrole{p}{:}\DUrole{w}{ }Mapped\DUrole{p}{{[}}int\DUrole{w}{ }\DUrole{p}{|}\DUrole{w}{ }None\DUrole{p}{{]}}}}}
\pysigstopsignatures
\end{fulllineitems}

\index{Neutral\_mass (models.model.Detection attribute)@\spxentry{Neutral\_mass}\spxextra{models.model.Detection attribute}}

\begin{fulllineitems}
\phantomsection\label{\detokenize{api:models.model.Detection.Neutral_mass}}
\pysigstartsignatures
\pysigline
{\sphinxbfcode{\sphinxupquote{Neutral\_mass}}\sphinxbfcode{\sphinxupquote{\DUrole{p}{:}\DUrole{w}{ }Mapped\DUrole{p}{{[}}float\DUrole{p}{{]}}}}}
\pysigstopsignatures
\end{fulllineitems}

\index{RT (models.model.Detection attribute)@\spxentry{RT}\spxextra{models.model.Detection attribute}}

\begin{fulllineitems}
\phantomsection\label{\detokenize{api:models.model.Detection.RT}}
\pysigstartsignatures
\pysigline
{\sphinxbfcode{\sphinxupquote{RT}}\sphinxbfcode{\sphinxupquote{\DUrole{p}{:}\DUrole{w}{ }Mapped\DUrole{p}{{[}}float\DUrole{w}{ }\DUrole{p}{|}\DUrole{w}{ }None\DUrole{p}{{]}}}}}
\pysigstopsignatures
\end{fulllineitems}

\index{CCS (models.model.Detection attribute)@\spxentry{CCS}\spxextra{models.model.Detection attribute}}

\begin{fulllineitems}
\phantomsection\label{\detokenize{api:models.model.Detection.CCS}}
\pysigstartsignatures
\pysigline
{\sphinxbfcode{\sphinxupquote{CCS}}\sphinxbfcode{\sphinxupquote{\DUrole{p}{:}\DUrole{w}{ }Mapped\DUrole{p}{{[}}float\DUrole{w}{ }\DUrole{p}{|}\DUrole{w}{ }None\DUrole{p}{{]}}}}}
\pysigstopsignatures
\end{fulllineitems}

\index{fragments (models.model.Detection attribute)@\spxentry{fragments}\spxextra{models.model.Detection attribute}}

\begin{fulllineitems}
\phantomsection\label{\detokenize{api:models.model.Detection.fragments}}
\pysigstartsignatures
\pysigline
{\sphinxbfcode{\sphinxupquote{fragments}}\sphinxbfcode{\sphinxupquote{\DUrole{p}{:}\DUrole{w}{ }Mapped\DUrole{p}{{[}}list\DUrole{p}{{[}}{\hyperref[\detokenize{api:models.model.Fragment}]{\sphinxcrossref{Fragment}}}\DUrole{p}{{]}}\DUrole{p}{{]}}}}}
\pysigstopsignatures
\end{fulllineitems}

\index{annotation (models.model.Detection attribute)@\spxentry{annotation}\spxextra{models.model.Detection attribute}}

\begin{fulllineitems}
\phantomsection\label{\detokenize{api:models.model.Detection.annotation}}
\pysigstartsignatures
\pysigline
{\sphinxbfcode{\sphinxupquote{annotation}}\sphinxbfcode{\sphinxupquote{\DUrole{p}{:}\DUrole{w}{ }Mapped\DUrole{p}{{[}}{\hyperref[\detokenize{api:models.model.Annotation}]{\sphinxcrossref{Annotation}}}\DUrole{p}{{]}}}}}
\pysigstopsignatures
\end{fulllineitems}


\end{fulllineitems}

\index{Fragment (class in models.model)@\spxentry{Fragment}\spxextra{class in models.model}}

\begin{fulllineitems}
\phantomsection\label{\detokenize{api:models.model.Fragment}}
\pysigstartsignatures
\pysiglinewithargsret
{\sphinxbfcode{\sphinxupquote{\DUrole{k}{class}\DUrole{w}{ }}}\sphinxcode{\sphinxupquote{models.model.}}\sphinxbfcode{\sphinxupquote{Fragment}}}
{\sphinxparam{\DUrole{o}{**}\DUrole{n}{kwargs}}}
{}
\pysigstopsignatures
\sphinxAtStartPar
Bases: {\hyperref[\detokenize{api:models.model.Base}]{\sphinxcrossref{\sphinxcode{\sphinxupquote{Base}}}}}

\sphinxAtStartPar
A fragment ion peak (m/z, intensity) belonging to an MS2 spectrum.

\sphinxAtStartPar
Each row is attached to the {\hyperref[\detokenize{api:models.model.Detection}]{\sphinxcrossref{\sphinxcode{\sphinxupquote{Detection}}}}} it was measured in.
\index{Fragment\_ID (models.model.Fragment attribute)@\spxentry{Fragment\_ID}\spxextra{models.model.Fragment attribute}}

\begin{fulllineitems}
\phantomsection\label{\detokenize{api:models.model.Fragment.Fragment_ID}}
\pysigstartsignatures
\pysigline
{\sphinxbfcode{\sphinxupquote{Fragment\_ID}}\sphinxbfcode{\sphinxupquote{\DUrole{p}{:}\DUrole{w}{ }Mapped\DUrole{p}{{[}}int\DUrole{p}{{]}}}}}
\pysigstopsignatures
\end{fulllineitems}

\index{Detection\_id (models.model.Fragment attribute)@\spxentry{Detection\_id}\spxextra{models.model.Fragment attribute}}

\begin{fulllineitems}
\phantomsection\label{\detokenize{api:models.model.Fragment.Detection_id}}
\pysigstartsignatures
\pysigline
{\sphinxbfcode{\sphinxupquote{Detection\_id}}\sphinxbfcode{\sphinxupquote{\DUrole{p}{:}\DUrole{w}{ }Mapped\DUrole{p}{{[}}int\DUrole{p}{{]}}}}}
\pysigstopsignatures
\end{fulllineitems}

\index{MZ (models.model.Fragment attribute)@\spxentry{MZ}\spxextra{models.model.Fragment attribute}}

\begin{fulllineitems}
\phantomsection\label{\detokenize{api:models.model.Fragment.MZ}}
\pysigstartsignatures
\pysigline
{\sphinxbfcode{\sphinxupquote{MZ}}\sphinxbfcode{\sphinxupquote{\DUrole{p}{:}\DUrole{w}{ }Mapped\DUrole{p}{{[}}float\DUrole{w}{ }\DUrole{p}{|}\DUrole{w}{ }None\DUrole{p}{{]}}}}}
\pysigstopsignatures
\end{fulllineitems}

\index{Intensity (models.model.Fragment attribute)@\spxentry{Intensity}\spxextra{models.model.Fragment attribute}}

\begin{fulllineitems}
\phantomsection\label{\detokenize{api:models.model.Fragment.Intensity}}
\pysigstartsignatures
\pysigline
{\sphinxbfcode{\sphinxupquote{Intensity}}\sphinxbfcode{\sphinxupquote{\DUrole{p}{:}\DUrole{w}{ }Mapped\DUrole{p}{{[}}float\DUrole{w}{ }\DUrole{p}{|}\DUrole{w}{ }None\DUrole{p}{{]}}}}}
\pysigstopsignatures
\end{fulllineitems}

\index{detection (models.model.Fragment attribute)@\spxentry{detection}\spxextra{models.model.Fragment attribute}}

\begin{fulllineitems}
\phantomsection\label{\detokenize{api:models.model.Fragment.detection}}
\pysigstartsignatures
\pysigline
{\sphinxbfcode{\sphinxupquote{detection}}\sphinxbfcode{\sphinxupquote{\DUrole{p}{:}\DUrole{w}{ }Mapped\DUrole{p}{{[}}{\hyperref[\detokenize{api:models.model.Detection}]{\sphinxcrossref{Detection}}}\DUrole{p}{{]}}}}}
\pysigstopsignatures
\end{fulllineitems}


\end{fulllineitems}

\index{Lipid (class in models.model)@\spxentry{Lipid}\spxextra{class in models.model}}

\begin{fulllineitems}
\phantomsection\label{\detokenize{api:models.model.Lipid}}
\pysigstartsignatures
\pysiglinewithargsret
{\sphinxbfcode{\sphinxupquote{\DUrole{k}{class}\DUrole{w}{ }}}\sphinxcode{\sphinxupquote{models.model.}}\sphinxbfcode{\sphinxupquote{Lipid}}}
{\sphinxparam{\DUrole{o}{**}\DUrole{n}{kwargs}}}
{}
\pysigstopsignatures
\sphinxAtStartPar
Bases: {\hyperref[\detokenize{api:models.model.Base}]{\sphinxcrossref{\sphinxcode{\sphinxupquote{Base}}}}}

\sphinxAtStartPar
A lipid entity identified by name, formula and LIPID MAPS classification.

\sphinxAtStartPar
Independent of any given detection; linked to the {\hyperref[\detokenize{api:models.model.Detection}]{\sphinxcrossref{\sphinxcode{\sphinxupquote{Detection}}}}} it
was identified in through a single {\hyperref[\detokenize{api:models.model.Annotation}]{\sphinxcrossref{\sphinxcode{\sphinxupquote{Annotation}}}}}.
\index{Lipid\_ID (models.model.Lipid attribute)@\spxentry{Lipid\_ID}\spxextra{models.model.Lipid attribute}}

\begin{fulllineitems}
\phantomsection\label{\detokenize{api:models.model.Lipid.Lipid_ID}}
\pysigstartsignatures
\pysigline
{\sphinxbfcode{\sphinxupquote{Lipid\_ID}}\sphinxbfcode{\sphinxupquote{\DUrole{p}{:}\DUrole{w}{ }Mapped\DUrole{p}{{[}}int\DUrole{p}{{]}}}}}
\pysigstopsignatures
\end{fulllineitems}

\index{Lipid\_name (models.model.Lipid attribute)@\spxentry{Lipid\_name}\spxextra{models.model.Lipid attribute}}

\begin{fulllineitems}
\phantomsection\label{\detokenize{api:models.model.Lipid.Lipid_name}}
\pysigstartsignatures
\pysigline
{\sphinxbfcode{\sphinxupquote{Lipid\_name}}\sphinxbfcode{\sphinxupquote{\DUrole{p}{:}\DUrole{w}{ }Mapped\DUrole{p}{{[}}str\DUrole{p}{{]}}}}}
\pysigstopsignatures
\end{fulllineitems}

\index{Lipid\_category (models.model.Lipid attribute)@\spxentry{Lipid\_category}\spxextra{models.model.Lipid attribute}}

\begin{fulllineitems}
\phantomsection\label{\detokenize{api:models.model.Lipid.Lipid_category}}
\pysigstartsignatures
\pysigline
{\sphinxbfcode{\sphinxupquote{Lipid\_category}}\sphinxbfcode{\sphinxupquote{\DUrole{p}{:}\DUrole{w}{ }Mapped\DUrole{p}{{[}}str\DUrole{w}{ }\DUrole{p}{|}\DUrole{w}{ }None\DUrole{p}{{]}}}}}
\pysigstopsignatures
\end{fulllineitems}

\index{Lipid\_class (models.model.Lipid attribute)@\spxentry{Lipid\_class}\spxextra{models.model.Lipid attribute}}

\begin{fulllineitems}
\phantomsection\label{\detokenize{api:models.model.Lipid.Lipid_class}}
\pysigstartsignatures
\pysigline
{\sphinxbfcode{\sphinxupquote{Lipid\_class}}\sphinxbfcode{\sphinxupquote{\DUrole{p}{:}\DUrole{w}{ }Mapped\DUrole{p}{{[}}str\DUrole{w}{ }\DUrole{p}{|}\DUrole{w}{ }None\DUrole{p}{{]}}}}}
\pysigstopsignatures
\end{fulllineitems}

\index{Lipid\_subclass (models.model.Lipid attribute)@\spxentry{Lipid\_subclass}\spxextra{models.model.Lipid attribute}}

\begin{fulllineitems}
\phantomsection\label{\detokenize{api:models.model.Lipid.Lipid_subclass}}
\pysigstartsignatures
\pysigline
{\sphinxbfcode{\sphinxupquote{Lipid\_subclass}}\sphinxbfcode{\sphinxupquote{\DUrole{p}{:}\DUrole{w}{ }Mapped\DUrole{p}{{[}}str\DUrole{w}{ }\DUrole{p}{|}\DUrole{w}{ }None\DUrole{p}{{]}}}}}
\pysigstopsignatures
\end{fulllineitems}

\index{Formula (models.model.Lipid attribute)@\spxentry{Formula}\spxextra{models.model.Lipid attribute}}

\begin{fulllineitems}
\phantomsection\label{\detokenize{api:models.model.Lipid.Formula}}
\pysigstartsignatures
\pysigline
{\sphinxbfcode{\sphinxupquote{Formula}}\sphinxbfcode{\sphinxupquote{\DUrole{p}{:}\DUrole{w}{ }Mapped\DUrole{p}{{[}}str\DUrole{p}{{]}}}}}
\pysigstopsignatures
\end{fulllineitems}

\index{FA\_composition (models.model.Lipid attribute)@\spxentry{FA\_composition}\spxextra{models.model.Lipid attribute}}

\begin{fulllineitems}
\phantomsection\label{\detokenize{api:models.model.Lipid.FA_composition}}
\pysigstartsignatures
\pysigline
{\sphinxbfcode{\sphinxupquote{FA\_composition}}\sphinxbfcode{\sphinxupquote{\DUrole{p}{:}\DUrole{w}{ }Mapped\DUrole{p}{{[}}str\DUrole{w}{ }\DUrole{p}{|}\DUrole{w}{ }None\DUrole{p}{{]}}}}}
\pysigstopsignatures
\end{fulllineitems}

\index{Molecular\_weight (models.model.Lipid attribute)@\spxentry{Molecular\_weight}\spxextra{models.model.Lipid attribute}}

\begin{fulllineitems}
\phantomsection\label{\detokenize{api:models.model.Lipid.Molecular_weight}}
\pysigstartsignatures
\pysigline
{\sphinxbfcode{\sphinxupquote{Molecular\_weight}}\sphinxbfcode{\sphinxupquote{\DUrole{p}{:}\DUrole{w}{ }Mapped\DUrole{p}{{[}}float\DUrole{p}{{]}}}}}
\pysigstopsignatures
\end{fulllineitems}

\index{Monoisotopic\_mass (models.model.Lipid attribute)@\spxentry{Monoisotopic\_mass}\spxextra{models.model.Lipid attribute}}

\begin{fulllineitems}
\phantomsection\label{\detokenize{api:models.model.Lipid.Monoisotopic_mass}}
\pysigstartsignatures
\pysigline
{\sphinxbfcode{\sphinxupquote{Monoisotopic\_mass}}\sphinxbfcode{\sphinxupquote{\DUrole{p}{:}\DUrole{w}{ }Mapped\DUrole{p}{{[}}float\DUrole{p}{{]}}}}}
\pysigstopsignatures
\end{fulllineitems}

\index{annotation (models.model.Lipid attribute)@\spxentry{annotation}\spxextra{models.model.Lipid attribute}}

\begin{fulllineitems}
\phantomsection\label{\detokenize{api:models.model.Lipid.annotation}}
\pysigstartsignatures
\pysigline
{\sphinxbfcode{\sphinxupquote{annotation}}\sphinxbfcode{\sphinxupquote{\DUrole{p}{:}\DUrole{w}{ }Mapped\DUrole{p}{{[}}{\hyperref[\detokenize{api:models.model.Annotation}]{\sphinxcrossref{Annotation}}}\DUrole{p}{{]}}}}}
\pysigstopsignatures
\end{fulllineitems}


\end{fulllineitems}

\index{Annotation (class in models.model)@\spxentry{Annotation}\spxextra{class in models.model}}

\begin{fulllineitems}
\phantomsection\label{\detokenize{api:models.model.Annotation}}
\pysigstartsignatures
\pysiglinewithargsret
{\sphinxbfcode{\sphinxupquote{\DUrole{k}{class}\DUrole{w}{ }}}\sphinxcode{\sphinxupquote{models.model.}}\sphinxbfcode{\sphinxupquote{Annotation}}}
{\sphinxparam{\DUrole{o}{**}\DUrole{n}{kwargs}}}
{}
\pysigstopsignatures
\sphinxAtStartPar
Bases: {\hyperref[\detokenize{api:models.model.Base}]{\sphinxcrossref{\sphinxcode{\sphinxupquote{Base}}}}}

\sphinxAtStartPar
Join table linking one {\hyperref[\detokenize{api:models.model.Lipid}]{\sphinxcrossref{\sphinxcode{\sphinxupquote{Lipid}}}}} to one {\hyperref[\detokenize{api:models.model.Detection}]{\sphinxcrossref{\sphinxcode{\sphinxupquote{Detection}}}}}.

\sphinxAtStartPar
Represents the identification of a detected ion as a given lipid.
\index{Annotation\_ID (models.model.Annotation attribute)@\spxentry{Annotation\_ID}\spxextra{models.model.Annotation attribute}}

\begin{fulllineitems}
\phantomsection\label{\detokenize{api:models.model.Annotation.Annotation_ID}}
\pysigstartsignatures
\pysigline
{\sphinxbfcode{\sphinxupquote{Annotation\_ID}}\sphinxbfcode{\sphinxupquote{\DUrole{p}{:}\DUrole{w}{ }Mapped\DUrole{p}{{[}}int\DUrole{p}{{]}}}}}
\pysigstopsignatures
\end{fulllineitems}

\index{Lipid\_id (models.model.Annotation attribute)@\spxentry{Lipid\_id}\spxextra{models.model.Annotation attribute}}

\begin{fulllineitems}
\phantomsection\label{\detokenize{api:models.model.Annotation.Lipid_id}}
\pysigstartsignatures
\pysigline
{\sphinxbfcode{\sphinxupquote{Lipid\_id}}\sphinxbfcode{\sphinxupquote{\DUrole{p}{:}\DUrole{w}{ }Mapped\DUrole{p}{{[}}int\DUrole{p}{{]}}}}}
\pysigstopsignatures
\end{fulllineitems}

\index{Detection\_id (models.model.Annotation attribute)@\spxentry{Detection\_id}\spxextra{models.model.Annotation attribute}}

\begin{fulllineitems}
\phantomsection\label{\detokenize{api:models.model.Annotation.Detection_id}}
\pysigstartsignatures
\pysigline
{\sphinxbfcode{\sphinxupquote{Detection\_id}}\sphinxbfcode{\sphinxupquote{\DUrole{p}{:}\DUrole{w}{ }Mapped\DUrole{p}{{[}}int\DUrole{p}{{]}}}}}
\pysigstopsignatures
\end{fulllineitems}

\index{lipid (models.model.Annotation attribute)@\spxentry{lipid}\spxextra{models.model.Annotation attribute}}

\begin{fulllineitems}
\phantomsection\label{\detokenize{api:models.model.Annotation.lipid}}
\pysigstartsignatures
\pysigline
{\sphinxbfcode{\sphinxupquote{lipid}}\sphinxbfcode{\sphinxupquote{\DUrole{p}{:}\DUrole{w}{ }Mapped\DUrole{p}{{[}}{\hyperref[\detokenize{api:models.model.Lipid}]{\sphinxcrossref{Lipid}}}\DUrole{p}{{]}}}}}
\pysigstopsignatures
\end{fulllineitems}

\index{detection (models.model.Annotation attribute)@\spxentry{detection}\spxextra{models.model.Annotation attribute}}

\begin{fulllineitems}
\phantomsection\label{\detokenize{api:models.model.Annotation.detection}}
\pysigstartsignatures
\pysigline
{\sphinxbfcode{\sphinxupquote{detection}}\sphinxbfcode{\sphinxupquote{\DUrole{p}{:}\DUrole{w}{ }Mapped\DUrole{p}{{[}}{\hyperref[\detokenize{api:models.model.Detection}]{\sphinxcrossref{Detection}}}\DUrole{p}{{]}}}}}
\pysigstopsignatures
\end{fulllineitems}


\end{fulllineitems}



\section{utils.loading\_MS1}
\label{\detokenize{api:module-utils.loading_MS1}}\label{\detokenize{api:utils-loading-ms1}}\index{module@\spxentry{module}!utils.loading\_MS1@\spxentry{utils.loading\_MS1}}\index{utils.loading\_MS1@\spxentry{utils.loading\_MS1}!module@\spxentry{module}}\index{DB\_MS1() (in module utils.loading\_MS1)@\spxentry{DB\_MS1()}\spxextra{in module utils.loading\_MS1}}

\begin{fulllineitems}
\phantomsection\label{\detokenize{api:utils.loading_MS1.DB_MS1}}
\pysigstartsignatures
\pysiglinewithargsret
{\sphinxcode{\sphinxupquote{utils.loading\_MS1.}}\sphinxbfcode{\sphinxupquote{DB\_MS1}}}
{\sphinxparam{\DUrole{n}{df}\DUrole{p}{:}\DUrole{w}{ }\DUrole{n}{DataFrame}}\sphinxparamcomma \sphinxparam{\DUrole{n}{filename}\DUrole{p}{:}\DUrole{w}{ }\DUrole{n}{str}}\sphinxparamcomma \sphinxparam{\DUrole{n}{integrator}\DUrole{p}{:}\DUrole{w}{ }\DUrole{n}{str}}\sphinxparamcomma \sphinxparam{\DUrole{n}{file\_row\_name}\DUrole{p}{:}\DUrole{w}{ }\DUrole{n}{str}}}
{{ $\rightarrow$ None}}
\pysigstopsignatures
\sphinxAtStartPar
Insert MS1 annotation data into the database.
For each row of the DataFrame, creates and inserts a record into the three tables: Detection, Lipid and Annotation.
\begin{quote}\begin{description}
\sphinxlineitem{Parameters}\begin{itemize}
\item {} 
\sphinxAtStartPar
\sphinxstyleliteralstrong{\sphinxupquote{df}} (\sphinxstyleliteralemphasis{\sphinxupquote{pandas.DataFrame}}) \textendash{} DataFrame containing the columns Lipid\_Name, Formula, Precursor\_MZ, Neutral\_mass, Adduct, Molecular\_weight, MS\_level, Num\_Peaks, Lipid\_category, Lipid\_class, Lipid\_subclass
and optionally RT and CCS.

\item {} 
\sphinxAtStartPar
\sphinxstyleliteralstrong{\sphinxupquote{filename}} (\sphinxstyleliteralemphasis{\sphinxupquote{str}}) \textendash{} name of the file being integrated.

\item {} 
\sphinxAtStartPar
\sphinxstyleliteralstrong{\sphinxupquote{integrator}} (\sphinxstyleliteralemphasis{\sphinxupquote{str}}) \textendash{} name of the person performing the integration.

\item {} 
\sphinxAtStartPar
\sphinxstyleliteralstrong{\sphinxupquote{file\_row\_name}} (\sphinxstyleliteralemphasis{\sphinxupquote{str}}) \textendash{} name of the file that provided the annotations.

\end{itemize}

\sphinxlineitem{Raises}\begin{itemize}
\item {} 
\sphinxAtStartPar
\sphinxstyleliteralstrong{\sphinxupquote{Exception}} \textendash{} on insertion error no row is committed (automatic ROLLBACK) and the error is logged.

\item {} 
\sphinxAtStartPar
\sphinxstyleliteralstrong{\sphinxupquote{PostIntegrationError}} \textendash{} if the rows were committed successfully but the post\sphinxhyphen{}commit
backup or history logging failed \sphinxhyphen{} the data is already in the database.

\end{itemize}

\end{description}\end{quote}

\end{fulllineitems}



\section{utils.molecular\_weight}
\label{\detokenize{api:module-utils.molecular_weight}}\label{\detokenize{api:utils-molecular-weight}}\index{module@\spxentry{module}!utils.molecular\_weight@\spxentry{utils.molecular\_weight}}\index{utils.molecular\_weight@\spxentry{utils.molecular\_weight}!module@\spxentry{module}}\index{molecular\_weight() (in module utils.molecular\_weight)@\spxentry{molecular\_weight()}\spxextra{in module utils.molecular\_weight}}

\begin{fulllineitems}
\phantomsection\label{\detokenize{api:utils.molecular_weight.molecular_weight}}
\pysigstartsignatures
\pysiglinewithargsret
{\sphinxcode{\sphinxupquote{utils.molecular\_weight.}}\sphinxbfcode{\sphinxupquote{molecular\_weight}}}
{\sphinxparam{\DUrole{n}{formula}\DUrole{p}{:}\DUrole{w}{ }\DUrole{n}{str}}}
{{ $\rightarrow$ float\DUrole{w}{ }\DUrole{p}{|}\DUrole{w}{ }None}}
\pysigstopsignatures
\sphinxAtStartPar
Compute the average molecular weight of the lipid from its chemical formula.
Uses molmass.Formula.mass for the calculation, then rounds to 6 decimal places.
Returns None without raising an exception if the formula is None, empty, or invalid.
\begin{quote}\begin{description}
\sphinxlineitem{Parameters}
\sphinxAtStartPar
\sphinxstyleliteralstrong{\sphinxupquote{formula}} (\sphinxstyleliteralemphasis{\sphinxupquote{str}}) \textendash{} chemical formula

\sphinxlineitem{Returns}
\sphinxAtStartPar
average molecular weight rounded to 6 decimal places, or None if invalid.

\sphinxlineitem{Return type}
\sphinxAtStartPar
float or None

\end{description}\end{quote}

\end{fulllineitems}



\section{utils.monoisotopic}
\label{\detokenize{api:module-utils.monoisotopic}}\label{\detokenize{api:utils-monoisotopic}}\index{module@\spxentry{module}!utils.monoisotopic@\spxentry{utils.monoisotopic}}\index{utils.monoisotopic@\spxentry{utils.monoisotopic}!module@\spxentry{module}}\index{monoisotopic\_mass() (in module utils.monoisotopic)@\spxentry{monoisotopic\_mass()}\spxextra{in module utils.monoisotopic}}

\begin{fulllineitems}
\phantomsection\label{\detokenize{api:utils.monoisotopic.monoisotopic_mass}}
\pysigstartsignatures
\pysiglinewithargsret
{\sphinxcode{\sphinxupquote{utils.monoisotopic.}}\sphinxbfcode{\sphinxupquote{monoisotopic\_mass}}}
{\sphinxparam{\DUrole{n}{formula}\DUrole{p}{:}\DUrole{w}{ }\DUrole{n}{str\DUrole{w}{ }\DUrole{p}{|}\DUrole{w}{ }None}}}
{{ $\rightarrow$ float\DUrole{w}{ }\DUrole{p}{|}\DUrole{w}{ }None}}
\pysigstopsignatures
\sphinxAtStartPar
Compute the monoisotopic mass of the lipid from its chemical formula.
Uses molmass.Formula.monoisotopic\_mass for the calculation.
Returns None without raising an exception if the formula is null, empty, or invalid.
\begin{quote}\begin{description}
\sphinxlineitem{Parameters}
\sphinxAtStartPar
\sphinxstyleliteralstrong{\sphinxupquote{formula}} (\sphinxstyleliteralemphasis{\sphinxupquote{str}}\sphinxstyleliteralemphasis{\sphinxupquote{ or }}\sphinxstyleliteralemphasis{\sphinxupquote{None}}) \textendash{} chemical formula

\sphinxlineitem{Returns}
\sphinxAtStartPar
monoisotopic mass rounded to 6 decimal places, or None if invalid.

\sphinxlineitem{Return type}
\sphinxAtStartPar
float or None

\end{description}\end{quote}

\end{fulllineitems}



\section{utils.msp\_export}
\label{\detokenize{api:module-utils.msp_export}}\label{\detokenize{api:utils-msp-export}}\index{module@\spxentry{module}!utils.msp\_export@\spxentry{utils.msp\_export}}\index{utils.msp\_export@\spxentry{utils.msp\_export}!module@\spxentry{module}}\index{generate\_msp() (in module utils.msp\_export)@\spxentry{generate\_msp()}\spxextra{in module utils.msp\_export}}

\begin{fulllineitems}
\phantomsection\label{\detokenize{api:utils.msp_export.generate_msp}}
\pysigstartsignatures
\pysiglinewithargsret
{\sphinxcode{\sphinxupquote{utils.msp\_export.}}\sphinxbfcode{\sphinxupquote{generate\_msp}}}
{\sphinxparam{\DUrole{n}{engine}}\sphinxparamcomma \sphinxparam{\DUrole{n}{categories}\DUrole{p}{:}\DUrole{w}{ }\DUrole{n}{list}}\sphinxparamcomma \sphinxparam{\DUrole{n}{classes}\DUrole{p}{:}\DUrole{w}{ }\DUrole{n}{list}}\sphinxparamcomma \sphinxparam{\DUrole{n}{sub\_classes}\DUrole{p}{:}\DUrole{w}{ }\DUrole{n}{list}}\sphinxparamcomma \sphinxparam{\DUrole{n}{mz\_range}\DUrole{p}{:}\DUrole{w}{ }\DUrole{n}{tuple}}\sphinxparamcomma \sphinxparam{\DUrole{n}{ionisation\_modes}\DUrole{p}{:}\DUrole{w}{ }\DUrole{n}{list\DUrole{w}{ }\DUrole{p}{|}\DUrole{w}{ }None}\DUrole{w}{ }\DUrole{o}{=}\DUrole{w}{ }\DUrole{default_value}{None}}}
{{ $\rightarrow$ str}}
\pysigstopsignatures
\sphinxAtStartPar
Returns a string in .msp format for MS2 detections matching the filters.
\begin{quote}\begin{description}
\sphinxlineitem{Parameters}\begin{itemize}
\item {} 
\sphinxAtStartPar
\sphinxstyleliteralstrong{\sphinxupquote{engine}} \textendash{} engine: SQLAlchemy engine connected to the database.

\item {} 
\sphinxAtStartPar
\sphinxstyleliteralstrong{\sphinxupquote{categories}} \textendash{} selected lipid category list (empty = all).

\item {} 
\sphinxAtStartPar
\sphinxstyleliteralstrong{\sphinxupquote{classes}} \textendash{} selected lipid class list (empty = all).

\item {} 
\sphinxAtStartPar
\sphinxstyleliteralstrong{\sphinxupquote{sub\_classes}} \textendash{} selected lipid subclass list (empty = all).

\item {} 
\sphinxAtStartPar
\sphinxstyleliteralstrong{\sphinxupquote{mz\_range}} \textendash{} tuple (mz\_min, mz\_max) for filtering on Precursor\_MZ.

\item {} 
\sphinxAtStartPar
\sphinxstyleliteralstrong{\sphinxupquote{ionisation\_modes}} \textendash{} selected ionisation mode list, “Positive”/”Negative”

\end{itemize}

\sphinxlineitem{Returns}
\sphinxAtStartPar
.msp file ready to be downloaded.

\sphinxlineitem{Return type}
\sphinxAtStartPar
str

\end{description}\end{quote}

\end{fulllineitems}



\section{utils.neutral\_mass}
\label{\detokenize{api:module-utils.neutral_mass}}\label{\detokenize{api:utils-neutral-mass}}\index{module@\spxentry{module}!utils.neutral\_mass@\spxentry{utils.neutral\_mass}}\index{utils.neutral\_mass@\spxentry{utils.neutral\_mass}!module@\spxentry{module}}\index{json\_adducts() (in module utils.neutral\_mass)@\spxentry{json\_adducts()}\spxextra{in module utils.neutral\_mass}}

\begin{fulllineitems}
\phantomsection\label{\detokenize{api:utils.neutral_mass.json_adducts}}
\pysigstartsignatures
\pysiglinewithargsret
{\sphinxcode{\sphinxupquote{utils.neutral\_mass.}}\sphinxbfcode{\sphinxupquote{json\_adducts}}}
{}
{{ $\rightarrow$ dict}}
\pysigstopsignatures
\sphinxAtStartPar
Load the custom adducts persisted in the JSON file at ADDUCTS\_PATH.
\begin{quote}\begin{description}
\sphinxlineitem{Returns}
\sphinxAtStartPar
mapping of adduct name to \{“shift”: float, “sign”: int\}, or an empty
dict if the file does not exist or is empty.

\sphinxlineitem{Return type}
\sphinxAtStartPar
dict

\end{description}\end{quote}

\sphinxAtStartPar
Example:

\begin{sphinxVerbatim}[commandchars=\\\{\}]
\PYG{g+gp}{\PYGZgt{}\PYGZgt{}\PYGZgt{} }\PYG{n}{json\PYGZus{}adducts}\PYG{p}{(}\PYG{p}{)}
\PYG{g+go}{\PYGZob{}\PYGZsq{}[M+Na]+\PYGZsq{}: \PYGZob{}\PYGZsq{}shift\PYGZsq{}: 22.989221, \PYGZsq{}sign\PYGZsq{}: \PYGZhy{}1\PYGZcb{}\PYGZcb{}}
\end{sphinxVerbatim}

\end{fulllineitems}

\index{list\_adducts() (in module utils.neutral\_mass)@\spxentry{list\_adducts()}\spxextra{in module utils.neutral\_mass}}

\begin{fulllineitems}
\phantomsection\label{\detokenize{api:utils.neutral_mass.list_adducts}}
\pysigstartsignatures
\pysiglinewithargsret
{\sphinxcode{\sphinxupquote{utils.neutral\_mass.}}\sphinxbfcode{\sphinxupquote{list\_adducts}}}
{}
{{ $\rightarrow$ dict}}
\pysigstopsignatures
\sphinxAtStartPar
Return the currently persisted custom adducts.
\begin{quote}\begin{description}
\sphinxlineitem{Returns}
\sphinxAtStartPar
mapping of adduct name to \{“shift”: float, “sign”: int\}.

\sphinxlineitem{Return type}
\sphinxAtStartPar
dict

\end{description}\end{quote}

\sphinxAtStartPar
Example:

\begin{sphinxVerbatim}[commandchars=\\\{\}]
\PYG{g+gp}{\PYGZgt{}\PYGZgt{}\PYGZgt{} }\PYG{n}{list\PYGZus{}adducts}\PYG{p}{(}\PYG{p}{)}
\PYG{g+go}{\PYGZob{}\PYGZsq{}[M+Na]+\PYGZsq{}: \PYGZob{}\PYGZsq{}shift\PYGZsq{}: 22.989221, \PYGZsq{}sign\PYGZsq{}: \PYGZhy{}1\PYGZcb{}\PYGZcb{}}
\end{sphinxVerbatim}

\end{fulllineitems}

\index{add\_adduct() (in module utils.neutral\_mass)@\spxentry{add\_adduct()}\spxextra{in module utils.neutral\_mass}}

\begin{fulllineitems}
\phantomsection\label{\detokenize{api:utils.neutral_mass.add_adduct}}
\pysigstartsignatures
\pysiglinewithargsret
{\sphinxcode{\sphinxupquote{utils.neutral\_mass.}}\sphinxbfcode{\sphinxupquote{add\_adduct}}}
{\sphinxparam{\DUrole{n}{name}\DUrole{p}{:}\DUrole{w}{ }\DUrole{n}{str}}\sphinxparamcomma \sphinxparam{\DUrole{n}{shift}\DUrole{p}{:}\DUrole{w}{ }\DUrole{n}{float}}\sphinxparamcomma \sphinxparam{\DUrole{n}{sign}\DUrole{p}{:}\DUrole{w}{ }\DUrole{n}{int}}}
{{ $\rightarrow$ None}}
\pysigstopsignatures
\sphinxAtStartPar
Validate and add a new custom adduct.
\begin{quote}\begin{description}
\sphinxlineitem{Parameters}\begin{itemize}
\item {} 
\sphinxAtStartPar
\sphinxstyleliteralstrong{\sphinxupquote{name}} (\sphinxstyleliteralemphasis{\sphinxupquote{str}}) \textendash{} adduct name. Comparison against existing adducts is
case\sphinxhyphen{}insensitive, matching the lookup done by adduct().

\item {} 
\sphinxAtStartPar
\sphinxstyleliteralstrong{\sphinxupquote{shift}} (\sphinxstyleliteralemphasis{\sphinxupquote{float}}) \textendash{} mass shift in Daltons, must be strictly positive.

\item {} 
\sphinxAtStartPar
\sphinxstyleliteralstrong{\sphinxupquote{sign}} (\sphinxstyleliteralemphasis{\sphinxupquote{int}}) \textendash{} \sphinxhyphen{}1 if the adduct is formed by addition to the neutral molecule
(shift is subtracted to recover the neutral mass), +1 if formed by loss
(shift is added back).

\end{itemize}

\sphinxlineitem{Raises}
\sphinxAtStartPar
\sphinxstyleliteralstrong{\sphinxupquote{ValueError}} \textendash{} if name is empty or already used (case\sphinxhyphen{}insensitively),
sign is not \sphinxhyphen{}1/1, or shift is not a strictly positive number.

\end{description}\end{quote}

\sphinxAtStartPar
Example:

\begin{sphinxVerbatim}[commandchars=\\\{\}]
\PYG{g+gp}{\PYGZgt{}\PYGZgt{}\PYGZgt{} }\PYG{c+c1}{\PYGZsh{} [M+Na]+ is formed by addition of sodium, so sign=\PYGZhy{}1 (the shift is}
\PYG{g+gp}{\PYGZgt{}\PYGZgt{}\PYGZgt{} }\PYG{c+c1}{\PYGZsh{} subtracted back to recover the neutral mass from the precursor m/z).}
\PYG{g+gp}{\PYGZgt{}\PYGZgt{}\PYGZgt{} }\PYG{n}{add\PYGZus{}adduct}\PYG{p}{(}\PYG{l+s+s2}{\PYGZdq{}}\PYG{l+s+s2}{[M+Na]+}\PYG{l+s+s2}{\PYGZdq{}}\PYG{p}{,} \PYG{n}{shift}\PYG{o}{=}\PYG{l+m+mf}{22.989221}\PYG{p}{,} \PYG{n}{sign}\PYG{o}{=}\PYG{o}{\PYGZhy{}}\PYG{l+m+mi}{1}\PYG{p}{)}
\PYG{g+gp}{\PYGZgt{}\PYGZgt{}\PYGZgt{} }\PYG{n}{list\PYGZus{}adducts}\PYG{p}{(}\PYG{p}{)}
\PYG{g+go}{\PYGZob{}\PYGZsq{}[M+Na]+\PYGZsq{}: \PYGZob{}\PYGZsq{}shift\PYGZsq{}: 22.989221, \PYGZsq{}sign\PYGZsq{}: \PYGZhy{}1\PYGZcb{}\PYGZcb{}}
\end{sphinxVerbatim}

\end{fulllineitems}

\index{remove\_adduct() (in module utils.neutral\_mass)@\spxentry{remove\_adduct()}\spxextra{in module utils.neutral\_mass}}

\begin{fulllineitems}
\phantomsection\label{\detokenize{api:utils.neutral_mass.remove_adduct}}
\pysigstartsignatures
\pysiglinewithargsret
{\sphinxcode{\sphinxupquote{utils.neutral\_mass.}}\sphinxbfcode{\sphinxupquote{remove\_adduct}}}
{\sphinxparam{\DUrole{n}{name}\DUrole{p}{:}\DUrole{w}{ }\DUrole{n}{str}}}
{{ $\rightarrow$ None}}
\pysigstopsignatures
\sphinxAtStartPar
Remove a previously added custom adduct. Built\sphinxhyphen{}in adducts cannot be removed.
\begin{quote}\begin{description}
\sphinxlineitem{Parameters}
\sphinxAtStartPar
\sphinxstyleliteralstrong{\sphinxupquote{name}} (\sphinxstyleliteralemphasis{\sphinxupquote{str}}) \textendash{} name of the custom adduct to remove.

\sphinxlineitem{Raises}
\sphinxAtStartPar
\sphinxstyleliteralstrong{\sphinxupquote{ValueError}} \textendash{} if name is not a persisted custom adduct.

\end{description}\end{quote}

\sphinxAtStartPar
Example:

\begin{sphinxVerbatim}[commandchars=\\\{\}]
\PYG{g+gp}{\PYGZgt{}\PYGZgt{}\PYGZgt{} }\PYG{n}{remove\PYGZus{}adduct}\PYG{p}{(}\PYG{l+s+s2}{\PYGZdq{}}\PYG{l+s+s2}{[M+Na]+}\PYG{l+s+s2}{\PYGZdq{}}\PYG{p}{)}
\PYG{g+gp}{\PYGZgt{}\PYGZgt{}\PYGZgt{} }\PYG{n}{list\PYGZus{}adducts}\PYG{p}{(}\PYG{p}{)}
\PYG{g+go}{\PYGZob{}\PYGZcb{}}
\end{sphinxVerbatim}

\end{fulllineitems}

\index{adduct() (in module utils.neutral\_mass)@\spxentry{adduct()}\spxextra{in module utils.neutral\_mass}}

\begin{fulllineitems}
\phantomsection\label{\detokenize{api:utils.neutral_mass.adduct}}
\pysigstartsignatures
\pysiglinewithargsret
{\sphinxcode{\sphinxupquote{utils.neutral\_mass.}}\sphinxbfcode{\sphinxupquote{adduct}}}
{\sphinxparam{\DUrole{n}{raw\_adduct}}}
{{ $\rightarrow$ str}}
\pysigstopsignatures
\sphinxAtStartPar
Validate and normalize a raw adduct value into one of the standard adduct strings. The adduct is required for the neutral mass calculation
and the completion of the “Adduct” column in the output table.
\begin{quote}\begin{description}
\sphinxlineitem{Parameters}
\sphinxAtStartPar
\sphinxstyleliteralstrong{\sphinxupquote{raw\_adduct}} \textendash{} raw value to resolve

\sphinxlineitem{Raises}
\sphinxAtStartPar
\sphinxstyleliteralstrong{\sphinxupquote{ValueError}} \textendash{} if raw\_adduct is missing or not a recognized adduct.

\sphinxlineitem{Returns}
\sphinxAtStartPar
standard adduct string

\sphinxlineitem{Return type}
\sphinxAtStartPar
str

\end{description}\end{quote}

\sphinxAtStartPar
Example:

\begin{sphinxVerbatim}[commandchars=\\\{\}]
\PYG{g+gp}{\PYGZgt{}\PYGZgt{}\PYGZgt{} }\PYG{c+c1}{\PYGZsh{} Case and surrounding spaces are ignored, only the standard form is returned.}
\PYG{g+gp}{\PYGZgt{}\PYGZgt{}\PYGZgt{} }\PYG{n}{adduct}\PYG{p}{(}\PYG{l+s+s2}{\PYGZdq{}}\PYG{l+s+s2}{  [m+h]+ }\PYG{l+s+s2}{\PYGZdq{}}\PYG{p}{)}
\PYG{g+go}{\PYGZsq{}[M+H]+\PYGZsq{}}
\end{sphinxVerbatim}

\end{fulllineitems}

\index{neutral\_mass() (in module utils.neutral\_mass)@\spxentry{neutral\_mass()}\spxextra{in module utils.neutral\_mass}}

\begin{fulllineitems}
\phantomsection\label{\detokenize{api:utils.neutral_mass.neutral_mass}}
\pysigstartsignatures
\pysiglinewithargsret
{\sphinxcode{\sphinxupquote{utils.neutral\_mass.}}\sphinxbfcode{\sphinxupquote{neutral\_mass}}}
{\sphinxparam{\DUrole{n}{Precursor\_MZ}\DUrole{p}{:}\DUrole{w}{ }\DUrole{n}{float}}\sphinxparamcomma \sphinxparam{\DUrole{n}{adduct}\DUrole{p}{:}\DUrole{w}{ }\DUrole{n}{str}}}
{{ $\rightarrow$ float}}
\pysigstopsignatures
\sphinxAtStartPar
Calculate the neutral mass from the precursor m/z ratio and the precursor adduct.
\begin{quote}\begin{description}
\sphinxlineitem{Parameters}\begin{itemize}
\item {} 
\sphinxAtStartPar
\sphinxstyleliteralstrong{\sphinxupquote{Precursor\_MZ}} (\sphinxstyleliteralemphasis{\sphinxupquote{float}}) \textendash{} m/z ratio of the precursor.

\item {} 
\sphinxAtStartPar
\sphinxstyleliteralstrong{\sphinxupquote{adduct}} (\sphinxstyleliteralemphasis{\sphinxupquote{str}}) \textendash{} precursor adduct

\end{itemize}

\sphinxlineitem{Returns}
\sphinxAtStartPar
neutral mass in Daltons, rounded to 6 decimal places.

\sphinxlineitem{Return type}
\sphinxAtStartPar
float

\end{description}\end{quote}

\end{fulllineitems}



\section{scripts.formula}
\label{\detokenize{api:scripts-formula}}

\section{scripts.init\_db}
\label{\detokenize{api:module-scripts.init_db}}\label{\detokenize{api:scripts-init-db}}\index{module@\spxentry{module}!scripts.init\_db@\spxentry{scripts.init\_db}}\index{scripts.init\_db@\spxentry{scripts.init\_db}!module@\spxentry{module}}

\subsection{init\_db.py}
\label{\detokenize{api:init-db-py}}
\sphinxAtStartPar
Create the database and all tables if they do not exist yet
Initialization script for the database “BacLipidDB.db”
Run this script only once to create the SQLite tables defined in models/model.py
\begin{description}
\sphinxlineitem{Usage :}
\sphinxAtStartPar
python scripts/init\_db.py

\end{description}
\index{main() (in module scripts.init\_db)@\spxentry{main()}\spxextra{in module scripts.init\_db}}

\begin{fulllineitems}
\phantomsection\label{\detokenize{api:scripts.init_db.main}}
\pysigstartsignatures
\pysiglinewithargsret
{\sphinxcode{\sphinxupquote{scripts.init\_db.}}\sphinxbfcode{\sphinxupquote{main}}}
{}
{}
\pysigstopsignatures
\end{fulllineitems}



\section{scripts.loading\_MS2}
\label{\detokenize{api:scripts-loading-ms2}}

\section{scripts.reset\_db}
\label{\detokenize{api:module-scripts.reset_db}}\label{\detokenize{api:scripts-reset-db}}\index{module@\spxentry{module}!scripts.reset\_db@\spxentry{scripts.reset\_db}}\index{scripts.reset\_db@\spxentry{scripts.reset\_db}!module@\spxentry{module}}

\subsection{reset\_db.py}
\label{\detokenize{api:reset-db-py}}
\sphinxAtStartPar
Delete all records from BacLipidDB without dropping the tables.
\begin{description}
\sphinxlineitem{The deletion order respects foreign key constraints:}\begin{enumerate}
\sphinxsetlistlabels{\arabic}{enumi}{enumii}{}{.}%
\item {} 
\sphinxAtStartPar
Annotation  (references to Lipid and Detection)

\item {} 
\sphinxAtStartPar
Fragment    (reference to Detection)

\item {} 
\sphinxAtStartPar
Detection

\item {} 
\sphinxAtStartPar
Lipid

\end{enumerate}

\sphinxlineitem{Usage :}
\sphinxAtStartPar
python scripts/reset\_db.py

\end{description}
\index{reset\_db() (in module scripts.reset\_db)@\spxentry{reset\_db()}\spxextra{in module scripts.reset\_db}}

\begin{fulllineitems}
\phantomsection\label{\detokenize{api:scripts.reset_db.reset_db}}
\pysigstartsignatures
\pysiglinewithargsret
{\sphinxcode{\sphinxupquote{scripts.reset\_db.}}\sphinxbfcode{\sphinxupquote{reset\_db}}}
{}
{}
\pysigstopsignatures
\end{fulllineitems}



\renewcommand{\indexname}{Python Module Index}
\begin{sphinxtheindex}
\let\bigletter\sphinxstyleindexlettergroup
\bigletter{c}
\item\relax\sphinxstyleindexentry{config}\sphinxstyleindexpageref{api:\detokenize{module-config}}
\indexspace
\bigletter{m}
\item\relax\sphinxstyleindexentry{models.model}\sphinxstyleindexpageref{api:\detokenize{module-models.model}}
\indexspace
\bigletter{s}
\item\relax\sphinxstyleindexentry{scripts.init\_db}\sphinxstyleindexpageref{api:\detokenize{module-scripts.init_db}}
\item\relax\sphinxstyleindexentry{scripts.reset\_db}\sphinxstyleindexpageref{api:\detokenize{module-scripts.reset_db}}
\indexspace
\bigletter{u}
\item\relax\sphinxstyleindexentry{utils.loading\_MS1}\sphinxstyleindexpageref{api:\detokenize{module-utils.loading_MS1}}
\item\relax\sphinxstyleindexentry{utils.molecular\_weight}\sphinxstyleindexpageref{api:\detokenize{module-utils.molecular_weight}}
\item\relax\sphinxstyleindexentry{utils.monoisotopic}\sphinxstyleindexpageref{api:\detokenize{module-utils.monoisotopic}}
\item\relax\sphinxstyleindexentry{utils.msp\_export}\sphinxstyleindexpageref{api:\detokenize{module-utils.msp_export}}
\item\relax\sphinxstyleindexentry{utils.neutral\_mass}\sphinxstyleindexpageref{api:\detokenize{module-utils.neutral_mass}}
\end{sphinxtheindex}

\renewcommand{\indexname}{Index}
\printindex
\end{document}